\documentclass[conference]{IEEEtran}
\usepackage[utf8]{inputenc}
\usepackage[T1]{fontenc}
\usepackage[french]{babel}
\usepackage{amsmath}
\usepackage{amsfonts}
\usepackage{amssymb}
\usepackage{amsthm}
\usepackage{graphicx}
\usepackage{mathrsfs}
\usepackage{cite}
% Packages ajoutés pour corriger les incohérences
\usepackage{algorithm2e}
\usepackage{booktabs}
\usepackage[caption=false,font=footnotesize]{subfig}
\usepackage{pifont}
\usepackage{multirow}
\usepackage{tabularx}
\usepackage{xcolor}
\usepackage{float}
\usepackage{stfloats}
\usepackage{hyperref}

% Environnements pour théorèmes et définitions
\newtheorem{definition}{Définition}
\newtheorem{theorem}{Théorème}
\newtheorem{lemma}{Lemme}
\newtheorem{proposition}{Proposition}

\title{Optimisation de l'Efficacité Énergétique dans les Réseaux 5G : Une Approche par Apprentissage par Renforcement Profond}
\author{
    \IEEEauthorblockN{Nada, Youssef IFQIR}
    \IEEEauthorblockA{Université ... , Faculté des Sciences et Techniques\\
    Année académique 2025--2026\\[0.5em]
    \includegraphics[height=1.8cm]{results/figures/Logofst.jpg}}}

\begin{document}

\maketitle

\begin{abstract}
L'essor des réseaux 5G s'accompagne d'une augmentation exponentielle de la consommation énergétique, ce qui rend crucial le développement de stratégies de contrôle de puissance énergétiquement efficaces. Cet article propose une approche novatrice basée sur l'apprentissage par renforcement profond (Deep Q-Network, DQN) pour le contrôle de puissance décentralisé dans les réseaux 5G, en vue de maximiser l'efficacité énergétique (EE) tout en respectant des contraintes de qualité de service (QoS). Nous commençons par une analyse approfondie de l'approche classique de Zappone et al. (2021), qui modélise le problème comme un jeu non-coopératif généralisé et propose un algorithme distribué fondé sur la théorie des jeux et l'équilibre de Nash généralisé (GNE). Nous détaillons les fondements mathématiques de cette approche : conditions de faisabilité (lemme 1), existence et unicité du GNE (propositions 3 et 4), calcul de la meilleure réponse via l'algorithme de Dinkelbach, et algorithme distribué itératif. Nous mettons en évidence les limites de cette méthode : nécessité d'un modèle parfait du canal, hypothèse de stationnarité, absence d'apprentissage et sensibilité aux erreurs d'estimation.

Pour surmonter ces limitations, nous reformulons le problème de contrôle de puissance comme un Processus de Décision Markovien (MDP) dans lequel chaque utilisateur est un agent intelligent. L'état de l'agent est composé d'observations locales (puissance courante, SINR, estimation du canal, historique des récompenses, etc.), l'action est le choix d'un niveau de puissance parmi 16 valeurs discrètes, et la récompense est une combinaison pondérée de l'EE instantanée, d'une pénalité en cas de violation de la QoS et d'une pénalité sur la puissance utilisée. Nous implémentons un réseau de neurones profond avec une architecture à trois couches cachées (256, 512, 256 neurones) et utilisons des techniques avancées de stabilisation de l'apprentissage : experience replay, réseau cible, Double DQN, Huber loss, gradient clipping et stratégie d'exploration $\varepsilon$-greedy décroissante.

Les expériences menées dans un environnement de simulation réaliste montrent que notre agent DQN atteint une EE moyenne de 515 000 bit/J, soit un gain de plus de 130 \% par rapport aux algorithmes distribués classiques, tout en maintenant un taux de violation de la QoS inférieur à 0,03 \%. De plus, l'approche DQN s'adapte dynamiquement aux variations des conditions radio et améliore ses performances avec l'expérience, contrairement aux méthodes analytiques. Nous visualisons également la topologie du réseau et la convergence des décisions à l'aide de la base de données graphe Neo4j, confirmant la stabilité et la cohérence de la politique apprise. Ces résultats démontrent la pertinence et l'efficacité de l'apprentissage par renforcement profond pour le contrôle de puissance dans les réseaux 5G, ouvrant la voie à des systèmes autonomes, adaptatifs et économes en énergie.
\end{abstract}

\begin{IEEEkeywords}
Réseaux 5G, Efficacité énergétique (EE), Contrôle de puissance décentralisé, Allocation de ressources, Contraintes de qualité de service (QoS), Théorie des jeux, Jeu non-coopératif, Jeu généralisé, Équilibre de Nash généralisé (GNE), Fonction standard, Algorithme de Dinkelbach, Meilleure réponse, Programmation fractionnaire, Apprentissage par renforcement profond (DRL), Deep Q-Network (DQN), Processus de Décision Markovien (MDP), Approche sans modèle, Adaptation dynamique, Environnement non stationnaire.
\end{IEEEkeywords}

%==============================================================================
\section{Introduction}
%==============================================================================

\subsection{Contexte et motivation}

Les réseaux de cinquième génération (5G) représentent une rupture technologique majeure dans l'histoire des communications sans fil. Avec des débits annoncés atteignant 10 Gbit/s, une latence inférieure à 1 ms et une densité de connexion pouvant atteindre 1 million de dispositifs par kilomètre carré \cite{3GPP}, la 5G doit répondre à des exigences sans précédent. Cependant, cette augmentation spectaculaire des performances s'accompagne d'une augmentation tout aussi spectaculaire de la consommation énergétique.

Selon \cite{Fehske2011}, le secteur des technologies de l'information et de la communication (TIC) représente déjà environ 5\% des émissions mondiales de CO$_2$, et ce pourcentage est en croissance rapide. Avec le déploiement massif des réseaux 5G, la situation risque de s'aggraver si des mesures ne sont pas prises. Il est estimé que le nombre de dispositifs connectés atteindra 50 milliards d'ici 2025 \cite{Ericsson2011}, et que les débits devront être multipliés par 1000 par rapport aux réseaux 4G \cite{Qualcomm}.

Une solution naïve consisterait à augmenter proportionnellement les puissances d'émission pour atteindre ces débits. Cependant, cette approche est non seulement irréaliste d'un point de vue énergétique, mais aussi inacceptable pour des raisons environnementales (émissions de gaz à effet de serre) et sanitaires (pollution électromagnétique). La seule solution viable est d'augmenter l'\textbf{efficacité énergétique} (EE) des réseaux, c'est-à-dire la quantité de données transmise par unité d'énergie consommée.

L'EE est généralement mesurée en bits par joule (bit/J) et peut être définie au niveau de l'utilisateur ou au niveau du réseau. L'objectif fondamental est de "transmettre plus de données en consommant moins d'énergie", ce qui nécessite des stratégies de contrôle de puissance sophistiquées.

\subsection{Le problème du contrôle de puissance}

Dans un réseau sans fil, chaque utilisateur doit choisir sa puissance d'émission. Ce choix apparemment simple a des conséquences complexes :

\begin{itemize}
    \item Il détermine le \textbf{débit} de l'utilisateur via le rapport signal-sur-bruit-plus-interférence (SINR).
    \item Il influence les \textbf{interférences} subies par les autres utilisateurs.
    \item Il affecte la \textbf{consommation énergétique} globale du réseau et de l'utilisateur.
\end{itemize}

Mathématiquement, pour un utilisateur $k$ utilisant $N$ blocs de ressources, le SINR sur le bloc $n$ s'exprime généralement par :

\begin{equation} \label{eq:sinr_general}
\gamma_{k,n} = \frac{\alpha_{k,n} p_{k,n}}{\sigma_{k,n}^2 + \phi_{k,n} p_{k,n} + \sum_{j \neq k} \omega_{k,j,n} p_{j,n}}
\end{equation}

où $p_{k,n}$ est la puissance d'émission de l'utilisateur $k$ sur le bloc $n$, $\alpha_{k,n}$ représente le gain du canal direct, $\phi_{k,n}$ modélise l'auto-interférence (due aux imperfections matérielles), $\omega_{k,j,n}$ est le coefficient d'interférence de l'utilisateur $j$ sur l'utilisateur $k$, et $\sigma_{k,n}^2$ est la puissance du bruit thermique.

L'efficacité énergétique individuelle de l'utilisateur $k$ est alors définie comme le rapport entre son débit total et sa puissance totale consommée :

\begin{equation} \label{eq:ee}
\eta_k = \frac{\displaystyle\sum_{n=1}^{N} B \log_2(1 + \gamma_{k,n})}{p_{c,k} + \displaystyle\sum_{n=1}^{N} p_{k,n}}
\end{equation}

où $B$ est la bande passante par bloc de ressources et $p_{c,k}$ est la puissance de circuit (consommation fixe).

Ce problème d'optimisation présente plusieurs difficultés majeures :
\begin{itemize}
    \item \textbf{Non-convexité} : la fonction objectif $\eta_k$ est non-convexe en raison de la présence du logarithme et du rapport.
    \item \textbf{Couplage fort} : les décisions des utilisateurs sont interdépendantes via les termes d'interférence $\sum_{j \neq k} \omega_{k,j,n} p_{j,n}$.
    \item \textbf{Contraintes de QoS} : chaque utilisateur doit généralement respecter un débit minimum, ce qui ajoute des contraintes supplémentaires.
    \item \textbf{Dynamique} : les canaux radio varient dans le temps (fading, mobilité).
\end{itemize}

\subsection{Concepts mathématiques fondamentaux}

Avant de présenter les approches existantes et notre contribution, il est nécessaire d'introduire certains concepts mathématiques qui seront utilisés tout au long de cet article. Ces concepts proviennent de la théorie des jeux et de l'optimisation, et constituent le cadre théorique sur lequel repose l'approche décentralisée classique.

\subsubsection{Théorie des jeux et équilibre de Nash}

La théorie des jeux fournit un cadre naturel pour modéliser les interactions stratégiques entre des décideurs autonomes. Dans notre contexte, chaque utilisateur du réseau est considéré comme un \textit{joueur} cherchant à maximiser sa propre efficacité énergétique. Formellement, un jeu non-coopératif est défini par un triplet $\mathcal{G} = \{\mathcal{K}, \{\mathcal{S}_k\}_{k\in\mathcal{K}}, \{u_k\}_{k\in\mathcal{K}}\}$, où $\mathcal{K} = \{1,\dots,K\}$ représente l'ensemble des joueurs, $\mathcal{S}_k$ l'ensemble des stratégies disponibles pour le joueur $k$, et $u_k : \mathcal{S}_1 \times \cdots \times \mathcal{S}_K \to \mathbb{R}$ sa fonction d'utilité.

Le concept central de la théorie des jeux non-coopératifs est l'\textit{équilibre de Nash}. Un profil de stratégies $\mathbf{s}^* = (s_1^*, \dots, s_K^*)$ constitue un équilibre de Nash si aucun joueur n'a intérêt à modifier unilatéralement sa stratégie, c'est-à-dire :
\[
u_k(s_k^*, \mathbf{s}_{-k}^*) \geq u_k(s_k, \mathbf{s}_{-k}^*) \quad \forall s_k \in \mathcal{S}_k, \forall k \in \mathcal{K},
\]
où $\mathbf{s}_{-k}^*$ désigne les stratégies de tous les joueurs excepté $k$. À l'équilibre, chaque joueur adopte la meilleure réponse possible compte tenu des choix des autres.

\subsubsection{Jeux généralisés et équilibre de Nash généralisé}

Dans notre problème de contrôle de puissance, une difficulté supplémentaire apparaît : l'ensemble des stratégies réalisables pour un utilisateur dépend des stratégies des autres. En effet, la contrainte de qualité de service (QoS) impose un SINR minimum, lequel dépend des interférences générées par les autres utilisateurs. On parle alors de \textit{jeu généralisé}, où l'ensemble des stratégies s'écrit $\mathcal{S}_k = \mathcal{S}_k(\mathbf{s}_{-k})$.

Dans ce cadre, la notion d'équilibre doit être adaptée. Un \textit{équilibre de Nash généralisé} (GNE) est un profil $\mathbf{p}^*$ tel que pour tout utilisateur $k$ :
\[
\mathbf{p}_k^* \in \mathcal{P}_k(\mathbf{p}_{-k}^*) \quad \text{et} \quad \eta_k(\mathbf{p}_k^*, \mathbf{p}_{-k}^*) \geq \eta_k(\mathbf{p}_k, \mathbf{p}_{-k}^*) \quad \forall \mathbf{p}_k \in \mathcal{P}_k(\mathbf{p}_{-k}^*),
\]
où $\mathcal{P}_k(\mathbf{p}_{-k})$ représente l'ensemble des puissances réalisables pour l'utilisateur $k$ étant donné les puissances des autres.

\subsubsection{Fonction standard et convergence}

Pour garantir l'unicité de l'équilibre et la convergence des algorithmes itératifs, Yates \cite{Yates1995} a introduit le concept de \textit{fonction standard}. Une fonction $f: \mathbb{R}_+^K \to \mathbb{R}_+^K$ est dite standard si elle satisfait trois propriétés essentielles :
\begin{itemize}
    \item \textbf{Positivité} : $f(\mathbf{p}) > 0$ pour tout $\mathbf{p} \geq 0$,
    \item \textbf{Monotonie} : si $\mathbf{p} \leq \mathbf{p}'$, alors $f(\mathbf{p}) \leq f(\mathbf{p}')$,
    \item \textbf{Scalabilité} : pour tout $\alpha > 1$, $\alpha f(\mathbf{p}) > f(\alpha \mathbf{p})$.
\end{itemize}
Si la fonction de meilleure réponse d'un jeu est standard, alors le jeu possède un unique équilibre et la mise à jour synchrone $p_k^{[t+1]} = B_k(\mathbf{p}_{-k}^{[t]})$ converge vers cet équilibre quel que soit le point de départ.

\subsubsection{Optimalité de Pareto}

L'équilibre de Nash n'est pas nécessairement optimal du point de vue collectif. La notion d'\textit{optimalité de Pareto} permet de caractériser les situations où aucune amélioration n'est possible sans détériorer la situation d'au moins un joueur. Formellement, un profil $\mathbf{p}$ est Pareto-optimal s'il n'existe pas d'autre profil $\mathbf{p}'$ tel que $\eta_k(\mathbf{p}') \geq \eta_k(\mathbf{p})$ pour tout $k$, avec au moins une inégalité stricte. Cette notion est importante car elle met en évidence la possible inefficacité des équilibres non-coopératifs.

\subsubsection{Conditions d'optimalité de Karush-Kuhn-Tucker (KKT)}

En optimisation sous contraintes, les \textit{conditions de Karush-Kuhn-Tucker} (KKT) jouent un rôle fondamental. Pour un problème de maximisation avec contraintes d'inégalité $g_i(\mathbf{x}) \leq 0$ et d'égalité $h_j(\mathbf{x}) = 0$, un point $\mathbf{x}^*$ satisfait les conditions KKT s'il existe des multiplicateurs $\lambda_i \geq 0$ et $\mu_j$ tels que :
\begin{align}
\nabla f(\mathbf{x}^*) + \sum_{i=1}^m \lambda_i \nabla g_i(\mathbf{x}^*) + \sum_{j=1}^p \mu_j \nabla h_j(\mathbf{x}^*) &= 0, \\
\lambda_i g_i(\mathbf{x}^*) &= 0 \quad \forall i, \\
g_i(\mathbf{x}^*) &\leq 0 \quad \forall i, \\
h_j(\mathbf{x}^*) &= 0 \quad \forall j.
\end{align}
Ces conditions sont nécessaires pour qu'un point soit un optimum local (sous certaines hypothèses de régularité). Dans le contexte de l'optimisation de l'EE, les algorithmes développés par Zappone et al. \cite{Zappone2021} convergent vers des points satisfaisant ces conditions.\subsection{Travaux connexes}

L'optimisation de l'efficacité énergétique dans les réseaux sans fil a fait l'objet de nombreux travaux. On peut distinguer deux grandes familles d'approches : les approches centralisées et les approches décentralisées.

\subsubsection{Approches centralisées}

Les approches centralisées supposent l'existence d'une entité centrale disposant d'une connaissance parfaite de tous les canaux et de toutes les puissances. Cette entité résout globalement le problème d'optimisation. Les premiers travaux dans ce domaine utilisaient la programmation fractionnaire \cite{Dinkelbach1967}, qui permet de traiter les problèmes d'optimisation de ratios.

Cependant, comme le soulignent Zappone et al. \cite{Zappone2021}, la présence d'interférences rend le numérateur de l'EE non-concave, ce qui empêche l'application directe de la programmation fractionnaire classique. Des solutions basées sur des schémas de transmission orthogonaux \cite{Ng2012, Xu2014} ou sur l'annulation d'interférences \cite{Xu2013} ont été proposées, mais elles conduisent à un faible facteur de réutilisation des ressources.

D'autres approches utilisent des techniques heuristiques comme l'optimisation alternée \cite{Du2014, He2013, He2014}, mais sans garantie de convergence. Venturino et al. \cite{Venturino2015} ont proposé un cadre unifié pour l'optimisation de l'EE dans les réseaux OFDMA.

\subsubsection{Approches décentralisées}

Les approches décentralisées modélisent le réseau comme un jeu non-coopératif où chaque utilisateur maximise sa propre EE. Miao et al. \cite{Miao2011} ont étudié le problème de l'équilibre de Nash dans ce contexte, sans contraintes de débit. Scutari et al. \cite{Scutari2008} ont proposé des stratégies de précodage optimales basées sur la théorie des jeux.

Zappone et al. \cite{Zappone2013} ont étendu ces travaux aux réseaux avec relais, tandis que Stupia et al. \cite{Stupia2015} ont utilisé une approche par inéquations variationnelles. Bacci et al. \cite{Bacci2015} ont introduit des contraintes de QoS dans le jeu, rendant le problème plus réaliste mais aussi plus complexe.

L'article de Zappone et al. \cite{Zappone2021} constitue une référence majeure dans ce domaine. Il propose un cadre unifié pour l'optimisation centralisée et décentralisée de l'EE avec contraintes de QoS et une expression généralisée du SINR englobant plusieurs technologies 5G (Massive MIMO avec imperfections matérielles, réseaux avec relais, etc.). Les auteurs développent :
\begin{itemize}
    \item Des conditions de faisabilité sous forme fermée (Lemme 1),
    \item Des algorithmes centralisés convergeant vers des points KKT,
    \item Une modélisation du problème décentralisé comme un jeu généralisé,
    \item Une preuve d'existence et d'unicité du GNE (Propositions 3 et 4),
    \item Un algorithme distribué basé sur les meilleures réponses,
    \item Une extension au cas multi-blocs de ressources.
\end{itemize}

\subsubsection{Limites des approches existantes}

Malgré leur rigueur mathématique, les approches de Zappone et al. présentent plusieurs limitations :

\begin{enumerate}
    \item \textbf{Nécessité d'un modèle parfait} : les algorithmes reposent sur la connaissance exacte des paramètres $\alpha_k$, $\phi_k$, $\omega_{kj}$. En pratique, ces paramètres sont estimés avec des erreurs.
    \item \textbf{Stationnarité} : la convergence des algorithmes suppose que les canaux varient lentement. Dans un environnement dynamique (mobilité, fading rapide), l'algorithme n'a pas le temps de converger.
    \item \textbf{Absence d'apprentissage} : les performances ne s'améliorent pas avec l'expérience. Chaque nouvelle situation est traitée indépendamment des précédentes.
    \item \textbf{Sensibilité aux erreurs} : les approximations et erreurs d'estimation peuvent dégrader significativement les performances.
\end{enumerate}

\subsection{Contributions de cet article}

Face à ces limitations, nous proposons une approche alternative basée sur l'apprentissage par renforcement profond (Deep Reinforcement Learning). Nos principales contributions sont :

\begin{enumerate}
    \item Une formulation du problème de contrôle de puissance comme un Processus de Décision Markovien (MDP) avec un état composé d'observations locales, des actions discrètes (niveaux de puissance) et une récompense combinant EE, pénalité QoS et pénalité de puissance.
    \item L'utilisation d'un réseau de neurones profond (DQN) pour apprendre une politique de contrôle de puissance optimale, sans nécessiter de modèle mathématique explicite du canal.
    \item Une architecture de réseau adaptée au problème (8 entrées, couches cachées [256,512,256], sortie 16 actions).
    \item Des techniques de stabilisation avancées (experience replay, target network, Double DQN, Huber loss, gradient clipping, $\varepsilon$-greedy décroissant).
    \item Une évaluation expérimentale montrant un gain significatif de l'EE par rapport aux approches classiques.
\end{enumerate}

\subsection{Organisation de l'article}

La suite de cet article est organisée comme suit. La section 2 présente en détail l'approche décentralisée classique de Zappone et al., avec toutes les dérivations mathématiques. La section 3 décrit notre approche DQN, incluant la modélisation MDP, l'architecture du réseau et les techniques d'entraînement. La section 4 présente les résultats expérimentaux et les compare aux approches classiques. Enfin, la section 5 conclut l'article et discute des perspectives futures.

%==============================================================================
\section{L'approche décentralisée classique : jeu généralisé et équilibre de Nash}
%==============================================================================

[Le contenu détaillé de cette section sera développé dans la prochaine réponse]

% ============================================================
% Chapitres DQN + Comparaison — Rapport PFE
% Contrôle de Puissance Décentralisé via Deep Q-Network
% pour Réseaux 5G Massive MIMO
% ============================================================
% À inclure dans le document principal via \documentclass[conference]{IEEEtran}
\usepackage[utf8]{inputenc}
\usepackage[T1]{fontenc}
\usepackage[french]{babel}
\usepackage{amsmath}
\usepackage{amsfonts}
\usepackage{amssymb}
\usepackage{amsthm}
\usepackage{graphicx}
\usepackage{mathrsfs}
\usepackage{cite}
% Packages ajoutés pour corriger les incohérences
\usepackage{algorithm2e}
\usepackage{booktabs}
\usepackage[caption=false,font=footnotesize]{subfig}
\usepackage{pifont}
\usepackage{multirow}
\usepackage{tabularx}
\usepackage{xcolor}
\usepackage{float}
\usepackage{stfloats}
\usepackage{hyperref}

% Environnements pour théorèmes et définitions
\newtheorem{definition}{Définition}
\newtheorem{theorem}{Théorème}
\newtheorem{lemma}{Lemme}
\newtheorem{proposition}{Proposition}

\title{Optimisation de l'Efficacité Énergétique dans les Réseaux 5G : Une Approche par Apprentissage par Renforcement Profond}
\author{
    \IEEEauthorblockN{Nada, Youssef IFQIR}
    \IEEEauthorblockA{Université ... , Faculté des Sciences et Techniques\\
    Année académique 2025--2026\\[0.5em]
    \includegraphics[height=1.8cm]{results/figures/Logofst.jpg}}}

\begin{document}

\maketitle

\begin{abstract}
L'essor des réseaux 5G s'accompagne d'une augmentation exponentielle de la consommation énergétique, ce qui rend crucial le développement de stratégies de contrôle de puissance énergétiquement efficaces. Cet article propose une approche novatrice basée sur l'apprentissage par renforcement profond (Deep Q-Network, DQN) pour le contrôle de puissance décentralisé dans les réseaux 5G, en vue de maximiser l'efficacité énergétique (EE) tout en respectant des contraintes de qualité de service (QoS). Nous commençons par une analyse approfondie de l'approche classique de Zappone et al. (2021), qui modélise le problème comme un jeu non-coopératif généralisé et propose un algorithme distribué fondé sur la théorie des jeux et l'équilibre de Nash généralisé (GNE). Nous détaillons les fondements mathématiques de cette approche : conditions de faisabilité (lemme 1), existence et unicité du GNE (propositions 3 et 4), calcul de la meilleure réponse via l'algorithme de Dinkelbach, et algorithme distribué itératif. Nous mettons en évidence les limites de cette méthode : nécessité d'un modèle parfait du canal, hypothèse de stationnarité, absence d'apprentissage et sensibilité aux erreurs d'estimation.

Pour surmonter ces limitations, nous reformulons le problème de contrôle de puissance comme un Processus de Décision Markovien (MDP) dans lequel chaque utilisateur est un agent intelligent. L'état de l'agent est composé d'observations locales (puissance courante, SINR, estimation du canal, historique des récompenses, etc.), l'action est le choix d'un niveau de puissance parmi 16 valeurs discrètes, et la récompense est une combinaison pondérée de l'EE instantanée, d'une pénalité en cas de violation de la QoS et d'une pénalité sur la puissance utilisée. Nous implémentons un réseau de neurones profond avec une architecture à trois couches cachées (256, 512, 256 neurones) et utilisons des techniques avancées de stabilisation de l'apprentissage : experience replay, réseau cible, Double DQN, Huber loss, gradient clipping et stratégie d'exploration $\varepsilon$-greedy décroissante.

Les expériences menées dans un environnement de simulation réaliste montrent que notre agent DQN atteint une EE moyenne de 515 000 bit/J, soit un gain de plus de 130 \% par rapport aux algorithmes distribués classiques, tout en maintenant un taux de violation de la QoS inférieur à 0,03 \%. De plus, l'approche DQN s'adapte dynamiquement aux variations des conditions radio et améliore ses performances avec l'expérience, contrairement aux méthodes analytiques. Nous visualisons également la topologie du réseau et la convergence des décisions à l'aide de la base de données graphe Neo4j, confirmant la stabilité et la cohérence de la politique apprise. Ces résultats démontrent la pertinence et l'efficacité de l'apprentissage par renforcement profond pour le contrôle de puissance dans les réseaux 5G, ouvrant la voie à des systèmes autonomes, adaptatifs et économes en énergie.
\end{abstract}

\begin{IEEEkeywords}
Réseaux 5G, Efficacité énergétique (EE), Contrôle de puissance décentralisé, Allocation de ressources, Contraintes de qualité de service (QoS), Théorie des jeux, Jeu non-coopératif, Jeu généralisé, Équilibre de Nash généralisé (GNE), Fonction standard, Algorithme de Dinkelbach, Meilleure réponse, Programmation fractionnaire, Apprentissage par renforcement profond (DRL), Deep Q-Network (DQN), Processus de Décision Markovien (MDP), Approche sans modèle, Adaptation dynamique, Environnement non stationnaire.
\end{IEEEkeywords}

%==============================================================================
\section{Introduction}
%==============================================================================

\subsection{Contexte et motivation}

Les réseaux de cinquième génération (5G) représentent une rupture technologique majeure dans l'histoire des communications sans fil. Avec des débits annoncés atteignant 10 Gbit/s, une latence inférieure à 1 ms et une densité de connexion pouvant atteindre 1 million de dispositifs par kilomètre carré \cite{3GPP}, la 5G doit répondre à des exigences sans précédent. Cependant, cette augmentation spectaculaire des performances s'accompagne d'une augmentation tout aussi spectaculaire de la consommation énergétique.

Selon \cite{Fehske2011}, le secteur des technologies de l'information et de la communication (TIC) représente déjà environ 5\% des émissions mondiales de CO$_2$, et ce pourcentage est en croissance rapide. Avec le déploiement massif des réseaux 5G, la situation risque de s'aggraver si des mesures ne sont pas prises. Il est estimé que le nombre de dispositifs connectés atteindra 50 milliards d'ici 2025 \cite{Ericsson2011}, et que les débits devront être multipliés par 1000 par rapport aux réseaux 4G \cite{Qualcomm}.

Une solution naïve consisterait à augmenter proportionnellement les puissances d'émission pour atteindre ces débits. Cependant, cette approche est non seulement irréaliste d'un point de vue énergétique, mais aussi inacceptable pour des raisons environnementales (émissions de gaz à effet de serre) et sanitaires (pollution électromagnétique). La seule solution viable est d'augmenter l'\textbf{efficacité énergétique} (EE) des réseaux, c'est-à-dire la quantité de données transmise par unité d'énergie consommée.

L'EE est généralement mesurée en bits par joule (bit/J) et peut être définie au niveau de l'utilisateur ou au niveau du réseau. L'objectif fondamental est de "transmettre plus de données en consommant moins d'énergie", ce qui nécessite des stratégies de contrôle de puissance sophistiquées.

\subsection{Le problème du contrôle de puissance}

Dans un réseau sans fil, chaque utilisateur doit choisir sa puissance d'émission. Ce choix apparemment simple a des conséquences complexes :

\begin{itemize}
    \item Il détermine le \textbf{débit} de l'utilisateur via le rapport signal-sur-bruit-plus-interférence (SINR).
    \item Il influence les \textbf{interférences} subies par les autres utilisateurs.
    \item Il affecte la \textbf{consommation énergétique} globale du réseau et de l'utilisateur.
\end{itemize}

Mathématiquement, pour un utilisateur $k$ utilisant $N$ blocs de ressources, le SINR sur le bloc $n$ s'exprime généralement par :

\begin{equation} \label{eq:sinr_general}
\gamma_{k,n} = \frac{\alpha_{k,n} p_{k,n}}{\sigma_{k,n}^2 + \phi_{k,n} p_{k,n} + \sum_{j \neq k} \omega_{k,j,n} p_{j,n}}
\end{equation}

où $p_{k,n}$ est la puissance d'émission de l'utilisateur $k$ sur le bloc $n$, $\alpha_{k,n}$ représente le gain du canal direct, $\phi_{k,n}$ modélise l'auto-interférence (due aux imperfections matérielles), $\omega_{k,j,n}$ est le coefficient d'interférence de l'utilisateur $j$ sur l'utilisateur $k$, et $\sigma_{k,n}^2$ est la puissance du bruit thermique.

L'efficacité énergétique individuelle de l'utilisateur $k$ est alors définie comme le rapport entre son débit total et sa puissance totale consommée :

\begin{equation} \label{eq:ee}
\eta_k = \frac{\displaystyle\sum_{n=1}^{N} B \log_2(1 + \gamma_{k,n})}{p_{c,k} + \displaystyle\sum_{n=1}^{N} p_{k,n}}
\end{equation}

où $B$ est la bande passante par bloc de ressources et $p_{c,k}$ est la puissance de circuit (consommation fixe).

Ce problème d'optimisation présente plusieurs difficultés majeures :
\begin{itemize}
    \item \textbf{Non-convexité} : la fonction objectif $\eta_k$ est non-convexe en raison de la présence du logarithme et du rapport.
    \item \textbf{Couplage fort} : les décisions des utilisateurs sont interdépendantes via les termes d'interférence $\sum_{j \neq k} \omega_{k,j,n} p_{j,n}$.
    \item \textbf{Contraintes de QoS} : chaque utilisateur doit généralement respecter un débit minimum, ce qui ajoute des contraintes supplémentaires.
    \item \textbf{Dynamique} : les canaux radio varient dans le temps (fading, mobilité).
\end{itemize}

\subsection{Concepts mathématiques fondamentaux}

Avant de présenter les approches existantes et notre contribution, il est nécessaire d'introduire certains concepts mathématiques qui seront utilisés tout au long de cet article. Ces concepts proviennent de la théorie des jeux et de l'optimisation, et constituent le cadre théorique sur lequel repose l'approche décentralisée classique.

\subsubsection{Théorie des jeux et équilibre de Nash}

La théorie des jeux fournit un cadre naturel pour modéliser les interactions stratégiques entre des décideurs autonomes. Dans notre contexte, chaque utilisateur du réseau est considéré comme un \textit{joueur} cherchant à maximiser sa propre efficacité énergétique. Formellement, un jeu non-coopératif est défini par un triplet $\mathcal{G} = \{\mathcal{K}, \{\mathcal{S}_k\}_{k\in\mathcal{K}}, \{u_k\}_{k\in\mathcal{K}}\}$, où $\mathcal{K} = \{1,\dots,K\}$ représente l'ensemble des joueurs, $\mathcal{S}_k$ l'ensemble des stratégies disponibles pour le joueur $k$, et $u_k : \mathcal{S}_1 \times \cdots \times \mathcal{S}_K \to \mathbb{R}$ sa fonction d'utilité.

Le concept central de la théorie des jeux non-coopératifs est l'\textit{équilibre de Nash}. Un profil de stratégies $\mathbf{s}^* = (s_1^*, \dots, s_K^*)$ constitue un équilibre de Nash si aucun joueur n'a intérêt à modifier unilatéralement sa stratégie, c'est-à-dire :
\[
u_k(s_k^*, \mathbf{s}_{-k}^*) \geq u_k(s_k, \mathbf{s}_{-k}^*) \quad \forall s_k \in \mathcal{S}_k, \forall k \in \mathcal{K},
\]
où $\mathbf{s}_{-k}^*$ désigne les stratégies de tous les joueurs excepté $k$. À l'équilibre, chaque joueur adopte la meilleure réponse possible compte tenu des choix des autres.

\subsubsection{Jeux généralisés et équilibre de Nash généralisé}

Dans notre problème de contrôle de puissance, une difficulté supplémentaire apparaît : l'ensemble des stratégies réalisables pour un utilisateur dépend des stratégies des autres. En effet, la contrainte de qualité de service (QoS) impose un SINR minimum, lequel dépend des interférences générées par les autres utilisateurs. On parle alors de \textit{jeu généralisé}, où l'ensemble des stratégies s'écrit $\mathcal{S}_k = \mathcal{S}_k(\mathbf{s}_{-k})$.

Dans ce cadre, la notion d'équilibre doit être adaptée. Un \textit{équilibre de Nash généralisé} (GNE) est un profil $\mathbf{p}^*$ tel que pour tout utilisateur $k$ :
\[
\mathbf{p}_k^* \in \mathcal{P}_k(\mathbf{p}_{-k}^*) \quad \text{et} \quad \eta_k(\mathbf{p}_k^*, \mathbf{p}_{-k}^*) \geq \eta_k(\mathbf{p}_k, \mathbf{p}_{-k}^*) \quad \forall \mathbf{p}_k \in \mathcal{P}_k(\mathbf{p}_{-k}^*),
\]
où $\mathcal{P}_k(\mathbf{p}_{-k})$ représente l'ensemble des puissances réalisables pour l'utilisateur $k$ étant donné les puissances des autres.

\subsubsection{Fonction standard et convergence}

Pour garantir l'unicité de l'équilibre et la convergence des algorithmes itératifs, Yates \cite{Yates1995} a introduit le concept de \textit{fonction standard}. Une fonction $f: \mathbb{R}_+^K \to \mathbb{R}_+^K$ est dite standard si elle satisfait trois propriétés essentielles :
\begin{itemize}
    \item \textbf{Positivité} : $f(\mathbf{p}) > 0$ pour tout $\mathbf{p} \geq 0$,
    \item \textbf{Monotonie} : si $\mathbf{p} \leq \mathbf{p}'$, alors $f(\mathbf{p}) \leq f(\mathbf{p}')$,
    \item \textbf{Scalabilité} : pour tout $\alpha > 1$, $\alpha f(\mathbf{p}) > f(\alpha \mathbf{p})$.
\end{itemize}
Si la fonction de meilleure réponse d'un jeu est standard, alors le jeu possède un unique équilibre et la mise à jour synchrone $p_k^{[t+1]} = B_k(\mathbf{p}_{-k}^{[t]})$ converge vers cet équilibre quel que soit le point de départ.

\subsubsection{Optimalité de Pareto}

L'équilibre de Nash n'est pas nécessairement optimal du point de vue collectif. La notion d'\textit{optimalité de Pareto} permet de caractériser les situations où aucune amélioration n'est possible sans détériorer la situation d'au moins un joueur. Formellement, un profil $\mathbf{p}$ est Pareto-optimal s'il n'existe pas d'autre profil $\mathbf{p}'$ tel que $\eta_k(\mathbf{p}') \geq \eta_k(\mathbf{p})$ pour tout $k$, avec au moins une inégalité stricte. Cette notion est importante car elle met en évidence la possible inefficacité des équilibres non-coopératifs.

\subsubsection{Conditions d'optimalité de Karush-Kuhn-Tucker (KKT)}

En optimisation sous contraintes, les \textit{conditions de Karush-Kuhn-Tucker} (KKT) jouent un rôle fondamental. Pour un problème de maximisation avec contraintes d'inégalité $g_i(\mathbf{x}) \leq 0$ et d'égalité $h_j(\mathbf{x}) = 0$, un point $\mathbf{x}^*$ satisfait les conditions KKT s'il existe des multiplicateurs $\lambda_i \geq 0$ et $\mu_j$ tels que :
\begin{align}
\nabla f(\mathbf{x}^*) + \sum_{i=1}^m \lambda_i \nabla g_i(\mathbf{x}^*) + \sum_{j=1}^p \mu_j \nabla h_j(\mathbf{x}^*) &= 0, \\
\lambda_i g_i(\mathbf{x}^*) &= 0 \quad \forall i, \\
g_i(\mathbf{x}^*) &\leq 0 \quad \forall i, \\
h_j(\mathbf{x}^*) &= 0 \quad \forall j.
\end{align}
Ces conditions sont nécessaires pour qu'un point soit un optimum local (sous certaines hypothèses de régularité). Dans le contexte de l'optimisation de l'EE, les algorithmes développés par Zappone et al. \cite{Zappone2021} convergent vers des points satisfaisant ces conditions.\subsection{Travaux connexes}

L'optimisation de l'efficacité énergétique dans les réseaux sans fil a fait l'objet de nombreux travaux. On peut distinguer deux grandes familles d'approches : les approches centralisées et les approches décentralisées.

\subsubsection{Approches centralisées}

Les approches centralisées supposent l'existence d'une entité centrale disposant d'une connaissance parfaite de tous les canaux et de toutes les puissances. Cette entité résout globalement le problème d'optimisation. Les premiers travaux dans ce domaine utilisaient la programmation fractionnaire \cite{Dinkelbach1967}, qui permet de traiter les problèmes d'optimisation de ratios.

Cependant, comme le soulignent Zappone et al. \cite{Zappone2021}, la présence d'interférences rend le numérateur de l'EE non-concave, ce qui empêche l'application directe de la programmation fractionnaire classique. Des solutions basées sur des schémas de transmission orthogonaux \cite{Ng2012, Xu2014} ou sur l'annulation d'interférences \cite{Xu2013} ont été proposées, mais elles conduisent à un faible facteur de réutilisation des ressources.

D'autres approches utilisent des techniques heuristiques comme l'optimisation alternée \cite{Du2014, He2013, He2014}, mais sans garantie de convergence. Venturino et al. \cite{Venturino2015} ont proposé un cadre unifié pour l'optimisation de l'EE dans les réseaux OFDMA.

\subsubsection{Approches décentralisées}

Les approches décentralisées modélisent le réseau comme un jeu non-coopératif où chaque utilisateur maximise sa propre EE. Miao et al. \cite{Miao2011} ont étudié le problème de l'équilibre de Nash dans ce contexte, sans contraintes de débit. Scutari et al. \cite{Scutari2008} ont proposé des stratégies de précodage optimales basées sur la théorie des jeux.

Zappone et al. \cite{Zappone2013} ont étendu ces travaux aux réseaux avec relais, tandis que Stupia et al. \cite{Stupia2015} ont utilisé une approche par inéquations variationnelles. Bacci et al. \cite{Bacci2015} ont introduit des contraintes de QoS dans le jeu, rendant le problème plus réaliste mais aussi plus complexe.

L'article de Zappone et al. \cite{Zappone2021} constitue une référence majeure dans ce domaine. Il propose un cadre unifié pour l'optimisation centralisée et décentralisée de l'EE avec contraintes de QoS et une expression généralisée du SINR englobant plusieurs technologies 5G (Massive MIMO avec imperfections matérielles, réseaux avec relais, etc.). Les auteurs développent :
\begin{itemize}
    \item Des conditions de faisabilité sous forme fermée (Lemme 1),
    \item Des algorithmes centralisés convergeant vers des points KKT,
    \item Une modélisation du problème décentralisé comme un jeu généralisé,
    \item Une preuve d'existence et d'unicité du GNE (Propositions 3 et 4),
    \item Un algorithme distribué basé sur les meilleures réponses,
    \item Une extension au cas multi-blocs de ressources.
\end{itemize}

\subsubsection{Limites des approches existantes}

Malgré leur rigueur mathématique, les approches de Zappone et al. présentent plusieurs limitations :

\begin{enumerate}
    \item \textbf{Nécessité d'un modèle parfait} : les algorithmes reposent sur la connaissance exacte des paramètres $\alpha_k$, $\phi_k$, $\omega_{kj}$. En pratique, ces paramètres sont estimés avec des erreurs.
    \item \textbf{Stationnarité} : la convergence des algorithmes suppose que les canaux varient lentement. Dans un environnement dynamique (mobilité, fading rapide), l'algorithme n'a pas le temps de converger.
    \item \textbf{Absence d'apprentissage} : les performances ne s'améliorent pas avec l'expérience. Chaque nouvelle situation est traitée indépendamment des précédentes.
    \item \textbf{Sensibilité aux erreurs} : les approximations et erreurs d'estimation peuvent dégrader significativement les performances.
\end{enumerate}

\subsection{Contributions de cet article}

Face à ces limitations, nous proposons une approche alternative basée sur l'apprentissage par renforcement profond (Deep Reinforcement Learning). Nos principales contributions sont :

\begin{enumerate}
    \item Une formulation du problème de contrôle de puissance comme un Processus de Décision Markovien (MDP) avec un état composé d'observations locales, des actions discrètes (niveaux de puissance) et une récompense combinant EE, pénalité QoS et pénalité de puissance.
    \item L'utilisation d'un réseau de neurones profond (DQN) pour apprendre une politique de contrôle de puissance optimale, sans nécessiter de modèle mathématique explicite du canal.
    \item Une architecture de réseau adaptée au problème (8 entrées, couches cachées [256,512,256], sortie 16 actions).
    \item Des techniques de stabilisation avancées (experience replay, target network, Double DQN, Huber loss, gradient clipping, $\varepsilon$-greedy décroissant).
    \item Une évaluation expérimentale montrant un gain significatif de l'EE par rapport aux approches classiques.
\end{enumerate}

\subsection{Organisation de l'article}

La suite de cet article est organisée comme suit. La section 2 présente en détail l'approche décentralisée classique de Zappone et al., avec toutes les dérivations mathématiques. La section 3 décrit notre approche DQN, incluant la modélisation MDP, l'architecture du réseau et les techniques d'entraînement. La section 4 présente les résultats expérimentaux et les compare aux approches classiques. Enfin, la section 5 conclut l'article et discute des perspectives futures.

%==============================================================================
\section{L'approche décentralisée classique : jeu généralisé et équilibre de Nash}
%==============================================================================

[Le contenu détaillé de cette section sera développé dans la prochaine réponse]

% ============================================================
% Chapitres DQN + Comparaison — Rapport PFE
% Contrôle de Puissance Décentralisé via Deep Q-Network
% pour Réseaux 5G Massive MIMO
% ============================================================
% À inclure dans le document principal via \documentclass[conference]{IEEEtran}
\usepackage[utf8]{inputenc}
\usepackage[T1]{fontenc}
\usepackage[french]{babel}
\usepackage{amsmath}
\usepackage{amsfonts}
\usepackage{amssymb}
\usepackage{amsthm}
\usepackage{graphicx}
\usepackage{mathrsfs}
\usepackage{cite}
% Packages ajoutés pour corriger les incohérences
\usepackage{algorithm2e}
\usepackage{booktabs}
\usepackage[caption=false,font=footnotesize]{subfig}
\usepackage{pifont}
\usepackage{multirow}
\usepackage{tabularx}
\usepackage{xcolor}
\usepackage{float}
\usepackage{stfloats}
\usepackage{hyperref}

% Environnements pour théorèmes et définitions
\newtheorem{definition}{Définition}
\newtheorem{theorem}{Théorème}
\newtheorem{lemma}{Lemme}
\newtheorem{proposition}{Proposition}

\title{Optimisation de l'Efficacité Énergétique dans les Réseaux 5G : Une Approche par Apprentissage par Renforcement Profond}
\author{
    \IEEEauthorblockN{Nada, Youssef IFQIR}
    \IEEEauthorblockA{Université ... , Faculté des Sciences et Techniques\\
    Année académique 2025--2026\\[0.5em]
    \includegraphics[height=1.8cm]{results/figures/Logofst.jpg}}}

\begin{document}

\maketitle

\begin{abstract}
L'essor des réseaux 5G s'accompagne d'une augmentation exponentielle de la consommation énergétique, ce qui rend crucial le développement de stratégies de contrôle de puissance énergétiquement efficaces. Cet article propose une approche novatrice basée sur l'apprentissage par renforcement profond (Deep Q-Network, DQN) pour le contrôle de puissance décentralisé dans les réseaux 5G, en vue de maximiser l'efficacité énergétique (EE) tout en respectant des contraintes de qualité de service (QoS). Nous commençons par une analyse approfondie de l'approche classique de Zappone et al. (2021), qui modélise le problème comme un jeu non-coopératif généralisé et propose un algorithme distribué fondé sur la théorie des jeux et l'équilibre de Nash généralisé (GNE). Nous détaillons les fondements mathématiques de cette approche : conditions de faisabilité (lemme 1), existence et unicité du GNE (propositions 3 et 4), calcul de la meilleure réponse via l'algorithme de Dinkelbach, et algorithme distribué itératif. Nous mettons en évidence les limites de cette méthode : nécessité d'un modèle parfait du canal, hypothèse de stationnarité, absence d'apprentissage et sensibilité aux erreurs d'estimation.

Pour surmonter ces limitations, nous reformulons le problème de contrôle de puissance comme un Processus de Décision Markovien (MDP) dans lequel chaque utilisateur est un agent intelligent. L'état de l'agent est composé d'observations locales (puissance courante, SINR, estimation du canal, historique des récompenses, etc.), l'action est le choix d'un niveau de puissance parmi 16 valeurs discrètes, et la récompense est une combinaison pondérée de l'EE instantanée, d'une pénalité en cas de violation de la QoS et d'une pénalité sur la puissance utilisée. Nous implémentons un réseau de neurones profond avec une architecture à trois couches cachées (256, 512, 256 neurones) et utilisons des techniques avancées de stabilisation de l'apprentissage : experience replay, réseau cible, Double DQN, Huber loss, gradient clipping et stratégie d'exploration $\varepsilon$-greedy décroissante.

Les expériences menées dans un environnement de simulation réaliste montrent que notre agent DQN atteint une EE moyenne de 515 000 bit/J, soit un gain de plus de 130 \% par rapport aux algorithmes distribués classiques, tout en maintenant un taux de violation de la QoS inférieur à 0,03 \%. De plus, l'approche DQN s'adapte dynamiquement aux variations des conditions radio et améliore ses performances avec l'expérience, contrairement aux méthodes analytiques. Nous visualisons également la topologie du réseau et la convergence des décisions à l'aide de la base de données graphe Neo4j, confirmant la stabilité et la cohérence de la politique apprise. Ces résultats démontrent la pertinence et l'efficacité de l'apprentissage par renforcement profond pour le contrôle de puissance dans les réseaux 5G, ouvrant la voie à des systèmes autonomes, adaptatifs et économes en énergie.
\end{abstract}

\begin{IEEEkeywords}
Réseaux 5G, Efficacité énergétique (EE), Contrôle de puissance décentralisé, Allocation de ressources, Contraintes de qualité de service (QoS), Théorie des jeux, Jeu non-coopératif, Jeu généralisé, Équilibre de Nash généralisé (GNE), Fonction standard, Algorithme de Dinkelbach, Meilleure réponse, Programmation fractionnaire, Apprentissage par renforcement profond (DRL), Deep Q-Network (DQN), Processus de Décision Markovien (MDP), Approche sans modèle, Adaptation dynamique, Environnement non stationnaire.
\end{IEEEkeywords}

%==============================================================================
\section{Introduction}
%==============================================================================

\subsection{Contexte et motivation}

Les réseaux de cinquième génération (5G) représentent une rupture technologique majeure dans l'histoire des communications sans fil. Avec des débits annoncés atteignant 10 Gbit/s, une latence inférieure à 1 ms et une densité de connexion pouvant atteindre 1 million de dispositifs par kilomètre carré \cite{3GPP}, la 5G doit répondre à des exigences sans précédent. Cependant, cette augmentation spectaculaire des performances s'accompagne d'une augmentation tout aussi spectaculaire de la consommation énergétique.

Selon \cite{Fehske2011}, le secteur des technologies de l'information et de la communication (TIC) représente déjà environ 5\% des émissions mondiales de CO$_2$, et ce pourcentage est en croissance rapide. Avec le déploiement massif des réseaux 5G, la situation risque de s'aggraver si des mesures ne sont pas prises. Il est estimé que le nombre de dispositifs connectés atteindra 50 milliards d'ici 2025 \cite{Ericsson2011}, et que les débits devront être multipliés par 1000 par rapport aux réseaux 4G \cite{Qualcomm}.

Une solution naïve consisterait à augmenter proportionnellement les puissances d'émission pour atteindre ces débits. Cependant, cette approche est non seulement irréaliste d'un point de vue énergétique, mais aussi inacceptable pour des raisons environnementales (émissions de gaz à effet de serre) et sanitaires (pollution électromagnétique). La seule solution viable est d'augmenter l'\textbf{efficacité énergétique} (EE) des réseaux, c'est-à-dire la quantité de données transmise par unité d'énergie consommée.

L'EE est généralement mesurée en bits par joule (bit/J) et peut être définie au niveau de l'utilisateur ou au niveau du réseau. L'objectif fondamental est de "transmettre plus de données en consommant moins d'énergie", ce qui nécessite des stratégies de contrôle de puissance sophistiquées.

\subsection{Le problème du contrôle de puissance}

Dans un réseau sans fil, chaque utilisateur doit choisir sa puissance d'émission. Ce choix apparemment simple a des conséquences complexes :

\begin{itemize}
    \item Il détermine le \textbf{débit} de l'utilisateur via le rapport signal-sur-bruit-plus-interférence (SINR).
    \item Il influence les \textbf{interférences} subies par les autres utilisateurs.
    \item Il affecte la \textbf{consommation énergétique} globale du réseau et de l'utilisateur.
\end{itemize}

Mathématiquement, pour un utilisateur $k$ utilisant $N$ blocs de ressources, le SINR sur le bloc $n$ s'exprime généralement par :

\begin{equation} \label{eq:sinr_general}
\gamma_{k,n} = \frac{\alpha_{k,n} p_{k,n}}{\sigma_{k,n}^2 + \phi_{k,n} p_{k,n} + \sum_{j \neq k} \omega_{k,j,n} p_{j,n}}
\end{equation}

où $p_{k,n}$ est la puissance d'émission de l'utilisateur $k$ sur le bloc $n$, $\alpha_{k,n}$ représente le gain du canal direct, $\phi_{k,n}$ modélise l'auto-interférence (due aux imperfections matérielles), $\omega_{k,j,n}$ est le coefficient d'interférence de l'utilisateur $j$ sur l'utilisateur $k$, et $\sigma_{k,n}^2$ est la puissance du bruit thermique.

L'efficacité énergétique individuelle de l'utilisateur $k$ est alors définie comme le rapport entre son débit total et sa puissance totale consommée :

\begin{equation} \label{eq:ee}
\eta_k = \frac{\displaystyle\sum_{n=1}^{N} B \log_2(1 + \gamma_{k,n})}{p_{c,k} + \displaystyle\sum_{n=1}^{N} p_{k,n}}
\end{equation}

où $B$ est la bande passante par bloc de ressources et $p_{c,k}$ est la puissance de circuit (consommation fixe).

Ce problème d'optimisation présente plusieurs difficultés majeures :
\begin{itemize}
    \item \textbf{Non-convexité} : la fonction objectif $\eta_k$ est non-convexe en raison de la présence du logarithme et du rapport.
    \item \textbf{Couplage fort} : les décisions des utilisateurs sont interdépendantes via les termes d'interférence $\sum_{j \neq k} \omega_{k,j,n} p_{j,n}$.
    \item \textbf{Contraintes de QoS} : chaque utilisateur doit généralement respecter un débit minimum, ce qui ajoute des contraintes supplémentaires.
    \item \textbf{Dynamique} : les canaux radio varient dans le temps (fading, mobilité).
\end{itemize}

\subsection{Concepts mathématiques fondamentaux}

Avant de présenter les approches existantes et notre contribution, il est nécessaire d'introduire certains concepts mathématiques qui seront utilisés tout au long de cet article. Ces concepts proviennent de la théorie des jeux et de l'optimisation, et constituent le cadre théorique sur lequel repose l'approche décentralisée classique.

\subsubsection{Théorie des jeux et équilibre de Nash}

La théorie des jeux fournit un cadre naturel pour modéliser les interactions stratégiques entre des décideurs autonomes. Dans notre contexte, chaque utilisateur du réseau est considéré comme un \textit{joueur} cherchant à maximiser sa propre efficacité énergétique. Formellement, un jeu non-coopératif est défini par un triplet $\mathcal{G} = \{\mathcal{K}, \{\mathcal{S}_k\}_{k\in\mathcal{K}}, \{u_k\}_{k\in\mathcal{K}}\}$, où $\mathcal{K} = \{1,\dots,K\}$ représente l'ensemble des joueurs, $\mathcal{S}_k$ l'ensemble des stratégies disponibles pour le joueur $k$, et $u_k : \mathcal{S}_1 \times \cdots \times \mathcal{S}_K \to \mathbb{R}$ sa fonction d'utilité.

Le concept central de la théorie des jeux non-coopératifs est l'\textit{équilibre de Nash}. Un profil de stratégies $\mathbf{s}^* = (s_1^*, \dots, s_K^*)$ constitue un équilibre de Nash si aucun joueur n'a intérêt à modifier unilatéralement sa stratégie, c'est-à-dire :
\[
u_k(s_k^*, \mathbf{s}_{-k}^*) \geq u_k(s_k, \mathbf{s}_{-k}^*) \quad \forall s_k \in \mathcal{S}_k, \forall k \in \mathcal{K},
\]
où $\mathbf{s}_{-k}^*$ désigne les stratégies de tous les joueurs excepté $k$. À l'équilibre, chaque joueur adopte la meilleure réponse possible compte tenu des choix des autres.

\subsubsection{Jeux généralisés et équilibre de Nash généralisé}

Dans notre problème de contrôle de puissance, une difficulté supplémentaire apparaît : l'ensemble des stratégies réalisables pour un utilisateur dépend des stratégies des autres. En effet, la contrainte de qualité de service (QoS) impose un SINR minimum, lequel dépend des interférences générées par les autres utilisateurs. On parle alors de \textit{jeu généralisé}, où l'ensemble des stratégies s'écrit $\mathcal{S}_k = \mathcal{S}_k(\mathbf{s}_{-k})$.

Dans ce cadre, la notion d'équilibre doit être adaptée. Un \textit{équilibre de Nash généralisé} (GNE) est un profil $\mathbf{p}^*$ tel que pour tout utilisateur $k$ :
\[
\mathbf{p}_k^* \in \mathcal{P}_k(\mathbf{p}_{-k}^*) \quad \text{et} \quad \eta_k(\mathbf{p}_k^*, \mathbf{p}_{-k}^*) \geq \eta_k(\mathbf{p}_k, \mathbf{p}_{-k}^*) \quad \forall \mathbf{p}_k \in \mathcal{P}_k(\mathbf{p}_{-k}^*),
\]
où $\mathcal{P}_k(\mathbf{p}_{-k})$ représente l'ensemble des puissances réalisables pour l'utilisateur $k$ étant donné les puissances des autres.

\subsubsection{Fonction standard et convergence}

Pour garantir l'unicité de l'équilibre et la convergence des algorithmes itératifs, Yates \cite{Yates1995} a introduit le concept de \textit{fonction standard}. Une fonction $f: \mathbb{R}_+^K \to \mathbb{R}_+^K$ est dite standard si elle satisfait trois propriétés essentielles :
\begin{itemize}
    \item \textbf{Positivité} : $f(\mathbf{p}) > 0$ pour tout $\mathbf{p} \geq 0$,
    \item \textbf{Monotonie} : si $\mathbf{p} \leq \mathbf{p}'$, alors $f(\mathbf{p}) \leq f(\mathbf{p}')$,
    \item \textbf{Scalabilité} : pour tout $\alpha > 1$, $\alpha f(\mathbf{p}) > f(\alpha \mathbf{p})$.
\end{itemize}
Si la fonction de meilleure réponse d'un jeu est standard, alors le jeu possède un unique équilibre et la mise à jour synchrone $p_k^{[t+1]} = B_k(\mathbf{p}_{-k}^{[t]})$ converge vers cet équilibre quel que soit le point de départ.

\subsubsection{Optimalité de Pareto}

L'équilibre de Nash n'est pas nécessairement optimal du point de vue collectif. La notion d'\textit{optimalité de Pareto} permet de caractériser les situations où aucune amélioration n'est possible sans détériorer la situation d'au moins un joueur. Formellement, un profil $\mathbf{p}$ est Pareto-optimal s'il n'existe pas d'autre profil $\mathbf{p}'$ tel que $\eta_k(\mathbf{p}') \geq \eta_k(\mathbf{p})$ pour tout $k$, avec au moins une inégalité stricte. Cette notion est importante car elle met en évidence la possible inefficacité des équilibres non-coopératifs.

\subsubsection{Conditions d'optimalité de Karush-Kuhn-Tucker (KKT)}

En optimisation sous contraintes, les \textit{conditions de Karush-Kuhn-Tucker} (KKT) jouent un rôle fondamental. Pour un problème de maximisation avec contraintes d'inégalité $g_i(\mathbf{x}) \leq 0$ et d'égalité $h_j(\mathbf{x}) = 0$, un point $\mathbf{x}^*$ satisfait les conditions KKT s'il existe des multiplicateurs $\lambda_i \geq 0$ et $\mu_j$ tels que :
\begin{align}
\nabla f(\mathbf{x}^*) + \sum_{i=1}^m \lambda_i \nabla g_i(\mathbf{x}^*) + \sum_{j=1}^p \mu_j \nabla h_j(\mathbf{x}^*) &= 0, \\
\lambda_i g_i(\mathbf{x}^*) &= 0 \quad \forall i, \\
g_i(\mathbf{x}^*) &\leq 0 \quad \forall i, \\
h_j(\mathbf{x}^*) &= 0 \quad \forall j.
\end{align}
Ces conditions sont nécessaires pour qu'un point soit un optimum local (sous certaines hypothèses de régularité). Dans le contexte de l'optimisation de l'EE, les algorithmes développés par Zappone et al. \cite{Zappone2021} convergent vers des points satisfaisant ces conditions.\subsection{Travaux connexes}

L'optimisation de l'efficacité énergétique dans les réseaux sans fil a fait l'objet de nombreux travaux. On peut distinguer deux grandes familles d'approches : les approches centralisées et les approches décentralisées.

\subsubsection{Approches centralisées}

Les approches centralisées supposent l'existence d'une entité centrale disposant d'une connaissance parfaite de tous les canaux et de toutes les puissances. Cette entité résout globalement le problème d'optimisation. Les premiers travaux dans ce domaine utilisaient la programmation fractionnaire \cite{Dinkelbach1967}, qui permet de traiter les problèmes d'optimisation de ratios.

Cependant, comme le soulignent Zappone et al. \cite{Zappone2021}, la présence d'interférences rend le numérateur de l'EE non-concave, ce qui empêche l'application directe de la programmation fractionnaire classique. Des solutions basées sur des schémas de transmission orthogonaux \cite{Ng2012, Xu2014} ou sur l'annulation d'interférences \cite{Xu2013} ont été proposées, mais elles conduisent à un faible facteur de réutilisation des ressources.

D'autres approches utilisent des techniques heuristiques comme l'optimisation alternée \cite{Du2014, He2013, He2014}, mais sans garantie de convergence. Venturino et al. \cite{Venturino2015} ont proposé un cadre unifié pour l'optimisation de l'EE dans les réseaux OFDMA.

\subsubsection{Approches décentralisées}

Les approches décentralisées modélisent le réseau comme un jeu non-coopératif où chaque utilisateur maximise sa propre EE. Miao et al. \cite{Miao2011} ont étudié le problème de l'équilibre de Nash dans ce contexte, sans contraintes de débit. Scutari et al. \cite{Scutari2008} ont proposé des stratégies de précodage optimales basées sur la théorie des jeux.

Zappone et al. \cite{Zappone2013} ont étendu ces travaux aux réseaux avec relais, tandis que Stupia et al. \cite{Stupia2015} ont utilisé une approche par inéquations variationnelles. Bacci et al. \cite{Bacci2015} ont introduit des contraintes de QoS dans le jeu, rendant le problème plus réaliste mais aussi plus complexe.

L'article de Zappone et al. \cite{Zappone2021} constitue une référence majeure dans ce domaine. Il propose un cadre unifié pour l'optimisation centralisée et décentralisée de l'EE avec contraintes de QoS et une expression généralisée du SINR englobant plusieurs technologies 5G (Massive MIMO avec imperfections matérielles, réseaux avec relais, etc.). Les auteurs développent :
\begin{itemize}
    \item Des conditions de faisabilité sous forme fermée (Lemme 1),
    \item Des algorithmes centralisés convergeant vers des points KKT,
    \item Une modélisation du problème décentralisé comme un jeu généralisé,
    \item Une preuve d'existence et d'unicité du GNE (Propositions 3 et 4),
    \item Un algorithme distribué basé sur les meilleures réponses,
    \item Une extension au cas multi-blocs de ressources.
\end{itemize}

\subsubsection{Limites des approches existantes}

Malgré leur rigueur mathématique, les approches de Zappone et al. présentent plusieurs limitations :

\begin{enumerate}
    \item \textbf{Nécessité d'un modèle parfait} : les algorithmes reposent sur la connaissance exacte des paramètres $\alpha_k$, $\phi_k$, $\omega_{kj}$. En pratique, ces paramètres sont estimés avec des erreurs.
    \item \textbf{Stationnarité} : la convergence des algorithmes suppose que les canaux varient lentement. Dans un environnement dynamique (mobilité, fading rapide), l'algorithme n'a pas le temps de converger.
    \item \textbf{Absence d'apprentissage} : les performances ne s'améliorent pas avec l'expérience. Chaque nouvelle situation est traitée indépendamment des précédentes.
    \item \textbf{Sensibilité aux erreurs} : les approximations et erreurs d'estimation peuvent dégrader significativement les performances.
\end{enumerate}

\subsection{Contributions de cet article}

Face à ces limitations, nous proposons une approche alternative basée sur l'apprentissage par renforcement profond (Deep Reinforcement Learning). Nos principales contributions sont :

\begin{enumerate}
    \item Une formulation du problème de contrôle de puissance comme un Processus de Décision Markovien (MDP) avec un état composé d'observations locales, des actions discrètes (niveaux de puissance) et une récompense combinant EE, pénalité QoS et pénalité de puissance.
    \item L'utilisation d'un réseau de neurones profond (DQN) pour apprendre une politique de contrôle de puissance optimale, sans nécessiter de modèle mathématique explicite du canal.
    \item Une architecture de réseau adaptée au problème (8 entrées, couches cachées [256,512,256], sortie 16 actions).
    \item Des techniques de stabilisation avancées (experience replay, target network, Double DQN, Huber loss, gradient clipping, $\varepsilon$-greedy décroissant).
    \item Une évaluation expérimentale montrant un gain significatif de l'EE par rapport aux approches classiques.
\end{enumerate}

\subsection{Organisation de l'article}

La suite de cet article est organisée comme suit. La section 2 présente en détail l'approche décentralisée classique de Zappone et al., avec toutes les dérivations mathématiques. La section 3 décrit notre approche DQN, incluant la modélisation MDP, l'architecture du réseau et les techniques d'entraînement. La section 4 présente les résultats expérimentaux et les compare aux approches classiques. Enfin, la section 5 conclut l'article et discute des perspectives futures.

%==============================================================================
\section{L'approche décentralisée classique : jeu généralisé et équilibre de Nash}
%==============================================================================

[Le contenu détaillé de cette section sera développé dans la prochaine réponse]

% ============================================================
% Chapitres DQN + Comparaison — Rapport PFE
% Contrôle de Puissance Décentralisé via Deep Q-Network
% pour Réseaux 5G Massive MIMO
% ============================================================
% À inclure dans le document principal via \documentclass[conference]{IEEEtran}
\usepackage[utf8]{inputenc}
\usepackage[T1]{fontenc}
\usepackage[french]{babel}
\usepackage{amsmath}
\usepackage{amsfonts}
\usepackage{amssymb}
\usepackage{amsthm}
\usepackage{graphicx}
\usepackage{mathrsfs}
\usepackage{cite}
% Packages ajoutés pour corriger les incohérences
\usepackage{algorithm2e}
\usepackage{booktabs}
\usepackage[caption=false,font=footnotesize]{subfig}
\usepackage{pifont}
\usepackage{multirow}
\usepackage{tabularx}
\usepackage{xcolor}
\usepackage{float}
\usepackage{stfloats}
\usepackage{hyperref}

% Environnements pour théorèmes et définitions
\newtheorem{definition}{Définition}
\newtheorem{theorem}{Théorème}
\newtheorem{lemma}{Lemme}
\newtheorem{proposition}{Proposition}

\title{Optimisation de l'Efficacité Énergétique dans les Réseaux 5G : Une Approche par Apprentissage par Renforcement Profond}
\author{
    \IEEEauthorblockN{Nada, Youssef IFQIR}
    \IEEEauthorblockA{Université ... , Faculté des Sciences et Techniques\\
    Année académique 2025--2026\\[0.5em]
    \includegraphics[height=1.8cm]{results/figures/Logofst.jpg}}}

\begin{document}

\maketitle

\begin{abstract}
L'essor des réseaux 5G s'accompagne d'une augmentation exponentielle de la consommation énergétique, ce qui rend crucial le développement de stratégies de contrôle de puissance énergétiquement efficaces. Cet article propose une approche novatrice basée sur l'apprentissage par renforcement profond (Deep Q-Network, DQN) pour le contrôle de puissance décentralisé dans les réseaux 5G, en vue de maximiser l'efficacité énergétique (EE) tout en respectant des contraintes de qualité de service (QoS). Nous commençons par une analyse approfondie de l'approche classique de Zappone et al. (2021), qui modélise le problème comme un jeu non-coopératif généralisé et propose un algorithme distribué fondé sur la théorie des jeux et l'équilibre de Nash généralisé (GNE). Nous détaillons les fondements mathématiques de cette approche : conditions de faisabilité (lemme 1), existence et unicité du GNE (propositions 3 et 4), calcul de la meilleure réponse via l'algorithme de Dinkelbach, et algorithme distribué itératif. Nous mettons en évidence les limites de cette méthode : nécessité d'un modèle parfait du canal, hypothèse de stationnarité, absence d'apprentissage et sensibilité aux erreurs d'estimation.

Pour surmonter ces limitations, nous reformulons le problème de contrôle de puissance comme un Processus de Décision Markovien (MDP) dans lequel chaque utilisateur est un agent intelligent. L'état de l'agent est composé d'observations locales (puissance courante, SINR, estimation du canal, historique des récompenses, etc.), l'action est le choix d'un niveau de puissance parmi 16 valeurs discrètes, et la récompense est une combinaison pondérée de l'EE instantanée, d'une pénalité en cas de violation de la QoS et d'une pénalité sur la puissance utilisée. Nous implémentons un réseau de neurones profond avec une architecture à trois couches cachées (256, 512, 256 neurones) et utilisons des techniques avancées de stabilisation de l'apprentissage : experience replay, réseau cible, Double DQN, Huber loss, gradient clipping et stratégie d'exploration $\varepsilon$-greedy décroissante.

Les expériences menées dans un environnement de simulation réaliste montrent que notre agent DQN atteint une EE moyenne de 515 000 bit/J, soit un gain de plus de 130 \% par rapport aux algorithmes distribués classiques, tout en maintenant un taux de violation de la QoS inférieur à 0,03 \%. De plus, l'approche DQN s'adapte dynamiquement aux variations des conditions radio et améliore ses performances avec l'expérience, contrairement aux méthodes analytiques. Nous visualisons également la topologie du réseau et la convergence des décisions à l'aide de la base de données graphe Neo4j, confirmant la stabilité et la cohérence de la politique apprise. Ces résultats démontrent la pertinence et l'efficacité de l'apprentissage par renforcement profond pour le contrôle de puissance dans les réseaux 5G, ouvrant la voie à des systèmes autonomes, adaptatifs et économes en énergie.
\end{abstract}

\begin{IEEEkeywords}
Réseaux 5G, Efficacité énergétique (EE), Contrôle de puissance décentralisé, Allocation de ressources, Contraintes de qualité de service (QoS), Théorie des jeux, Jeu non-coopératif, Jeu généralisé, Équilibre de Nash généralisé (GNE), Fonction standard, Algorithme de Dinkelbach, Meilleure réponse, Programmation fractionnaire, Apprentissage par renforcement profond (DRL), Deep Q-Network (DQN), Processus de Décision Markovien (MDP), Approche sans modèle, Adaptation dynamique, Environnement non stationnaire.
\end{IEEEkeywords}

%==============================================================================
\section{Introduction}
%==============================================================================

\subsection{Contexte et motivation}

Les réseaux de cinquième génération (5G) représentent une rupture technologique majeure dans l'histoire des communications sans fil. Avec des débits annoncés atteignant 10 Gbit/s, une latence inférieure à 1 ms et une densité de connexion pouvant atteindre 1 million de dispositifs par kilomètre carré \cite{3GPP}, la 5G doit répondre à des exigences sans précédent. Cependant, cette augmentation spectaculaire des performances s'accompagne d'une augmentation tout aussi spectaculaire de la consommation énergétique.

Selon \cite{Fehske2011}, le secteur des technologies de l'information et de la communication (TIC) représente déjà environ 5\% des émissions mondiales de CO$_2$, et ce pourcentage est en croissance rapide. Avec le déploiement massif des réseaux 5G, la situation risque de s'aggraver si des mesures ne sont pas prises. Il est estimé que le nombre de dispositifs connectés atteindra 50 milliards d'ici 2025 \cite{Ericsson2011}, et que les débits devront être multipliés par 1000 par rapport aux réseaux 4G \cite{Qualcomm}.

Une solution naïve consisterait à augmenter proportionnellement les puissances d'émission pour atteindre ces débits. Cependant, cette approche est non seulement irréaliste d'un point de vue énergétique, mais aussi inacceptable pour des raisons environnementales (émissions de gaz à effet de serre) et sanitaires (pollution électromagnétique). La seule solution viable est d'augmenter l'\textbf{efficacité énergétique} (EE) des réseaux, c'est-à-dire la quantité de données transmise par unité d'énergie consommée.

L'EE est généralement mesurée en bits par joule (bit/J) et peut être définie au niveau de l'utilisateur ou au niveau du réseau. L'objectif fondamental est de "transmettre plus de données en consommant moins d'énergie", ce qui nécessite des stratégies de contrôle de puissance sophistiquées.

\subsection{Le problème du contrôle de puissance}

Dans un réseau sans fil, chaque utilisateur doit choisir sa puissance d'émission. Ce choix apparemment simple a des conséquences complexes :

\begin{itemize}
    \item Il détermine le \textbf{débit} de l'utilisateur via le rapport signal-sur-bruit-plus-interférence (SINR).
    \item Il influence les \textbf{interférences} subies par les autres utilisateurs.
    \item Il affecte la \textbf{consommation énergétique} globale du réseau et de l'utilisateur.
\end{itemize}

Mathématiquement, pour un utilisateur $k$ utilisant $N$ blocs de ressources, le SINR sur le bloc $n$ s'exprime généralement par :

\begin{equation} \label{eq:sinr_general}
\gamma_{k,n} = \frac{\alpha_{k,n} p_{k,n}}{\sigma_{k,n}^2 + \phi_{k,n} p_{k,n} + \sum_{j \neq k} \omega_{k,j,n} p_{j,n}}
\end{equation}

où $p_{k,n}$ est la puissance d'émission de l'utilisateur $k$ sur le bloc $n$, $\alpha_{k,n}$ représente le gain du canal direct, $\phi_{k,n}$ modélise l'auto-interférence (due aux imperfections matérielles), $\omega_{k,j,n}$ est le coefficient d'interférence de l'utilisateur $j$ sur l'utilisateur $k$, et $\sigma_{k,n}^2$ est la puissance du bruit thermique.

L'efficacité énergétique individuelle de l'utilisateur $k$ est alors définie comme le rapport entre son débit total et sa puissance totale consommée :

\begin{equation} \label{eq:ee}
\eta_k = \frac{\displaystyle\sum_{n=1}^{N} B \log_2(1 + \gamma_{k,n})}{p_{c,k} + \displaystyle\sum_{n=1}^{N} p_{k,n}}
\end{equation}

où $B$ est la bande passante par bloc de ressources et $p_{c,k}$ est la puissance de circuit (consommation fixe).

Ce problème d'optimisation présente plusieurs difficultés majeures :
\begin{itemize}
    \item \textbf{Non-convexité} : la fonction objectif $\eta_k$ est non-convexe en raison de la présence du logarithme et du rapport.
    \item \textbf{Couplage fort} : les décisions des utilisateurs sont interdépendantes via les termes d'interférence $\sum_{j \neq k} \omega_{k,j,n} p_{j,n}$.
    \item \textbf{Contraintes de QoS} : chaque utilisateur doit généralement respecter un débit minimum, ce qui ajoute des contraintes supplémentaires.
    \item \textbf{Dynamique} : les canaux radio varient dans le temps (fading, mobilité).
\end{itemize}

\subsection{Concepts mathématiques fondamentaux}

Avant de présenter les approches existantes et notre contribution, il est nécessaire d'introduire certains concepts mathématiques qui seront utilisés tout au long de cet article. Ces concepts proviennent de la théorie des jeux et de l'optimisation, et constituent le cadre théorique sur lequel repose l'approche décentralisée classique.

\subsubsection{Théorie des jeux et équilibre de Nash}

La théorie des jeux fournit un cadre naturel pour modéliser les interactions stratégiques entre des décideurs autonomes. Dans notre contexte, chaque utilisateur du réseau est considéré comme un \textit{joueur} cherchant à maximiser sa propre efficacité énergétique. Formellement, un jeu non-coopératif est défini par un triplet $\mathcal{G} = \{\mathcal{K}, \{\mathcal{S}_k\}_{k\in\mathcal{K}}, \{u_k\}_{k\in\mathcal{K}}\}$, où $\mathcal{K} = \{1,\dots,K\}$ représente l'ensemble des joueurs, $\mathcal{S}_k$ l'ensemble des stratégies disponibles pour le joueur $k$, et $u_k : \mathcal{S}_1 \times \cdots \times \mathcal{S}_K \to \mathbb{R}$ sa fonction d'utilité.

Le concept central de la théorie des jeux non-coopératifs est l'\textit{équilibre de Nash}. Un profil de stratégies $\mathbf{s}^* = (s_1^*, \dots, s_K^*)$ constitue un équilibre de Nash si aucun joueur n'a intérêt à modifier unilatéralement sa stratégie, c'est-à-dire :
\[
u_k(s_k^*, \mathbf{s}_{-k}^*) \geq u_k(s_k, \mathbf{s}_{-k}^*) \quad \forall s_k \in \mathcal{S}_k, \forall k \in \mathcal{K},
\]
où $\mathbf{s}_{-k}^*$ désigne les stratégies de tous les joueurs excepté $k$. À l'équilibre, chaque joueur adopte la meilleure réponse possible compte tenu des choix des autres.

\subsubsection{Jeux généralisés et équilibre de Nash généralisé}

Dans notre problème de contrôle de puissance, une difficulté supplémentaire apparaît : l'ensemble des stratégies réalisables pour un utilisateur dépend des stratégies des autres. En effet, la contrainte de qualité de service (QoS) impose un SINR minimum, lequel dépend des interférences générées par les autres utilisateurs. On parle alors de \textit{jeu généralisé}, où l'ensemble des stratégies s'écrit $\mathcal{S}_k = \mathcal{S}_k(\mathbf{s}_{-k})$.

Dans ce cadre, la notion d'équilibre doit être adaptée. Un \textit{équilibre de Nash généralisé} (GNE) est un profil $\mathbf{p}^*$ tel que pour tout utilisateur $k$ :
\[
\mathbf{p}_k^* \in \mathcal{P}_k(\mathbf{p}_{-k}^*) \quad \text{et} \quad \eta_k(\mathbf{p}_k^*, \mathbf{p}_{-k}^*) \geq \eta_k(\mathbf{p}_k, \mathbf{p}_{-k}^*) \quad \forall \mathbf{p}_k \in \mathcal{P}_k(\mathbf{p}_{-k}^*),
\]
où $\mathcal{P}_k(\mathbf{p}_{-k})$ représente l'ensemble des puissances réalisables pour l'utilisateur $k$ étant donné les puissances des autres.

\subsubsection{Fonction standard et convergence}

Pour garantir l'unicité de l'équilibre et la convergence des algorithmes itératifs, Yates \cite{Yates1995} a introduit le concept de \textit{fonction standard}. Une fonction $f: \mathbb{R}_+^K \to \mathbb{R}_+^K$ est dite standard si elle satisfait trois propriétés essentielles :
\begin{itemize}
    \item \textbf{Positivité} : $f(\mathbf{p}) > 0$ pour tout $\mathbf{p} \geq 0$,
    \item \textbf{Monotonie} : si $\mathbf{p} \leq \mathbf{p}'$, alors $f(\mathbf{p}) \leq f(\mathbf{p}')$,
    \item \textbf{Scalabilité} : pour tout $\alpha > 1$, $\alpha f(\mathbf{p}) > f(\alpha \mathbf{p})$.
\end{itemize}
Si la fonction de meilleure réponse d'un jeu est standard, alors le jeu possède un unique équilibre et la mise à jour synchrone $p_k^{[t+1]} = B_k(\mathbf{p}_{-k}^{[t]})$ converge vers cet équilibre quel que soit le point de départ.

\subsubsection{Optimalité de Pareto}

L'équilibre de Nash n'est pas nécessairement optimal du point de vue collectif. La notion d'\textit{optimalité de Pareto} permet de caractériser les situations où aucune amélioration n'est possible sans détériorer la situation d'au moins un joueur. Formellement, un profil $\mathbf{p}$ est Pareto-optimal s'il n'existe pas d'autre profil $\mathbf{p}'$ tel que $\eta_k(\mathbf{p}') \geq \eta_k(\mathbf{p})$ pour tout $k$, avec au moins une inégalité stricte. Cette notion est importante car elle met en évidence la possible inefficacité des équilibres non-coopératifs.

\subsubsection{Conditions d'optimalité de Karush-Kuhn-Tucker (KKT)}

En optimisation sous contraintes, les \textit{conditions de Karush-Kuhn-Tucker} (KKT) jouent un rôle fondamental. Pour un problème de maximisation avec contraintes d'inégalité $g_i(\mathbf{x}) \leq 0$ et d'égalité $h_j(\mathbf{x}) = 0$, un point $\mathbf{x}^*$ satisfait les conditions KKT s'il existe des multiplicateurs $\lambda_i \geq 0$ et $\mu_j$ tels que :
\begin{align}
\nabla f(\mathbf{x}^*) + \sum_{i=1}^m \lambda_i \nabla g_i(\mathbf{x}^*) + \sum_{j=1}^p \mu_j \nabla h_j(\mathbf{x}^*) &= 0, \\
\lambda_i g_i(\mathbf{x}^*) &= 0 \quad \forall i, \\
g_i(\mathbf{x}^*) &\leq 0 \quad \forall i, \\
h_j(\mathbf{x}^*) &= 0 \quad \forall j.
\end{align}
Ces conditions sont nécessaires pour qu'un point soit un optimum local (sous certaines hypothèses de régularité). Dans le contexte de l'optimisation de l'EE, les algorithmes développés par Zappone et al. \cite{Zappone2021} convergent vers des points satisfaisant ces conditions.\subsection{Travaux connexes}

L'optimisation de l'efficacité énergétique dans les réseaux sans fil a fait l'objet de nombreux travaux. On peut distinguer deux grandes familles d'approches : les approches centralisées et les approches décentralisées.

\subsubsection{Approches centralisées}

Les approches centralisées supposent l'existence d'une entité centrale disposant d'une connaissance parfaite de tous les canaux et de toutes les puissances. Cette entité résout globalement le problème d'optimisation. Les premiers travaux dans ce domaine utilisaient la programmation fractionnaire \cite{Dinkelbach1967}, qui permet de traiter les problèmes d'optimisation de ratios.

Cependant, comme le soulignent Zappone et al. \cite{Zappone2021}, la présence d'interférences rend le numérateur de l'EE non-concave, ce qui empêche l'application directe de la programmation fractionnaire classique. Des solutions basées sur des schémas de transmission orthogonaux \cite{Ng2012, Xu2014} ou sur l'annulation d'interférences \cite{Xu2013} ont été proposées, mais elles conduisent à un faible facteur de réutilisation des ressources.

D'autres approches utilisent des techniques heuristiques comme l'optimisation alternée \cite{Du2014, He2013, He2014}, mais sans garantie de convergence. Venturino et al. \cite{Venturino2015} ont proposé un cadre unifié pour l'optimisation de l'EE dans les réseaux OFDMA.

\subsubsection{Approches décentralisées}

Les approches décentralisées modélisent le réseau comme un jeu non-coopératif où chaque utilisateur maximise sa propre EE. Miao et al. \cite{Miao2011} ont étudié le problème de l'équilibre de Nash dans ce contexte, sans contraintes de débit. Scutari et al. \cite{Scutari2008} ont proposé des stratégies de précodage optimales basées sur la théorie des jeux.

Zappone et al. \cite{Zappone2013} ont étendu ces travaux aux réseaux avec relais, tandis que Stupia et al. \cite{Stupia2015} ont utilisé une approche par inéquations variationnelles. Bacci et al. \cite{Bacci2015} ont introduit des contraintes de QoS dans le jeu, rendant le problème plus réaliste mais aussi plus complexe.

L'article de Zappone et al. \cite{Zappone2021} constitue une référence majeure dans ce domaine. Il propose un cadre unifié pour l'optimisation centralisée et décentralisée de l'EE avec contraintes de QoS et une expression généralisée du SINR englobant plusieurs technologies 5G (Massive MIMO avec imperfections matérielles, réseaux avec relais, etc.). Les auteurs développent :
\begin{itemize}
    \item Des conditions de faisabilité sous forme fermée (Lemme 1),
    \item Des algorithmes centralisés convergeant vers des points KKT,
    \item Une modélisation du problème décentralisé comme un jeu généralisé,
    \item Une preuve d'existence et d'unicité du GNE (Propositions 3 et 4),
    \item Un algorithme distribué basé sur les meilleures réponses,
    \item Une extension au cas multi-blocs de ressources.
\end{itemize}

\subsubsection{Limites des approches existantes}

Malgré leur rigueur mathématique, les approches de Zappone et al. présentent plusieurs limitations :

\begin{enumerate}
    \item \textbf{Nécessité d'un modèle parfait} : les algorithmes reposent sur la connaissance exacte des paramètres $\alpha_k$, $\phi_k$, $\omega_{kj}$. En pratique, ces paramètres sont estimés avec des erreurs.
    \item \textbf{Stationnarité} : la convergence des algorithmes suppose que les canaux varient lentement. Dans un environnement dynamique (mobilité, fading rapide), l'algorithme n'a pas le temps de converger.
    \item \textbf{Absence d'apprentissage} : les performances ne s'améliorent pas avec l'expérience. Chaque nouvelle situation est traitée indépendamment des précédentes.
    \item \textbf{Sensibilité aux erreurs} : les approximations et erreurs d'estimation peuvent dégrader significativement les performances.
\end{enumerate}

\subsection{Contributions de cet article}

Face à ces limitations, nous proposons une approche alternative basée sur l'apprentissage par renforcement profond (Deep Reinforcement Learning). Nos principales contributions sont :

\begin{enumerate}
    \item Une formulation du problème de contrôle de puissance comme un Processus de Décision Markovien (MDP) avec un état composé d'observations locales, des actions discrètes (niveaux de puissance) et une récompense combinant EE, pénalité QoS et pénalité de puissance.
    \item L'utilisation d'un réseau de neurones profond (DQN) pour apprendre une politique de contrôle de puissance optimale, sans nécessiter de modèle mathématique explicite du canal.
    \item Une architecture de réseau adaptée au problème (8 entrées, couches cachées [256,512,256], sortie 16 actions).
    \item Des techniques de stabilisation avancées (experience replay, target network, Double DQN, Huber loss, gradient clipping, $\varepsilon$-greedy décroissant).
    \item Une évaluation expérimentale montrant un gain significatif de l'EE par rapport aux approches classiques.
\end{enumerate}

\subsection{Organisation de l'article}

La suite de cet article est organisée comme suit. La section 2 présente en détail l'approche décentralisée classique de Zappone et al., avec toutes les dérivations mathématiques. La section 3 décrit notre approche DQN, incluant la modélisation MDP, l'architecture du réseau et les techniques d'entraînement. La section 4 présente les résultats expérimentaux et les compare aux approches classiques. Enfin, la section 5 conclut l'article et discute des perspectives futures.

%==============================================================================
\section{L'approche décentralisée classique : jeu généralisé et équilibre de Nash}
%==============================================================================

[Le contenu détaillé de cette section sera développé dans la prochaine réponse]

% ============================================================
% Chapitres DQN + Comparaison — Rapport PFE
% Contrôle de Puissance Décentralisé via Deep Q-Network
% pour Réseaux 5G Massive MIMO
% ============================================================
% À inclure dans le document principal via \input{rapport_dqn_chapters}
% Prérequis packages : amsmath, amssymb, algorithm2e, graphicx,
%   booktabs, multirow, tabularx, subcaption, hyperref, xcolor

% ============================================================
%            CHAPITRE : APPROCHE PAR DEEP Q-NETWORK
% ============================================================
\section{Contrôle de Puissance par Deep Q-Network}
\label{chap:dqn}

\subsection{Introduction}

Le chapitre précédent a présenté l'approche analytique de Zappone \textit{et al.}~\cite{zappone2015} pour la maximisation décentralisée de l'efficacité énergétique (EE) dans les réseaux 5G Massive MIMO, reposant sur la dynamique de meilleure réponse (BRD) couplée à l'algorithme de Dinkelbach. Bien que cette méthode converge vers l'Équilibre de Nash Généralisé (GNE) sous certaines conditions, elle nécessite une connaissance explicite du modèle de canal et la résolution itérative d'un problème d'optimisation fractionnaire à chaque mise à jour.

Dans ce chapitre, nous proposons une approche alternative basée sur l'apprentissage par renforcement profond (\textit{Deep Reinforcement Learning}), et plus spécifiquement sur le \textbf{Deep Q-Network (DQN)}~\cite{mnih2015}. L'objectif est de permettre à chaque utilisateur $k$ d'apprendre de manière autonome une politique de contrôle de puissance qui maximise son EE individuelle, sans connaissance explicite du modèle mathématique du canal. L'approche est intrinsèquement décentralisée : chaque utilisateur dispose de son propre agent DQN qui prend des décisions indépendantes, tout en interagissant avec les autres agents via l'interférence partagée dans l'environnement sans fil.

% ============================================================
\subsection{Formulation en Processus Décisionnel de Markov}
\label{sec:mdp}

Pour appliquer l'apprentissage par renforcement, nous reformulons le problème de contrôle de puissance comme un \textbf{Processus Décisionnel de Markov} (MDP). Chaque utilisateur $k \in \{1, \dots, K\}$ est modélisé comme un agent autonome interagissant avec un environnement partagé. Le MDP de l'utilisateur $k$ est défini par le tuple $(\mathcal{S}_k, \mathcal{A}_k, \mathcal{P}, \mathcal{R}_k, \gamma)$, dont les composantes sont détaillées ci-après.

\subsubsection{Espace d'états}
\label{sec:state_space}

L'état observé par l'utilisateur $k$ à l'instant $t$ est un vecteur de dimension 6 :
\begin{equation}
    \mathbf{s}_k^{(t)} = \begin{bmatrix}
        s_1^{(t)} \\ s_2^{(t)} \\ s_3^{(t)} \\ s_4^{(t)} \\ s_5^{(t)} \\ s_6^{(t)}
    \end{bmatrix}
    \label{eq:state_vector}
\end{equation}
dont les composantes sont définies comme suit :

\begin{enumerate}
    \item \textbf{SINR normalisé en dB} :
    \begin{equation}
        s_1^{(t)} = \text{clip}\!\left(\frac{\text{SINR}_k^{\text{dB}}}{25}, -1, 1\right), \quad \text{SINR}_k^{\text{dB}} = 10\log_{10}(\gamma_k^{(t)})
        \label{eq:state_sinr}
    \end{equation}
    où la normalisation par 25~dB correspond à la plage pratique du SINR dans le système considéré.

    \item \textbf{Puissance normalisée (échelle logarithmique)} :
    \begin{equation}
        s_2^{(t)} = \frac{\log_{10}(p_k^{(t)}) - \log_{10}(p_{\min})}{\log_{10}(p_{\max}) - \log_{10}(p_{\min})}
        \label{eq:state_power}
    \end{equation}
    Cette représentation en échelle log permet une résolution uniforme sur toute la plage de puissance.

    \item \textbf{Interférence normalisée} :
    \begin{equation}
        s_3^{(t)} = \frac{\log_{10}\left(\sum_{j \neq k} p_j^{(t)}\right) - \log_{10}\left((K\!-\!1)\,p_{\min}\right)}{\log_{10}\left((K\!-\!1)\,p_{\max}\right) - \log_{10}\left((K\!-\!1)\,p_{\min}\right)}
        \label{eq:state_interf}
    \end{equation}

    \item \textbf{Gain de canal équivalent normalisé} :
    \begin{equation}
        s_4^{(t)} = \text{clip}\!\left(\frac{\log_{10}(\mu_k^{(t)})}{10}, 0, 1\right)
        \label{eq:state_mu}
    \end{equation}
    où $\mu_k = \frac{\alpha_k}{\sigma^2 + \sum_{j \neq k} \omega_{kj}\,p_j}$ est le gain de canal équivalent utilisé dans la meilleure réponse de Dinkelbach. Cette composante est cruciale car elle encode l'information nécessaire pour déterminer la puissance optimale.

    \item \textbf{Débit normalisé} :
    \begin{equation}
        s_5^{(t)} = \frac{R_k^{(t)}}{R_{\max}}, \quad R_{\max} = \log_2(1 + \bar{\gamma})
        \label{eq:state_rate}
    \end{equation}
    où $R_{\max}$ est le débit maximal théorique limité par les imperfections matérielles ($\bar{\gamma} = 1/\varepsilon^2$).

    \item \textbf{Marge QoS} :
    \begin{equation}
        s_6^{(t)} = \text{clip}\!\left(\frac{R_k^{(t)} - \theta_k}{\theta_k}, -1, 1\right)
        \label{eq:state_qos}
    \end{equation}
    où $\theta_k$ est le débit minimal requis par la contrainte de qualité de service. Une valeur positive indique que la contrainte est satisfaite.
\end{enumerate}

Toutes les composantes sont normalisées dans l'intervalle $[-1, 1]$ ou $[0, 1]$ afin de faciliter l'apprentissage du réseau de neurones. Le choix d'une représentation en échelle logarithmique pour la puissance et l'interférence reflète la nature multiplicative des signaux radio et améliore la sensibilité de l'agent aux variations dans les régimes de faible puissance, où se situe typiquement le point de fonctionnement optimal en termes d'EE.

\subsubsection{Espace d'actions}
\label{sec:action_space}

L'espace d'actions est \textbf{discret} et composé de $|\mathcal{A}| = 32$ niveaux de puissance répartis de manière logarithmique entre $p_{\min}$ et $p_{\max}$ :
\begin{equation}
    \mathcal{A} = \left\{ p_{\min} \cdot \left(\frac{p_{\max}}{p_{\min}}\right)^{\frac{i}{N_a - 1}} \;\middle|\; i = 0, 1, \dots, N_a - 1 \right\}
    \label{eq:action_space}
\end{equation}
avec $p_{\min} = -20$~dBm ($0{,}01$~mW) et $p_{\max} = 23$~dBm ($200$~mW), soit un rapport de 1:20\,000.

L'espacement logarithmique (géométrique) est motivé par le fait que l'efficacité énergétique optimale se situe typiquement dans la partie inférieure de la plage de puissance. Cette discrétisation offre ainsi une \textbf{résolution fine} dans la zone optimale (entre 0,01 et 5~mW) tout en couvrant l'ensemble de la plage autorisée. L'action $a_k^{(t)} \in \{0, 1, \dots, 31\}$ sélectionne directement le niveau de puissance d'émission $p_k^{(t)} = \mathcal{A}[a_k^{(t)}]$.

\subsubsection{Fonction de récompense}
\label{sec:reward}

La fonction de récompense est conçue pour guider l'agent vers la maximisation de l'efficacité énergétique tout en respectant les contraintes de QoS et en encourageant l'économie de puissance :

\begin{equation}
    r_k^{(t)} = \text{clip}\!\left(\alpha_{\text{ee}} \cdot \frac{\eta_k^{(t)}}{\eta_{\text{ref}}} - \beta_{\text{qos}} \cdot \mathbb{1}_{\text{QoS}} \cdot \left(0{,}3 + 0{,}7 \cdot \min\!\left(\frac{\theta_k - R_k^{(t)}}{\theta_k}, 1\right)\right) - \gamma_p \cdot \frac{p_k^{(t)}}{p_{\max}},\; -3,\; 2\right)
    \label{eq:reward}
\end{equation}

où :
\begin{itemize}
    \item $\eta_k^{(t)} = \frac{B \cdot \log_2(1 + \gamma_k^{(t)})}{p_c + p_k^{(t)}}$ est l'efficacité énergétique en bit/Joule ;
    \item $\eta_{\text{ref}} = 3 \times 10^8$ bit/J (= 300~Mbit/J) est la référence de normalisation, choisie pour que le terme EE soit dans l'intervalle $[0, 1]$ au voisinage de la solution optimale ;
    \item $\alpha_{\text{ee}} = 1{,}0$ est le coefficient de récompense EE ;
    \item $\beta_{\text{qos}} = 5{,}0$ est le coefficient de pénalité QoS, volontairement élevé pour fortement décourager les violations de la contrainte de débit minimal ;
    \item $\gamma_p = 0{,}02$ est le coefficient de pénalité de puissance, encourageant l'économie d'énergie ;
    \item $\mathbb{1}_{\text{QoS}} = \mathbb{1}_{R_k < \theta_k}$ est l'indicatrice de violation QoS.
\end{itemize}

La pénalité QoS est composée d'un terme fixe ($0{,}3$) activé dès qu'il y a violation, et d'un terme proportionnel au déficit relatif $(\theta_k - R_k)/\theta_k$, ce qui permet de différencier les violations légères des violations sévères. Le \textit{soft clipping} dans $[-3, 2]$ assure la stabilité numérique de l'entraînement en évitant les récompenses extrêmes.

\subsubsection{Dynamique de l'environnement}
\label{sec:env_dynamics}

L'environnement simule un réseau 5G Massive MIMO avec les caractéristiques suivantes :

\begin{itemize}
    \item \textbf{Canal dynamique} : à chaque pas de temps, le canal est mis à jour par un évanouissement rapide résiduel :
    \begin{equation}
        d_k^{\text{eff}}(t) = d_k \cdot (1 + \delta_k(t)), \quad \delta_k \sim \mathcal{N}\!\left(0, \frac{1}{\sqrt{M}}\right)
        \label{eq:channel_hardening}
    \end{equation}
    Ce modèle capture l'effet de \textit{channel hardening} typique du Massive MIMO avec $M = 50$ antennes, donnant des variations de $\pm 1$--$2$~dB.

    \item \textbf{Réinitialisation par épisode} : à chaque épisode, les $K = 5$ utilisateurs sont repositionnés aléatoirement dans la cellule (rayon $R_{\text{cell}} = 400$~m, distance minimale $d_{\min} = 35$~m), et un nouveau réalisation de l'affaiblissement grande échelle (perte de trajet + \textit{shadowing}) est tirée :
    \begin{equation}
        d_k = 10^{-(L_k + S_k)/10}, \quad L_k = 10\alpha\log_{10}(r_k), \quad S_k \sim \mathcal{N}(0, \sigma_s^2)
        \label{eq:large_scale_fading}
    \end{equation}
    avec $\alpha = 3{,}76$ et $\sigma_s = 8$~dB.

    \item \textbf{Durée d'épisode} : chaque épisode comprend $T = 200$ pas de temps, durant lesquels les agents interagissent simultanément.
\end{itemize}

% ============================================================
\subsection{Architecture du Deep Q-Network}
\label{sec:dqn_architecture}

\subsubsection{Principe du DQN}

Le Deep Q-Network~\cite{mnih2015} approxime la fonction de valeur action-état $Q^*(s, a)$ par un réseau de neurones profond $Q(s, a; \boldsymbol{\theta})$ paramétré par $\boldsymbol{\theta}$. La politique gloutonne (\textit{greedy}) correspondante est :
\begin{equation}
    \pi^*(s) = \arg\max_{a \in \mathcal{A}} Q(s, a; \boldsymbol{\theta})
    \label{eq:greedy_policy}
\end{equation}

L'objectif est de minimiser l'erreur temporelle différentielle (\textit{TD error}) entre la valeur Q courante et la valeur cible :
\begin{equation}
    \mathcal{L}(\boldsymbol{\theta}) = \mathbb{E}_{(s,a,r,s') \sim \mathcal{D}} \left[ \ell\!\left(r + \gamma \max_{a'} Q(s', a'; \boldsymbol{\theta}^{-}) - Q(s, a; \boldsymbol{\theta})\right) \right]
    \label{eq:dqn_loss}
\end{equation}
où $\mathcal{D}$ est le \textit{replay buffer}, $\boldsymbol{\theta}^{-}$ les paramètres du réseau cible, $\gamma = 0{,}95$ le facteur d'actualisation, et $\ell(\cdot)$ la fonction de perte.

\subsubsection{Architecture du réseau de neurones}

Le réseau de neurones utilisé pour chaque agent DQN est un perceptron multicouche (\textit{Multi-Layer Perceptron}, MLP) composé de trois couches cachées entièrement connectées (\textit{fully connected}). L'architecture est la suivante :

\begin{equation}
    \underbrace{\mathbb{R}^6}_{\text{Entrée (état)}} \xrightarrow{\mathbf{W}_1, \mathbf{b}_1} \underbrace{\mathbb{R}^{128}}_{\text{Couche 1}} \xrightarrow{\text{ReLU}} \underbrace{\mathbb{R}^{256}}_{\text{Couche 2}} \xrightarrow{\text{ReLU}} \underbrace{\mathbb{R}^{128}}_{\text{Couche 3}} \xrightarrow{\text{ReLU}} \underbrace{\mathbb{R}^{32}}_{\text{Sortie (Q-values)}}
    \label{eq:architecture}
\end{equation}

\begin{itemize}
    \item \textbf{Couche d'entrée} : 6 neurones correspondant aux 6 composantes de l'état $\mathbf{s}_k$.
    \item \textbf{Couche cachée 1} : 128 neurones avec activation ReLU.
    \item \textbf{Couche cachée 2} : 256 neurones avec activation ReLU.
    \item \textbf{Couche cachée 3} : 128 neurones avec activation ReLU.
    \item \textbf{Couche de sortie} : 32 neurones (un par action), sans activation — les sorties représentent directement les Q-valeurs $Q(s, a_i; \boldsymbol{\theta})$ pour $i = 0, \dots, 31$.
\end{itemize}

L'architecture en forme de \textit{bottleneck} (128 $\to$ 256 $\to$ 128) permet d'abord une expansion de la représentation dans l'espace latent, puis une compression vers la dimension de sortie, favorisant l'extraction de caractéristiques pertinentes.

Les poids sont initialisés selon la méthode de \textbf{Xavier uniforme}~\cite{glorot2010} :
\begin{equation}
    W_{ij} \sim \mathcal{U}\left(-\sqrt{\frac{6}{n_{\text{in}} + n_{\text{out}}}},\; +\sqrt{\frac{6}{n_{\text{in}} + n_{\text{out}}}}\right)
    \label{eq:xavier_init}
\end{equation}
et les biais sont initialisés à zéro. Le nombre total de paramètres entraînables par agent est de :
\begin{equation}
    N_{\text{params}} = (6 \times 128 + 128) + (128 \times 256 + 256) + (256 \times 128 + 128) + (128 \times 32 + 32) = 70\,944
    \label{eq:n_params}
\end{equation}

\begin{figure*}[htbp]
    \centering
    \includegraphics[width=0.85\textwidth]{figures/dqn_architecture.png}
    \caption{Architecture du réseau DQN pour le contrôle de puissance. L'entrée est le vecteur d'état à 6 dimensions et la sortie fournit les Q-valeurs pour les 32 niveaux de puissance discrets.}
    \label{fig:dqn_architecture}
\end{figure*}

% ============================================================
\subsection{Mécanismes d'apprentissage}
\label{sec:learning_mechanisms}

\subsubsection{Double DQN}
\label{sec:double_dqn}

Afin de réduire la surestimation des Q-valeurs inhérente au DQN standard~\cite{vanHasselt2016}, nous implémentons la variante \textbf{Double DQN}. Le principe consiste à découpler la \textbf{sélection} de l'action (effectuée par le réseau de politique $\boldsymbol{\theta}$) de son \textbf{évaluation} (effectuée par le réseau cible $\boldsymbol{\theta}^{-}$) :

\begin{equation}
    y^{\text{DDQN}} = r + \gamma \cdot Q\!\left(s',\; \underbrace{\arg\max_{a'} Q(s', a'; \boldsymbol{\theta})}_{\text{sélection par réseau politique}};\; \boldsymbol{\theta}^{-}\right)
    \label{eq:double_dqn_target}
\end{equation}

Cette technique réduit significativement le biais d'estimation positif et améliore la stabilité de l'apprentissage, ce qui est particulièrement important dans notre contexte multi-agent où les environnements sont non-stationnaires du point de vue de chaque agent.

\subsubsection{Replay d'expérience}
\label{sec:experience_replay}

Chaque agent maintient un \textbf{tampon de rejeu} (\textit{experience replay buffer}) $\mathcal{D}_k$ de capacité $|\mathcal{D}| = 100\,000$ transitions. Chaque transition stockée est un tuple :
\begin{equation}
    e^{(t)} = \left(s_k^{(t)},\; a_k^{(t)},\; r_k^{(t)},\; s_k^{(t+1)},\; d^{(t)}\right)
    \label{eq:transition}
\end{equation}
où $d^{(t)} \in \{0, 1\}$ est l'indicateur de fin d'épisode. Le tampon fonctionne selon le principe FIFO (\textit{First-In, First-Out}) : lorsque la capacité maximale est atteinte, les transitions les plus anciennes sont remplacées.

À chaque pas d'entraînement, un \textbf{mini-batch} de $B = 256$ transitions est échantillonné uniformément de manière aléatoire :
\begin{equation}
    \{(s_i, a_i, r_i, s'_i, d_i)\}_{i=1}^{B} \sim \text{Uniforme}(\mathcal{D}_k)
    \label{eq:sampling}
\end{equation}

Le replay d'expérience confère deux avantages fondamentaux :
\begin{itemize}
    \item \textbf{Décorrélation temporelle} : les transitions échantillonnées aléatoirement brisent les corrélations séquentielles, satisfaisant l'hypothèse i.i.d. de la descente de gradient stochastique.
    \item \textbf{Efficacité en données} : chaque transition peut être réutilisée dans plusieurs mises à jour, améliorant l'exploitation des données collectées.
\end{itemize}

L'entraînement ne commence qu'une fois le tampon contenant au moins $|\mathcal{D}|_{\min} = 2\,000$ transitions, garantissant une diversité suffisante dans les premiers mini-batchs.

\subsubsection{Réseau cible et mise à jour douce}
\label{sec:target_network}

Pour stabiliser l'apprentissage, nous utilisons un \textbf{réseau cible} (\textit{target network}) $Q(s, a; \boldsymbol{\theta}^{-})$ distinct du réseau de politique $Q(s, a; \boldsymbol{\theta})$. Les paramètres du réseau cible sont mis à jour par \textbf{mise à jour douce} (\textit{soft update} ou \textit{Polyak averaging}) :

\begin{equation}
    \boldsymbol{\theta}^{-} \leftarrow \tau \cdot \boldsymbol{\theta} + (1 - \tau) \cdot \boldsymbol{\theta}^{-}, \quad \tau = 0{,}005
    \label{eq:soft_update}
\end{equation}

Cette mise à jour est effectuée après chaque pas d'optimisation. Le coefficient $\tau = 0{,}005$ assure une transition lente et régulière du réseau cible vers le réseau de politique, évitant les oscillations et divergences liées à une cible en mouvement rapide. Par rapport à la mise à jour dure (\textit{hard update}) périodique, la mise à jour douce offre une stabilité accrue car le réseau cible évolue de manière continue plutôt que par sauts discontinus.

\subsubsection{Fonction de perte de Huber}
\label{sec:huber_loss}

La fonction de perte utilisée pour l'optimisation est la \textbf{perte de Huber} (aussi appelée \textit{Smooth L1 Loss}) :
\begin{equation}
    \ell_{\delta}(x) = \begin{cases}
        \frac{1}{2}x^2 & \text{si } |x| \leq \delta \\
        \delta\left(|x| - \frac{1}{2}\delta\right) & \text{sinon}
    \end{cases}
    \label{eq:huber_loss}
\end{equation}
avec $\delta = 1{,}0$. La perte de Huber combine les avantages de la MSE (sensibilité aux petites erreurs) et de la MAE (robustesse aux valeurs aberrantes), ce qui est crucial dans notre environnement multi-agent où les récompenses peuvent être bruitées par les actions concurrentes des autres agents.

\subsubsection{Politique d'exploration $\varepsilon$-gloutonne}
\label{sec:epsilon_greedy}

L'exploration du space d'actions est réalisée par une politique $\varepsilon$-gloutonne (\textit{$\varepsilon$-greedy}) :
\begin{equation}
    a_k^{(t)} = \begin{cases}
        \text{Uniforme}(\mathcal{A}) & \text{avec probabilité } \varepsilon_t \\
        \arg\max_{a} Q(s_k^{(t)}, a; \boldsymbol{\theta}_k) & \text{avec probabilité } 1 - \varepsilon_t
    \end{cases}
    \label{eq:epsilon_greedy}
\end{equation}

Le taux d'exploration $\varepsilon_t$ décroît de manière exponentielle au fil des épisodes :
\begin{equation}
    \varepsilon_t = \max\!\left(\varepsilon_{\text{fin}},\; \varepsilon_0 \cdot \lambda^t\right), \quad \varepsilon_0 = 1{,}0, \quad \varepsilon_{\text{fin}} = 0{,}02, \quad \lambda = 0{,}9995
    \label{eq:epsilon_decay}
\end{equation}

Ce schéma de décroissance assure :
\begin{itemize}
    \item Une \textbf{phase d'exploration intensive} initiale (épisodes 1--1\,000, $\varepsilon > 0{,}6$) permettant une couverture large de l'espace état-action.
    \item Une \textbf{transition progressive} (épisodes 1\,000--7\,000) où l'agent bascule graduellement vers l'exploitation.
    \item Une \textbf{phase d'exploitation} finale (épisodes 7\,000--10\,000, $\varepsilon \approx 0{,}02$) où l'agent utilise quasi-exclusivement la politique apprise, tout en conservant une exploration résiduelle de 2\% pour maintenir une certaine adaptabilité.
\end{itemize}

L'épisode où $\varepsilon$ passe sous le seuil $0{,}1$ (exploitation dominante) se situe aux alentours de l'épisode $4\,600$, ce qui correspond au régime où les Q-valeurs sont suffisamment stables pour guider les décisions.

\begin{figure}[htbp]
    \centering
    \includegraphics[width=1\textwidth]{figures/epsilon_decay.png}
    \caption{Décroissance exponentielle du taux d'exploration $\varepsilon$ au cours des $10\,000$ épisodes d'entraînement. La ligne pointillée indique le seuil $\varepsilon = 0{,}1$.}
    \label{fig:epsilon_decay}
\end{figure}

\subsubsection{Gradient clipping}
\label{sec:grad_clip}

Pour prévenir l'instabilité numérique causée par des gradients de grande magnitude, une \textbf{troncature du gradient} (\textit{gradient clipping}) est appliquée avant chaque mise à jour des poids :
\begin{equation}
    \nabla_{\boldsymbol{\theta}} \mathcal{L} \leftarrow \frac{\nabla_{\boldsymbol{\theta}} \mathcal{L}}{\max\!\left(1, \frac{\|\nabla_{\boldsymbol{\theta}} \mathcal{L}\|_2}{G_{\max}}\right)}, \quad G_{\max} = 1{,}0
    \label{eq:grad_clip}
\end{equation}

Cette technique est d'autant plus importante dans le contexte multi-agent, où la non-stationnarité induite par l'apprentissage simultané de tous les agents peut engendrer des gradients instables.

% ============================================================
\subsection{Entraînement multi-agent décentralisé}
\label{sec:multi_agent}

\subsubsection{Independent Q-Learning (IQL)}

L'approche adoptée est le \textbf{Independent Q-Learning} (IQL)~\cite{tan1993}, où chaque utilisateur $k$ possède son propre agent DQN qui apprend indépendamment. Les agents interagissent uniquement à travers l'environnement partagé : le SINR de l'utilisateur $k$ dépend des puissances d'émission de tous les utilisateurs via l'interférence :
\begin{equation}
    \gamma_k = \frac{\alpha_k \, p_k}{\sigma^2 + \phi_k \, p_k + \sum_{j \neq k} \omega_{kj} \, p_j}
    \label{eq:sinr_coupling}
\end{equation}

Du point de vue de l'utilisateur $k$, les autres utilisateurs font partie de l'environnement, qui devient \textbf{non-stationnaire} puisque les politiques des autres agents évoluent au cours de l'entraînement. Bien que cela viole formellement les hypothèses de stationnarité du MDP, l'IQL fonctionne remarquablement bien en pratique pour les jeux coopératifs et compétitifs de taille modérée~\cite{lanctot2017}.

\subsubsection{Protocole d'entraînement}

Le protocole d'entraînement est résumé dans l'Algorithme~\ref{alg:training}.

\begin{algorithm}[htbp]
\caption{Entraînement multi-agent DQN décentralisé}
\label{alg:training}
\SetKwInOut{Input}{Entrées}
\SetKwInOut{Output}{Sorties}
\Input{$K = 5$ agents, $N_{\text{ep}} = 10\,000$ épisodes, $T = 200$ pas/épisode}
\Output{Politiques optimales $\{\boldsymbol{\theta}_k^*\}_{k=1}^{K}$}
\BlankLine
Initialiser $K$ réseaux de politique $\{Q(\cdot; \boldsymbol{\theta}_k)\}$ et cibles $\{Q(\cdot; \boldsymbol{\theta}_k^{-})\}$\;
Initialiser $K$ tampons de rejeu $\{\mathcal{D}_k\}$ de capacité $100\,000$\;
$\varepsilon \leftarrow 1{,}0$\;
\BlankLine
\For{épisode $= 1$ \KwTo $N_{\text{ep}}$}{
    Réinitialiser l'environnement (nouvelles positions + canal)\;
    Observer les états initiaux $\{s_k^{(0)}\}_{k=1}^{K}$\;
    \BlankLine
    \For{$t = 0$ \KwTo $T-1$}{
        \For{$k = 1$ \KwTo $K$ \textrm{(en parallèle)}}{
            Sélectionner $a_k^{(t)}$ par politique $\varepsilon$-gloutonne (Éq.~\ref{eq:epsilon_greedy})\;
        }
        Exécuter les actions $\{a_k^{(t)}\}$ dans l'environnement\;
        Observer récompenses $\{r_k^{(t)}\}$ et nouveaux états $\{s_k^{(t+1)}\}$\;
        \BlankLine
        \For{$k = 1$ \KwTo $K$}{
            Stocker $(s_k^{(t)}, a_k^{(t)}, r_k^{(t)}, s_k^{(t+1)}, d^{(t)})$ dans $\mathcal{D}_k$\;
            \If{$|\mathcal{D}_k| \geq 2\,000$ \textbf{et} $t \bmod 16 = 0$}{
                Échantillonner un mini-batch de 256 transitions de $\mathcal{D}_k$\;
                Calculer les cibles Double DQN (Éq.~\ref{eq:double_dqn_target})\;
                Mettre à jour $\boldsymbol{\theta}_k$ par descente de gradient (perte de Huber)\;
                Mise à jour douce : $\boldsymbol{\theta}_k^{-} \leftarrow 0{,}005 \cdot \boldsymbol{\theta}_k + 0{,}995 \cdot \boldsymbol{\theta}_k^{-}$\;
            }
        }
    }
    $\varepsilon \leftarrow \max(0{,}02,\; \varepsilon \times 0{,}9995)$\;
    \BlankLine
    \If{épisode $\bmod$ 500 $= 0$}{
        Sauvegarder les checkpoints $\{\boldsymbol{\theta}_k\}$\;
    }
}
\end{algorithm}

Un détail d'implémentation important est la \textbf{fréquence d'entraînement réduite} : l'optimisation du réseau n'est effectuée que tous les 16 pas de temps ($t \bmod 16 = 0$), ce qui réduit le surcoût de transfert CPU$\leftrightarrow$GPU et évite l'overfitting sur des transitions fortement corrélées au sein d'un même épisode.

\subsubsection{Optimisation}

L'optimiseur utilisé est \textbf{Adam}~\cite{kingma2015} avec un taux d'apprentissage $\eta = 10^{-4}$ :
\begin{equation}
    \boldsymbol{\theta}_{t+1} = \boldsymbol{\theta}_t - \eta \cdot \frac{\hat{m}_t}{\sqrt{\hat{v}_t} + \epsilon_{\text{Adam}}}
    \label{eq:adam}
\end{equation}
où $\hat{m}_t$ et $\hat{v}_t$ sont les estimateurs corrigés du premier et second moment des gradients. Le taux d'apprentissage relativement faible ($10^{-4}$) est choisi pour assurer une convergence stable dans l'environnement multi-agent non-stationnaire.

% ============================================================
\subsection{Sauvegarde et évaluation des modèles}
\label{sec:checkpoints}

\subsubsection{Système de checkpoints}

Des points de contrôle (\textit{checkpoints}) sont sauvegardés tous les 500 épisodes, contenant :
\begin{itemize}
    \item Les paramètres du réseau de politique $\boldsymbol{\theta}_k$ et du réseau cible $\boldsymbol{\theta}_k^{-}$
    \item L'état de l'optimiseur Adam (moments)
    \item La valeur courante de $\varepsilon$
    \item Le compteur d'étapes d'entraînement
    \item L'historique des pertes
\end{itemize}

Cela permet de reprendre l'entraînement en cas d'interruption et d'évaluer les performances intermédiaires de la politique à différents stades de l'apprentissage.

\subsubsection{Protocole d'évaluation}

L'évaluation des modèles entraînés est réalisée en mode \textbf{glouton} ($\varepsilon = 0$) sur $N_{\text{eval}} = 50$ réalisations de canal indépendantes. Pour chaque réalisation, les agents sont exécutés pendant $T = 200$ pas afin de stabiliser les puissances d'émission avant de mesurer l'EE résultante.

\begin{figure*}[htbp]
    \centering
    \includegraphics[width=0.95\textwidth]{figures/training_curves.png}
    \caption{Courbes d'entraînement des agents DQN décentralisés sur $10\,000$ épisodes : (a) efficacité énergétique moyenne, (b) reward cumulatif, (c) perte d'entraînement, (d) SINR moyen, (e) puissance d'émission moyenne, (f) taux d'exploration $\varepsilon$.}
    \label{fig:training_curves}
\end{figure*}

% ============================================================
\subsection{Récapitulatif des hyperparamètres}
\label{sec:hyperparams_summary}

Le Tableau~\ref{tab:hyperparams} regroupe l'ensemble des hyperparamètres utilisés pour l'entraînement.

\begin{table}[H]
    \centering
    \caption{Hyperparamètres de l'entraînement DQN}
    \label{tab:hyperparams}
    \begin{tabular}{@{}llr@{}}
        \toprule
        \textbf{Catégorie} & \textbf{Paramètre} & \textbf{Valeur} \\
        \midrule
        \multirow{4}{*}{Architecture} 
            & Couches cachées & $[128, 256, 128]$ \\
            & Activation & ReLU \\
            & Initialisation & Xavier uniforme \\
            & Params/agent & 70\,944 \\
        \midrule
        \multirow{5}{*}{Optimisation}
            & Optimiseur & Adam \\
            & Taux d'apprentissage $\eta$ & $10^{-4}$ \\
            & Taille du batch $B$ & 256 \\
            & Fonction de perte & Huber ($\delta = 1{,}0$) \\
            & Gradient clipping & $G_{\max} = 1{,}0$ \\
        \midrule
        \multirow{3}{*}{Exploration}
            & $\varepsilon_0$ & 1,0 \\
            & $\varepsilon_{\text{fin}}$ & 0,02 \\
            & Décroissance $\lambda$ & 0,9995/épisode \\
        \midrule
        \multirow{4}{*}{Replay \& cible}
            & Taille du buffer $|\mathcal{D}|$ & 100\,000 \\
            & Buffer minimal $|\mathcal{D}|_{\min}$ & 2\,000 \\
            & Soft update $\tau$ & 0,005 \\
            & Fréquence entraîn. & 1/16 pas \\
        \midrule
        \multirow{3}{*}{Entraînement}
            & Facteur de discount $\gamma$ & 0,95 \\
            & Épisodes $N_{\text{ep}}$ & 10\,000 \\
            & Pas/épisode $T$ & 200 \\
        \midrule
        \multirow{3}{*}{Récompense}
            & $\alpha_{\text{ee}}$ & 1,0 \\
            & $\beta_{\text{qos}}$ & 5,0 \\
            & $\gamma_p$ & 0,02 \\
        \bottomrule
    \end{tabular}
\end{table}


% ============================================================
% ============================================================
%     CHAPITRE : RÉSULTATS ET COMPARAISONS
% ============================================================
% ============================================================

\section{Résultats Expérimentaux et Analyse Comparative}

\label{chap:results}

\section{Introduction}

Ce chapitre présente une analyse comparative détaillée entre l'approche DQN décentralisée développée dans le Chapitre~\ref{chap:dqn} et les méthodes de référence. Nous évaluons les performances sur $N = 50$ réalisations de canal indépendantes, garantissant une comparaison statistiquement significative. Toutes les méthodes sont testées sur les \textbf{mêmes} réalisations de canal (échantillons appariés), assurant l'équité de la comparaison.

\section{Méthodes comparées}
\label{sec:methods}

Cinq méthodes de contrôle de puissance sont évaluées :

\begin{enumerate}
    \item \textbf{Article (BRD + Dinkelbach)} : l'algorithme de référence de Zappone \textit{et al.}~\cite{zappone2015}, basé sur la dynamique de meilleure réponse avec résolution interne par l'algorithme de Dinkelbach (Chapitre~\ref{chap:article_algo}).

    \item \textbf{DQN (entraîné)} : notre approche par apprentissage par renforcement profond, avec $K = 5$ agents DQN indépendants entraînés sur $10\,000$ épisodes (Chapitre~\ref{chap:dqn}).

    \item \textbf{Optimal Uniform} : recherche exhaustive de la puissance uniforme $p^*$ optimale parmi 100 niveaux géométriquement espacés, où tous les utilisateurs émettent à la même puissance. Cette méthode sert de baseline analytique simple.

    \item \textbf{Max Power} : tous les utilisateurs émettent à la puissance maximale $p_{\max} = 23$~dBm ($200$~mW). Ce baseline naïf illustre la sous-optimalité de l'approche brute en termes d'EE.

    \item \textbf{Random Power} : chaque utilisateur choisit une puissance aléatoire uniforme dans $[p_{\min}, p_{\max}]$. Ce baseline représente l'absence totale de stratégie.
\end{enumerate}

\section{Résultats en efficacité énergétique}
\label{sec:ee_results}

\subsection{Performances globales}

Le Tableau~\ref{tab:ee_comparison} présente les statistiques descriptives de l'efficacité énergétique pour chaque méthode sur les 50 essais.

\begin{table*}[htbp]
    \centering
    \caption{Comparaison de l'efficacité énergétique (Mbit/J) — 50 essais}
    \label{tab:ee_comparison}
    \begin{tabular}{@{}lrrrrr@{}}
        \toprule
        \textbf{Méthode} & \textbf{Moyenne} & \textbf{Écart-type} & \textbf{Min} & \textbf{Max} & \textbf{Médiane} \\
        \midrule
        \textbf{DQN (entraîné)}         & \textbf{255,16} & 47,31 & 141,27 & 313,18 & --- \\
        Article (BRD+Dinkelbach)         & 207,01          & 30,67 & 137,99 & 256,16 & --- \\
        Optimal Uniform                  & 213,04          & 52,76 & ---    & ---    & --- \\
        Random Power                     & 23,06           & 15,59 & ---    & ---    & --- \\
        Max Power                        & 10,22           & 2,55  & ---    & ---    & --- \\
        \bottomrule
    \end{tabular}
\end{table*}

L'approche DQN atteint une EE moyenne de $\boldsymbol{255{,}16}$~\textbf{Mbit/J}, surpassant significativement l'algorithme de l'article ($207{,}01$~Mbit/J) avec un \textbf{gain relatif de $+23{,}3$~\%}. Les méthodes naïves (Max Power, Random) sont nettement inférieures, confirmant que le contrôle intelligent de puissance est essentiel pour l'efficacité énergétique.

\begin{figure}[htbp]
    \centering
    \includegraphics[width=0.9\textwidth]{figures/comparison_bar.png}
    \caption{Efficacité énergétique moyenne (Mbit/J) par méthode avec barres d'erreur ($\pm 1$ écart-type). Le DQN surpasse toutes les méthodes de référence.}
    \label{fig:ee_barplot}
\end{figure}

\subsection{Distribution de l'efficacité énergétique}

La Figure~\ref{fig:ee_distribution} présente la distribution de l'EE pour l'approche DQN et l'algorithme de l'article sur les 50 essais.

\begin{figure*}[htbp]
    \centering
    \subfloat[Histogramme comparatif\label{fig:ee_histogram}]{%
        \includegraphics[width=0.48\textwidth]{figures/ee_histogram_dqn_vs_article.png}}
    \hfill
    \subfloat[Fonction de répartition cumulative (CDF)\label{fig:ee_cdf}]{%
        \includegraphics[width=0.48\textwidth]{figures/comparison_cdf.png}}
    \caption{Distribution de l'efficacité énergétique : (a) histogrammes superposés montrant le décalage favorable du DQN ; (b) CDF confirmant la dominance stochastique du DQN sur l'algorithme de l'article.}
    \label{fig:ee_distribution}
\end{figure*}

La CDF (Figure~\ref{fig:ee_cdf}) montre que pour tout seuil d'EE, la probabilité que le DQN dépasse ce seuil est plus élevée que celle de l'algorithme de l'article, ce qui constitue une \textbf{dominance stochastique du premier ordre}. La médiane du DQN se situe aux alentours de 255~Mbit/J contre 218~Mbit/J pour l'article.

\subsection{Comparaison essai par essai}

La Figure~\ref{fig:ee_scatter_trial} présente un graphique de dispersion comparant l'EE du DQN à celle de l'article pour chaque essai individuellement.

\begin{figure}[htbp]
    \centering
    \includegraphics[width=0.65\textwidth]{figures/ee_scatter_dqn_vs_article.png}
    \caption{Comparaison essai par essai de l'EE : DQN (axe Y) vs Article (axe X). Les points au-dessus de la diagonale ($y = x$) indiquent une supériorité du DQN. La grande majorité des essais favorise le DQN.}
    \label{fig:ee_scatter_trial}
\end{figure}

% ============================================================
\section{Analyse de la consommation de puissance}
\label{sec:power_results}

\subsection{Puissance moyenne par méthode}

Le Tableau~\ref{tab:power_comparison} détaille la consommation de puissance moyenne par méthode.

\begin{table}[htbp]
    \centering
    \caption{Consommation de puissance moyenne (mW) par méthode}
    \label{tab:power_comparison}
    \begin{tabular}{@{}lr@{}}
        \toprule
        \textbf{Méthode} & \textbf{Puissance moyenne (mW)} \\
        \midrule
        \textbf{DQN (entraîné)}         & \textbf{15,27} \\
        Article (BRD+Dinkelbach)         & 28,73 \\
        Optimal Uniform                  & 0,02 \\
        Max Power                        & 199,53 \\
        Random Power                     & 100,57 \\
        \bottomrule
    \end{tabular}
\end{table}

L'approche DQN consomme en moyenne $\boldsymbol{15{,}27}$~\textbf{mW}, soit une \textbf{réduction de $46{,}8$~\%} par rapport à l'algorithme de l'article ($28{,}73$~mW). Cette économie de puissance significative est obtenue tout en améliorant l'EE, ce qui démontre que le DQN apprend à opérer dans la zone optimale du compromis débit/consommation.

\begin{figure}[htbp]
    \centering
    \includegraphics[width=0.9\textwidth]{figures/comparison_boxplot.png}
    \caption{Distribution de la puissance d'émission (mW) par méthode sous forme de boîtes à moustaches (échelle logarithmique). Le DQN opère à des puissances significativement inférieures à l'algorithme de l'article.}
    \label{fig:power_boxplot}
\end{figure}

\subsection{Compromis puissance--EE}

La Figure~\ref{fig:power_vs_ee} illustre le compromis fondamental entre la puissance d'émission et l'efficacité énergétique pour toutes les méthodes.

\begin{figure}[htbp]
    \centering
    \includegraphics[width=0.9\textwidth]{figures/scatter_power_vs_ee.png}
    \caption{Diagramme de dispersion puissance (mW, échelle log) vs EE (Mbit/J) pour les 50 essais. Chaque point représente un essai; la couleur indique la méthode. Le DQN (vert) se concentre dans la zone optimale : faible puissance, haute EE.}
    \label{fig:power_vs_ee}
\end{figure}

Ce graphique met en évidence le positionnement optimal du DQN dans le coin supérieur gauche (haute EE, faible puissance), confirmant que l'agent a appris la structure du compromis débit-puissance décrit dans l'article de Zappone \textit{et al.} : l'EE est une fonction quasi-concave de la puissance d'émission, avec un maximum bien défini à des puissances relativement basses.

% ============================================================
\section{Analyse de l'algorithme de l'article}
\label{sec:article_results}

L'algorithme de référence (BRD + Dinkelbach) présente les caractéristiques suivantes sur les 50 essais :

\subsection{Convergence}

\begin{itemize}
    \item \textbf{Fraction convergée} : $100\%$ des essais convergent (tolérance relative $10^{-4}$).
    \item \textbf{Nombre moyen d'itérations} : $5{,}28$ itérations de BRD par essai.
    \item \textbf{Nombre maximum d'itérations} : $12$ (convergence toujours rapide).
\end{itemize}

La convergence rapide et systématique confirme les résultats théoriques de l'article : sous les conditions de faisabilité, la dynamique de meilleure réponse converge vers l'unique GNE.

\subsection{Efficacité énergétique détaillée}

Les statistiques descriptives de l'EE obtenue par l'algorithme de l'article sont présentées dans le Tableau~\ref{tab:article_ee_stats}.

\begin{table}[htbp]
    \centering
    \caption{Statistiques détaillées de l'EE — Algorithme de l'article (50 essais)}
    \label{tab:article_ee_stats}
    \begin{tabular}{@{}lr@{}}
        \toprule
        \textbf{Statistique} & \textbf{Valeur (Mbit/J)} \\
        \midrule
        Moyenne              & 218,07 \\
        Écart-type           & 30,10 \\
        Minimum              & 137,00 \\
        Maximum              & 268,39 \\
        Médiane              & 218,08 \\
        Percentile 5\%       & 179,47 \\
        Percentile 95\%      & 264,72 \\
        \bottomrule
    \end{tabular}
\end{table}

\subsection{Débit spectral et QoS}

\begin{itemize}
    \item \textbf{Débit spectral moyen} : $\bar{R} = 2{,}951$~bits/s/Hz (écart-type~$= 0{,}347$).
    \item \textbf{Débit maximal théorique} : $R_{\max} = \log_2(1 + 125) = 6{,}977$~bits/s/Hz.
    \item \textbf{Contrainte QoS} : $\theta_k = 20\% \times R_{\max} = 1{,}395$~bits/s/Hz.
    \item \textbf{Taux de satisfaction QoS} : $\boldsymbol{74\%}$ des utilisateurs satisfont la contrainte de débit minimal.
\end{itemize}

Le taux de satisfaction QoS de 74\% indique que même l'algorithme optimal ne parvient pas toujours à satisfaire les contraintes de débit dans des conditions de canal défavorables, ce qui souligne la difficulté inhérente du problème de faisabilité dans les systèmes à interférence.

\begin{figure*}[htbp]
    \centering
    \subfloat[Distribution de l'EE\label{fig:article_ee_dist}]{%
        \includegraphics[width=0.48\textwidth]{figures/article_ee_distribution.png}}
    \hfill
    \subfloat[Distribution de la puissance\label{fig:article_power_dist}]{%
        \includegraphics[width=0.48\textwidth]{figures/article_power_distribution.png}}
    \caption{Distributions des métriques de l'algorithme de l'article sur 50 essais : (a) EE avec annotations des percentiles P5 et P95 ; (b) puissance émise montrant une forte variance entre les réalisations de canal.}
    \label{fig:article_distributions}
\end{figure*}

% ============================================================
\section{Analyse de l'entraînement DQN}
\label{sec:training_analysis}

\subsection{Convergence de l'efficacité énergétique}

La Figure~\ref{fig:training_ee} illustre l'évolution de l'EE moyenne au cours de l'entraînement. On distingue trois phases :

\begin{enumerate}
    \item \textbf{Phase d'exploration} (épisodes 1--2\,000) : l'EE est faible et volatile, les agents explorent l'espace d'actions de manière quasi-aléatoire.
    \item \textbf{Phase d'apprentissage rapide} (épisodes 2\,000--6\,000) : l'EE croît rapidement à mesure que les Q-valeurs convergent et que l'exploitation commence à dominer.
    \item \textbf{Phase de stabilisation} (épisodes 6\,000--10\,000) : l'EE se stabilise autour de 250--300~Mbit/J, les agents ayant convergé vers des politiques quasi-optimales.
\end{enumerate}

\begin{figure}[htbp]
    \centering
    \includegraphics[width=0.8\textwidth]{figures/training_ee_convergence.png}
    \caption{Convergence de l'efficacité énergétique moyenne durant l'entraînement. La courbe lissée (fenêtre de 100 épisodes) montre une convergence progressive vers $\sim 250$~Mbit/J. La ligne pointillée rouge indique le niveau atteint par l'algorithme de l'article ($\sim 218$~Mbit/J).}
    \label{fig:training_ee}
\end{figure}

\subsection{Évolution de la perte d'entraînement}

La Figure~\ref{fig:training_loss} montre l'évolution de la perte de Huber au cours de l'entraînement. Après un pic initial correspondant aux premières mises à jour des Q-valeurs (lorsque le réseau cible est encore initialisé aléatoirement), la perte décroît progressivement, indiquant une réduction de l'erreur de prédiction temporelle. La stabilisation de la perte aux alentours de l'épisode 5\,000 confirme que les Q-valeurs convergent.

\begin{figure}[htbp]
    \centering
    \includegraphics[width=0.75\textwidth]{figures/training_loss.png}
    \caption{Perte d'entraînement (Huber loss) lissée sur une fenêtre de 100 épisodes. La décroissance progressive confirme la convergence des Q-valeurs.}
    \label{fig:training_loss}
\end{figure}

\subsection{Évolution de la puissance d'émission}

L'analyse de la puissance d'émission moyenne au cours de l'entraînement (Figure~\ref{fig:training_power}) révèle un comportement caractéristique : les agents apprennent progressivement à réduire leur puissance d'émission, passant d'une puissance élevée (exploration aléatoire) à des puissances optimales basses, cohérentes avec la maximisation de l'EE.

\begin{figure}[htbp]
    \centering
    \includegraphics[width=0.75\textwidth]{figures/training_power_evolution.png}
    \caption{Évolution de la puissance d'émission moyenne (dBm) au cours de l'entraînement. Les agents apprennent à réduire progressivement leur puissance vers le point de fonctionnement optimal.}
    \label{fig:training_power}
\end{figure}

\subsection{Évolution du SINR}

La Figure~\ref{fig:training_sinr} présente l'évolution du SINR moyen au cours de l'entraînement. On observe une stabilisation progressive du SINR à mesure que les agents convergent vers des puissances cohérentes, réduisant les fluctuations d'interférence mutuelle.

\begin{figure}[htbp]
    \centering
    \includegraphics[width=0.75\textwidth]{figures/training_sinr_evolution.png}
    \caption{Évolution du SINR moyen (dB) durant l'entraînement. La stabilisation du SINR reflète l'émergence d'un équilibre entre les agents.}
    \label{fig:training_sinr}
\end{figure}

% ============================================================
\section{Comparaison multi-critères}
\label{sec:multi_criteria}

Le Tableau~\ref{tab:multi_criteria} présente une comparaison synthétique sur l'ensemble des critères de performance.

\begin{table*}[htbp]
    \centering
    \caption{Comparaison multi-critères des méthodes de contrôle de puissance}
    \label{tab:multi_criteria}
    \renewcommand{\arraystretch}{1.2}
    \begin{tabular}{@{}lccccc@{}}
        \toprule
        \textbf{Critère} & \textbf{DQN} & \textbf{Article} & \textbf{Opt. Unif.} & \textbf{Max Power} & \textbf{Random} \\
        \midrule
        EE (Mbit/J) $\uparrow$      & \textbf{255,16} & 207,01 & 213,04 & 10,22 & 23,06 \\
        Puissance (mW) $\downarrow$  & \textbf{15,27}  & 28,73  & 0,02   & 199,53 & 100,57 \\
        Stabilité ($1/\sigma$) $\uparrow$ & 0,021 & \textbf{0,033} & 0,019 & 0,392 & 0,064 \\
        Connaissance modèle          & \ding{55}       & \ding{51} & \ding{55} & \ding{55} & \ding{55} \\
        Décentralisé                  & \ding{51}       & \ding{51} & \ding{55} & N/A   & N/A \\
        Adaptabilité                  & \textbf{Haute}  & Basse    & Nulle & Nulle & Nulle \\
        \bottomrule
    \end{tabular}
\end{table*}

\begin{figure}[htbp]
    \centering
    \includegraphics[width=0.7\textwidth]{figures/radar_chart_comparison.png}
    \caption{Diagramme radar comparatif des méthodes sur quatre critères normalisés : EE, économie de puissance (~$1/P$), stabilité ($1/\sigma$), et débit spectral.}
    \label{fig:radar_comparison}
\end{figure}

% ============================================================
\section{Discussion}
\label{sec:discussion}

\subsection{Avantages du DQN}

\begin{enumerate}
    \item \textbf{Performances supérieures} : le DQN surpasse l'algorithme analytique de $+23{,}3\%$ en EE tout en consommant $46{,}8\%$ moins de puissance. Cet avantage est remarquable car il est obtenu \textit{sans connaissance explicite du modèle de canal}.

    \item \textbf{Absence de modèle} (\textit{model-free}) : contrairement à l'approche BRD + Dinkelbach qui nécessite la connaissance exacte de la structure du SINR et la résolution analytique de l'optimisation fractionnaire, le DQN apprend directement à partir des observations et des récompenses. Cela le rend robuste aux inexactitudes de modélisation et adaptable à des scénarios non couverts par la théorie.

    \item \textbf{Adaptabilité} : les agents DQN peuvent s'adapter à des changements dans l'environnement (nombre d'utilisateurs, topologie, conditions de propagation) par ré-entraînement ou \textit{fine-tuning}, sans nécessiter de re-dérivation des solutions analytiques.
\end{enumerate}

\subsection{Limitations et travaux futurs}

\begin{enumerate}
    \item \textbf{Variance plus élevée} : l'écart-type de l'EE du DQN ($47{,}31$~Mbit/J) est supérieur à celui de l'algorithme de l'article ($30{,}67$~Mbit/J), indiquant une plus grande variabilité des performances selon les réalisations de canal. Cela pourrait être atténué par des techniques de régularisation ou un entraînement plus long.

    \item \textbf{Coût d'entraînement} : l'entraînement de $10\,000$ épisodes requiert un temps significatif, notamment avec $K = 5$ agents indépendants. Cependant, ce coût est assumé hors-ligne (\textit{offline}) et la politique résultante est exécutée en temps réel avec une simple passe avant dans le réseau de neurones.

    \item \textbf{Garanties théoriques} : contrairement à l'algorithme de l'article qui fournit des garanties de convergence vers le GNE, l'approche IQL n'a pas de preuve formelle de convergence dans le cas multi-agent non-stationnaire. Néanmoins, les résultats empiriques démontrent une convergence effective.

    \item \textbf{Passage à l'échelle} : des travaux futurs pourraient explorer des architectures de \textit{parameter sharing} entre agents ou des approches de type \textit{Mean Field} pour étendre le système à un plus grand nombre d'utilisateurs.
\end{enumerate}

% ============================================================
\section{Conclusion du chapitre}

Les résultats expérimentaux démontrent que l'approche par Deep Q-Network constitue une alternative viable et performante à l'algorithme analytique de Zappone \textit{et al.} pour le contrôle de puissance décentralisé dans les réseaux 5G Massive MIMO. Avec un gain de $+23{,}3\%$ en efficacité énergétique et une réduction de $46{,}8\%$ en consommation de puissance, le DQN valide le potentiel de l'apprentissage par renforcement profond pour l'optimisation des ressources radio. L'avantage principal réside dans sa nature \textit{model-free}, qui le rend applicable à des scénarios réalistes où le modèle mathématique exact du canal n'est pas disponible ou est trop complexe pour être exploité analytiquement.

% ============================================================
% BIBLIOGRAPHIE (extraits pertinents)
% ============================================================
% À inclure dans le fichier .bib principal :
%
% @article{zappone2015,
%   author  = {Zappone, A. and Björnson, E. and Sanguinetti, L. and Jorswieck, E.},
%   title   = {Energy-Efficient Power Control: A Look at 5G Wireless Technologies},
%   journal = {IEEE Transactions on Signal Processing},
%   year    = {2015},
%   volume  = {64},
%   number  = {7},
%   pages   = {1668--1683}
% }
%
% @article{mnih2015,
%   author  = {Mnih, V. and Kavukcuoglu, K. and Silver, D. and others},
%   title   = {Human-level control through deep reinforcement learning},
%   journal = {Nature},
%   year    = {2015},
%   volume  = {518},
%   pages   = {529--533}
% }
%
% @inproceedings{vanHasselt2016,
%   author  = {van Hasselt, H. and Guez, A. and Silver, D.},
%   title   = {Deep Reinforcement Learning with Double Q-learning},
%   booktitle = {AAAI},
%   year    = {2016}
% }
%
% @article{kingma2015,
%   author  = {Kingma, D. P. and Ba, J.},
%   title   = {Adam: A Method for Stochastic Optimization},
%   journal = {ICLR},
%   year    = {2015}
% }
%
% @inproceedings{glorot2010,
%   author  = {Glorot, X. and Bengio, Y.},
%   title   = {Understanding the difficulty of training deep feedforward neural networks},
%   booktitle = {AISTATS},
%   year    = {2010}
% }
%
% @inproceedings{tan1993,
%   author  = {Tan, M.},
%   title   = {Multi-Agent Reinforcement Learning: Independent vs. Cooperative Agents},
%   booktitle = {ICML},
%   year    = {1993}
% }
%
% @article{lanctot2017,
%   author  = {Lanctot, M. and Zambaldi, V. and Gruslys, A. and others},
%   title   = {A Unified Game-Theoretic Approach to Multiagent Reinforcement Learning},
%   journal = {NeurIPS},
%   year    = {2017}
% }


%==============================================================================
\section{Conclusion et perspectives}
%==============================================================================

[À développer]

\begin{thebibliography}{99}

\bibitem{3GPP} 3GPP, "Technical Specification Group Radio Access Network; NR; Physical layer; General description (Release 15)," 3GPP TS 38.201, 2018.

\bibitem{Fehske2011} A. Fehske, J. Malmodin, G. Biczok, and G. Fettweis, "The global footprint of mobile communications: The ecological and economic perspective," \textit{IEEE Commun. Mag.}, pp. 55-62, Aug. 2011.

\bibitem{Ericsson2011} Ericsson White Paper, "More than 50 billion connected devices," Ericsson, Tech. Rep. 284 23-3149 Uen, Feb. 2011.

\bibitem{Qualcomm} Qualcomm, "The 1000x data challenge," Tech. Rep.

\bibitem{Zappone2021} A. Zappone, L. Sanguinetti, G. Bacci, E. Jorswieck, and M. Debbah, "Energy-Efficient Power Control: A Look at 5G Wireless Technologies," \textit{IEEE Transactions on Signal Processing}, 2021.

\bibitem{Dinkelbach1967} W. Dinkelbach, "On nonlinear fractional programming," \textit{Management Science}, vol. 13, no. 7, pp. 492-498, Mar. 1967.

\bibitem{Ng2012} D. W. K. Ng, E. S. Lo, and R. Schober, "Energy-efficient resource allocation in multi-cell OFDMA systems with limited backhaul capacity," \textit{IEEE Trans. Wireless Commun.}, vol. 11, no. 10, pp. 3618-3631, Oct. 2012.

\bibitem{Xu2014} Q. Xu, X. Li, H. Ji, and X. Du, "Energy-efficient resource allocation for heterogeneous services in OFDMA downlink networks: Systematic perspective," \textit{IEEE Trans. Veh. Technol.}, vol. 63, no. 5, pp. 2071-2082, June 2014.

\bibitem{Xu2013} J. Xu and L. Qiu, "Energy efficiency optimization for MIMO broadcast channels," \textit{IEEE Trans. Wireless Commun.}, vol. 12, no. 2, pp. 690-701, Feb. 2013.

\bibitem{Du2014} B. Du, C. Pan, W. Zhang, and M. Chen, "Distributed energy-efficient power optimization for CoMP systems with max-min fairness," \textit{IEEE Commun. Lett.}, vol. 18, no. 6, pp. 999-1002, Jun. 2014.

\bibitem{He2013} S. He, Y. Huang, S. Jin, and L. Yang, "Coordinated beamforming for energy efficient transmission in multicell multiuser systems," \textit{IEEE Trans. Commun.}, vol. 61, no. 12, pp. 4961-4971, Dec. 2013.

\bibitem{He2014} S. He, Y. Huang, L. Yang, and B. Ottersten, "Coordinated multicell multiuser precoding for maximizing weighted sum energy efficiency," \textit{IEEE Trans. Signal Process.}, vol. 62, no. 3, pp. 741-751, Feb. 2014.

\bibitem{Venturino2015} L. Venturino, A. Zappone, C. Risi, and S. Buzzi, "Energy-efficient scheduling and power allocation in downlink OFDMA networks with base station coordination," \textit{IEEE Trans. Wireless Commun.}, vol. 14, no. 1, pp. 1-14, 2015.

\bibitem{Miao2011} G. Miao, N. Himayat, G. Li, and S. Talwar, "Distributed interference-aware energy-efficient power optimization," \textit{IEEE Trans. Wireless Commun.}, vol. 10, no. 4, pp. 1323-1333, Apr. 2011.

\bibitem{Scutari2008} G. Scutari, D. Palomar, and S. Barbarossa, "Optimal linear precoding strategies for wideband noncooperative systems based on game theory part I: Nash equilibria," \textit{IEEE Trans. Signal Process.}, vol. 56, no. 3, pp. 1230-1249, March 2008.

\bibitem{Zappone2013} A. Zappone, Z. Chong, E. Jorswieck, and S. Buzzi, "Energy-aware competitive power control in relay-assisted interference wireless networks," \textit{IEEE Trans. Wireless Commun.}, vol. 12, no. 4, pp. 1860-1871, Apr. 2013.

\bibitem{Stupia2015} I. Stupia, L. Sanguinetti, G. Bacci, and L. Vandendorpe, "Power control in networks with heterogeneous users: A quasi-variational inequality approach," \textit{IEEE Trans. Signal Process.}, vol. 63, no. 21, pp. 5691-5705, November 2015.

\bibitem{Bacci2015} G. Bacci, E. Belmega, P. Mertikopoulos, and L. Sanguinetti, "Energy-aware competitive power allocation for heterogeneous networks under QoS constraints," \textit{IEEE Trans. Wireless Commun.}, vol. 14, no. 9, pp. 4728-4742, Sept 2015.

\bibitem{Yates1995} R. D. Yates, "A framework for uplink power control in cellular radio systems," \textit{IEEE J. Sel. Areas Commun.}, vol. 13, no. 7, pp. 1341-1347, September 1995.

\end{thebibliography}

\end{document}
% Prérequis packages : amsmath, amssymb, algorithm2e, graphicx,
%   booktabs, multirow, tabularx, subcaption, hyperref, xcolor

% ============================================================
%            CHAPITRE : APPROCHE PAR DEEP Q-NETWORK
% ============================================================
\section{Contrôle de Puissance par Deep Q-Network}
\label{chap:dqn}

\subsection{Introduction}

Le chapitre précédent a présenté l'approche analytique de Zappone \textit{et al.}~\cite{zappone2015} pour la maximisation décentralisée de l'efficacité énergétique (EE) dans les réseaux 5G Massive MIMO, reposant sur la dynamique de meilleure réponse (BRD) couplée à l'algorithme de Dinkelbach. Bien que cette méthode converge vers l'Équilibre de Nash Généralisé (GNE) sous certaines conditions, elle nécessite une connaissance explicite du modèle de canal et la résolution itérative d'un problème d'optimisation fractionnaire à chaque mise à jour.

Dans ce chapitre, nous proposons une approche alternative basée sur l'apprentissage par renforcement profond (\textit{Deep Reinforcement Learning}), et plus spécifiquement sur le \textbf{Deep Q-Network (DQN)}~\cite{mnih2015}. L'objectif est de permettre à chaque utilisateur $k$ d'apprendre de manière autonome une politique de contrôle de puissance qui maximise son EE individuelle, sans connaissance explicite du modèle mathématique du canal. L'approche est intrinsèquement décentralisée : chaque utilisateur dispose de son propre agent DQN qui prend des décisions indépendantes, tout en interagissant avec les autres agents via l'interférence partagée dans l'environnement sans fil.

% ============================================================
\subsection{Formulation en Processus Décisionnel de Markov}
\label{sec:mdp}

Pour appliquer l'apprentissage par renforcement, nous reformulons le problème de contrôle de puissance comme un \textbf{Processus Décisionnel de Markov} (MDP). Chaque utilisateur $k \in \{1, \dots, K\}$ est modélisé comme un agent autonome interagissant avec un environnement partagé. Le MDP de l'utilisateur $k$ est défini par le tuple $(\mathcal{S}_k, \mathcal{A}_k, \mathcal{P}, \mathcal{R}_k, \gamma)$, dont les composantes sont détaillées ci-après.

\subsubsection{Espace d'états}
\label{sec:state_space}

L'état observé par l'utilisateur $k$ à l'instant $t$ est un vecteur de dimension 6 :
\begin{equation}
    \mathbf{s}_k^{(t)} = \begin{bmatrix}
        s_1^{(t)} \\ s_2^{(t)} \\ s_3^{(t)} \\ s_4^{(t)} \\ s_5^{(t)} \\ s_6^{(t)}
    \end{bmatrix}
    \label{eq:state_vector}
\end{equation}
dont les composantes sont définies comme suit :

\begin{enumerate}
    \item \textbf{SINR normalisé en dB} :
    \begin{equation}
        s_1^{(t)} = \text{clip}\!\left(\frac{\text{SINR}_k^{\text{dB}}}{25}, -1, 1\right), \quad \text{SINR}_k^{\text{dB}} = 10\log_{10}(\gamma_k^{(t)})
        \label{eq:state_sinr}
    \end{equation}
    où la normalisation par 25~dB correspond à la plage pratique du SINR dans le système considéré.

    \item \textbf{Puissance normalisée (échelle logarithmique)} :
    \begin{equation}
        s_2^{(t)} = \frac{\log_{10}(p_k^{(t)}) - \log_{10}(p_{\min})}{\log_{10}(p_{\max}) - \log_{10}(p_{\min})}
        \label{eq:state_power}
    \end{equation}
    Cette représentation en échelle log permet une résolution uniforme sur toute la plage de puissance.

    \item \textbf{Interférence normalisée} :
    \begin{equation}
        s_3^{(t)} = \frac{\log_{10}\left(\sum_{j \neq k} p_j^{(t)}\right) - \log_{10}\left((K\!-\!1)\,p_{\min}\right)}{\log_{10}\left((K\!-\!1)\,p_{\max}\right) - \log_{10}\left((K\!-\!1)\,p_{\min}\right)}
        \label{eq:state_interf}
    \end{equation}

    \item \textbf{Gain de canal équivalent normalisé} :
    \begin{equation}
        s_4^{(t)} = \text{clip}\!\left(\frac{\log_{10}(\mu_k^{(t)})}{10}, 0, 1\right)
        \label{eq:state_mu}
    \end{equation}
    où $\mu_k = \frac{\alpha_k}{\sigma^2 + \sum_{j \neq k} \omega_{kj}\,p_j}$ est le gain de canal équivalent utilisé dans la meilleure réponse de Dinkelbach. Cette composante est cruciale car elle encode l'information nécessaire pour déterminer la puissance optimale.

    \item \textbf{Débit normalisé} :
    \begin{equation}
        s_5^{(t)} = \frac{R_k^{(t)}}{R_{\max}}, \quad R_{\max} = \log_2(1 + \bar{\gamma})
        \label{eq:state_rate}
    \end{equation}
    où $R_{\max}$ est le débit maximal théorique limité par les imperfections matérielles ($\bar{\gamma} = 1/\varepsilon^2$).

    \item \textbf{Marge QoS} :
    \begin{equation}
        s_6^{(t)} = \text{clip}\!\left(\frac{R_k^{(t)} - \theta_k}{\theta_k}, -1, 1\right)
        \label{eq:state_qos}
    \end{equation}
    où $\theta_k$ est le débit minimal requis par la contrainte de qualité de service. Une valeur positive indique que la contrainte est satisfaite.
\end{enumerate}

Toutes les composantes sont normalisées dans l'intervalle $[-1, 1]$ ou $[0, 1]$ afin de faciliter l'apprentissage du réseau de neurones. Le choix d'une représentation en échelle logarithmique pour la puissance et l'interférence reflète la nature multiplicative des signaux radio et améliore la sensibilité de l'agent aux variations dans les régimes de faible puissance, où se situe typiquement le point de fonctionnement optimal en termes d'EE.

\subsubsection{Espace d'actions}
\label{sec:action_space}

L'espace d'actions est \textbf{discret} et composé de $|\mathcal{A}| = 32$ niveaux de puissance répartis de manière logarithmique entre $p_{\min}$ et $p_{\max}$ :
\begin{equation}
    \mathcal{A} = \left\{ p_{\min} \cdot \left(\frac{p_{\max}}{p_{\min}}\right)^{\frac{i}{N_a - 1}} \;\middle|\; i = 0, 1, \dots, N_a - 1 \right\}
    \label{eq:action_space}
\end{equation}
avec $p_{\min} = -20$~dBm ($0{,}01$~mW) et $p_{\max} = 23$~dBm ($200$~mW), soit un rapport de 1:20\,000.

L'espacement logarithmique (géométrique) est motivé par le fait que l'efficacité énergétique optimale se situe typiquement dans la partie inférieure de la plage de puissance. Cette discrétisation offre ainsi une \textbf{résolution fine} dans la zone optimale (entre 0,01 et 5~mW) tout en couvrant l'ensemble de la plage autorisée. L'action $a_k^{(t)} \in \{0, 1, \dots, 31\}$ sélectionne directement le niveau de puissance d'émission $p_k^{(t)} = \mathcal{A}[a_k^{(t)}]$.

\subsubsection{Fonction de récompense}
\label{sec:reward}

La fonction de récompense est conçue pour guider l'agent vers la maximisation de l'efficacité énergétique tout en respectant les contraintes de QoS et en encourageant l'économie de puissance :

\begin{equation}
    r_k^{(t)} = \text{clip}\!\left(\alpha_{\text{ee}} \cdot \frac{\eta_k^{(t)}}{\eta_{\text{ref}}} - \beta_{\text{qos}} \cdot \mathbb{1}_{\text{QoS}} \cdot \left(0{,}3 + 0{,}7 \cdot \min\!\left(\frac{\theta_k - R_k^{(t)}}{\theta_k}, 1\right)\right) - \gamma_p \cdot \frac{p_k^{(t)}}{p_{\max}},\; -3,\; 2\right)
    \label{eq:reward}
\end{equation}

où :
\begin{itemize}
    \item $\eta_k^{(t)} = \frac{B \cdot \log_2(1 + \gamma_k^{(t)})}{p_c + p_k^{(t)}}$ est l'efficacité énergétique en bit/Joule ;
    \item $\eta_{\text{ref}} = 3 \times 10^8$ bit/J (= 300~Mbit/J) est la référence de normalisation, choisie pour que le terme EE soit dans l'intervalle $[0, 1]$ au voisinage de la solution optimale ;
    \item $\alpha_{\text{ee}} = 1{,}0$ est le coefficient de récompense EE ;
    \item $\beta_{\text{qos}} = 5{,}0$ est le coefficient de pénalité QoS, volontairement élevé pour fortement décourager les violations de la contrainte de débit minimal ;
    \item $\gamma_p = 0{,}02$ est le coefficient de pénalité de puissance, encourageant l'économie d'énergie ;
    \item $\mathbb{1}_{\text{QoS}} = \mathbb{1}_{R_k < \theta_k}$ est l'indicatrice de violation QoS.
\end{itemize}

La pénalité QoS est composée d'un terme fixe ($0{,}3$) activé dès qu'il y a violation, et d'un terme proportionnel au déficit relatif $(\theta_k - R_k)/\theta_k$, ce qui permet de différencier les violations légères des violations sévères. Le \textit{soft clipping} dans $[-3, 2]$ assure la stabilité numérique de l'entraînement en évitant les récompenses extrêmes.

\subsubsection{Dynamique de l'environnement}
\label{sec:env_dynamics}

L'environnement simule un réseau 5G Massive MIMO avec les caractéristiques suivantes :

\begin{itemize}
    \item \textbf{Canal dynamique} : à chaque pas de temps, le canal est mis à jour par un évanouissement rapide résiduel :
    \begin{equation}
        d_k^{\text{eff}}(t) = d_k \cdot (1 + \delta_k(t)), \quad \delta_k \sim \mathcal{N}\!\left(0, \frac{1}{\sqrt{M}}\right)
        \label{eq:channel_hardening}
    \end{equation}
    Ce modèle capture l'effet de \textit{channel hardening} typique du Massive MIMO avec $M = 50$ antennes, donnant des variations de $\pm 1$--$2$~dB.

    \item \textbf{Réinitialisation par épisode} : à chaque épisode, les $K = 5$ utilisateurs sont repositionnés aléatoirement dans la cellule (rayon $R_{\text{cell}} = 400$~m, distance minimale $d_{\min} = 35$~m), et un nouveau réalisation de l'affaiblissement grande échelle (perte de trajet + \textit{shadowing}) est tirée :
    \begin{equation}
        d_k = 10^{-(L_k + S_k)/10}, \quad L_k = 10\alpha\log_{10}(r_k), \quad S_k \sim \mathcal{N}(0, \sigma_s^2)
        \label{eq:large_scale_fading}
    \end{equation}
    avec $\alpha = 3{,}76$ et $\sigma_s = 8$~dB.

    \item \textbf{Durée d'épisode} : chaque épisode comprend $T = 200$ pas de temps, durant lesquels les agents interagissent simultanément.
\end{itemize}

% ============================================================
\subsection{Architecture du Deep Q-Network}
\label{sec:dqn_architecture}

\subsubsection{Principe du DQN}

Le Deep Q-Network~\cite{mnih2015} approxime la fonction de valeur action-état $Q^*(s, a)$ par un réseau de neurones profond $Q(s, a; \boldsymbol{\theta})$ paramétré par $\boldsymbol{\theta}$. La politique gloutonne (\textit{greedy}) correspondante est :
\begin{equation}
    \pi^*(s) = \arg\max_{a \in \mathcal{A}} Q(s, a; \boldsymbol{\theta})
    \label{eq:greedy_policy}
\end{equation}

L'objectif est de minimiser l'erreur temporelle différentielle (\textit{TD error}) entre la valeur Q courante et la valeur cible :
\begin{equation}
    \mathcal{L}(\boldsymbol{\theta}) = \mathbb{E}_{(s,a,r,s') \sim \mathcal{D}} \left[ \ell\!\left(r + \gamma \max_{a'} Q(s', a'; \boldsymbol{\theta}^{-}) - Q(s, a; \boldsymbol{\theta})\right) \right]
    \label{eq:dqn_loss}
\end{equation}
où $\mathcal{D}$ est le \textit{replay buffer}, $\boldsymbol{\theta}^{-}$ les paramètres du réseau cible, $\gamma = 0{,}95$ le facteur d'actualisation, et $\ell(\cdot)$ la fonction de perte.

\subsubsection{Architecture du réseau de neurones}

Le réseau de neurones utilisé pour chaque agent DQN est un perceptron multicouche (\textit{Multi-Layer Perceptron}, MLP) composé de trois couches cachées entièrement connectées (\textit{fully connected}). L'architecture est la suivante :

\begin{equation}
    \underbrace{\mathbb{R}^6}_{\text{Entrée (état)}} \xrightarrow{\mathbf{W}_1, \mathbf{b}_1} \underbrace{\mathbb{R}^{128}}_{\text{Couche 1}} \xrightarrow{\text{ReLU}} \underbrace{\mathbb{R}^{256}}_{\text{Couche 2}} \xrightarrow{\text{ReLU}} \underbrace{\mathbb{R}^{128}}_{\text{Couche 3}} \xrightarrow{\text{ReLU}} \underbrace{\mathbb{R}^{32}}_{\text{Sortie (Q-values)}}
    \label{eq:architecture}
\end{equation}

\begin{itemize}
    \item \textbf{Couche d'entrée} : 6 neurones correspondant aux 6 composantes de l'état $\mathbf{s}_k$.
    \item \textbf{Couche cachée 1} : 128 neurones avec activation ReLU.
    \item \textbf{Couche cachée 2} : 256 neurones avec activation ReLU.
    \item \textbf{Couche cachée 3} : 128 neurones avec activation ReLU.
    \item \textbf{Couche de sortie} : 32 neurones (un par action), sans activation — les sorties représentent directement les Q-valeurs $Q(s, a_i; \boldsymbol{\theta})$ pour $i = 0, \dots, 31$.
\end{itemize}

L'architecture en forme de \textit{bottleneck} (128 $\to$ 256 $\to$ 128) permet d'abord une expansion de la représentation dans l'espace latent, puis une compression vers la dimension de sortie, favorisant l'extraction de caractéristiques pertinentes.

Les poids sont initialisés selon la méthode de \textbf{Xavier uniforme}~\cite{glorot2010} :
\begin{equation}
    W_{ij} \sim \mathcal{U}\left(-\sqrt{\frac{6}{n_{\text{in}} + n_{\text{out}}}},\; +\sqrt{\frac{6}{n_{\text{in}} + n_{\text{out}}}}\right)
    \label{eq:xavier_init}
\end{equation}
et les biais sont initialisés à zéro. Le nombre total de paramètres entraînables par agent est de :
\begin{equation}
    N_{\text{params}} = (6 \times 128 + 128) + (128 \times 256 + 256) + (256 \times 128 + 128) + (128 \times 32 + 32) = 70\,944
    \label{eq:n_params}
\end{equation}

\begin{figure*}[htbp]
    \centering
    \includegraphics[width=0.85\textwidth]{figures/dqn_architecture.png}
    \caption{Architecture du réseau DQN pour le contrôle de puissance. L'entrée est le vecteur d'état à 6 dimensions et la sortie fournit les Q-valeurs pour les 32 niveaux de puissance discrets.}
    \label{fig:dqn_architecture}
\end{figure*}

% ============================================================
\subsection{Mécanismes d'apprentissage}
\label{sec:learning_mechanisms}

\subsubsection{Double DQN}
\label{sec:double_dqn}

Afin de réduire la surestimation des Q-valeurs inhérente au DQN standard~\cite{vanHasselt2016}, nous implémentons la variante \textbf{Double DQN}. Le principe consiste à découpler la \textbf{sélection} de l'action (effectuée par le réseau de politique $\boldsymbol{\theta}$) de son \textbf{évaluation} (effectuée par le réseau cible $\boldsymbol{\theta}^{-}$) :

\begin{equation}
    y^{\text{DDQN}} = r + \gamma \cdot Q\!\left(s',\; \underbrace{\arg\max_{a'} Q(s', a'; \boldsymbol{\theta})}_{\text{sélection par réseau politique}};\; \boldsymbol{\theta}^{-}\right)
    \label{eq:double_dqn_target}
\end{equation}

Cette technique réduit significativement le biais d'estimation positif et améliore la stabilité de l'apprentissage, ce qui est particulièrement important dans notre contexte multi-agent où les environnements sont non-stationnaires du point de vue de chaque agent.

\subsubsection{Replay d'expérience}
\label{sec:experience_replay}

Chaque agent maintient un \textbf{tampon de rejeu} (\textit{experience replay buffer}) $\mathcal{D}_k$ de capacité $|\mathcal{D}| = 100\,000$ transitions. Chaque transition stockée est un tuple :
\begin{equation}
    e^{(t)} = \left(s_k^{(t)},\; a_k^{(t)},\; r_k^{(t)},\; s_k^{(t+1)},\; d^{(t)}\right)
    \label{eq:transition}
\end{equation}
où $d^{(t)} \in \{0, 1\}$ est l'indicateur de fin d'épisode. Le tampon fonctionne selon le principe FIFO (\textit{First-In, First-Out}) : lorsque la capacité maximale est atteinte, les transitions les plus anciennes sont remplacées.

À chaque pas d'entraînement, un \textbf{mini-batch} de $B = 256$ transitions est échantillonné uniformément de manière aléatoire :
\begin{equation}
    \{(s_i, a_i, r_i, s'_i, d_i)\}_{i=1}^{B} \sim \text{Uniforme}(\mathcal{D}_k)
    \label{eq:sampling}
\end{equation}

Le replay d'expérience confère deux avantages fondamentaux :
\begin{itemize}
    \item \textbf{Décorrélation temporelle} : les transitions échantillonnées aléatoirement brisent les corrélations séquentielles, satisfaisant l'hypothèse i.i.d. de la descente de gradient stochastique.
    \item \textbf{Efficacité en données} : chaque transition peut être réutilisée dans plusieurs mises à jour, améliorant l'exploitation des données collectées.
\end{itemize}

L'entraînement ne commence qu'une fois le tampon contenant au moins $|\mathcal{D}|_{\min} = 2\,000$ transitions, garantissant une diversité suffisante dans les premiers mini-batchs.

\subsubsection{Réseau cible et mise à jour douce}
\label{sec:target_network}

Pour stabiliser l'apprentissage, nous utilisons un \textbf{réseau cible} (\textit{target network}) $Q(s, a; \boldsymbol{\theta}^{-})$ distinct du réseau de politique $Q(s, a; \boldsymbol{\theta})$. Les paramètres du réseau cible sont mis à jour par \textbf{mise à jour douce} (\textit{soft update} ou \textit{Polyak averaging}) :

\begin{equation}
    \boldsymbol{\theta}^{-} \leftarrow \tau \cdot \boldsymbol{\theta} + (1 - \tau) \cdot \boldsymbol{\theta}^{-}, \quad \tau = 0{,}005
    \label{eq:soft_update}
\end{equation}

Cette mise à jour est effectuée après chaque pas d'optimisation. Le coefficient $\tau = 0{,}005$ assure une transition lente et régulière du réseau cible vers le réseau de politique, évitant les oscillations et divergences liées à une cible en mouvement rapide. Par rapport à la mise à jour dure (\textit{hard update}) périodique, la mise à jour douce offre une stabilité accrue car le réseau cible évolue de manière continue plutôt que par sauts discontinus.

\subsubsection{Fonction de perte de Huber}
\label{sec:huber_loss}

La fonction de perte utilisée pour l'optimisation est la \textbf{perte de Huber} (aussi appelée \textit{Smooth L1 Loss}) :
\begin{equation}
    \ell_{\delta}(x) = \begin{cases}
        \frac{1}{2}x^2 & \text{si } |x| \leq \delta \\
        \delta\left(|x| - \frac{1}{2}\delta\right) & \text{sinon}
    \end{cases}
    \label{eq:huber_loss}
\end{equation}
avec $\delta = 1{,}0$. La perte de Huber combine les avantages de la MSE (sensibilité aux petites erreurs) et de la MAE (robustesse aux valeurs aberrantes), ce qui est crucial dans notre environnement multi-agent où les récompenses peuvent être bruitées par les actions concurrentes des autres agents.

\subsubsection{Politique d'exploration $\varepsilon$-gloutonne}
\label{sec:epsilon_greedy}

L'exploration du space d'actions est réalisée par une politique $\varepsilon$-gloutonne (\textit{$\varepsilon$-greedy}) :
\begin{equation}
    a_k^{(t)} = \begin{cases}
        \text{Uniforme}(\mathcal{A}) & \text{avec probabilité } \varepsilon_t \\
        \arg\max_{a} Q(s_k^{(t)}, a; \boldsymbol{\theta}_k) & \text{avec probabilité } 1 - \varepsilon_t
    \end{cases}
    \label{eq:epsilon_greedy}
\end{equation}

Le taux d'exploration $\varepsilon_t$ décroît de manière exponentielle au fil des épisodes :
\begin{equation}
    \varepsilon_t = \max\!\left(\varepsilon_{\text{fin}},\; \varepsilon_0 \cdot \lambda^t\right), \quad \varepsilon_0 = 1{,}0, \quad \varepsilon_{\text{fin}} = 0{,}02, \quad \lambda = 0{,}9995
    \label{eq:epsilon_decay}
\end{equation}

Ce schéma de décroissance assure :
\begin{itemize}
    \item Une \textbf{phase d'exploration intensive} initiale (épisodes 1--1\,000, $\varepsilon > 0{,}6$) permettant une couverture large de l'espace état-action.
    \item Une \textbf{transition progressive} (épisodes 1\,000--7\,000) où l'agent bascule graduellement vers l'exploitation.
    \item Une \textbf{phase d'exploitation} finale (épisodes 7\,000--10\,000, $\varepsilon \approx 0{,}02$) où l'agent utilise quasi-exclusivement la politique apprise, tout en conservant une exploration résiduelle de 2\% pour maintenir une certaine adaptabilité.
\end{itemize}

L'épisode où $\varepsilon$ passe sous le seuil $0{,}1$ (exploitation dominante) se situe aux alentours de l'épisode $4\,600$, ce qui correspond au régime où les Q-valeurs sont suffisamment stables pour guider les décisions.

\begin{figure}[htbp]
    \centering
    \includegraphics[width=1\textwidth]{figures/epsilon_decay.png}
    \caption{Décroissance exponentielle du taux d'exploration $\varepsilon$ au cours des $10\,000$ épisodes d'entraînement. La ligne pointillée indique le seuil $\varepsilon = 0{,}1$.}
    \label{fig:epsilon_decay}
\end{figure}

\subsubsection{Gradient clipping}
\label{sec:grad_clip}

Pour prévenir l'instabilité numérique causée par des gradients de grande magnitude, une \textbf{troncature du gradient} (\textit{gradient clipping}) est appliquée avant chaque mise à jour des poids :
\begin{equation}
    \nabla_{\boldsymbol{\theta}} \mathcal{L} \leftarrow \frac{\nabla_{\boldsymbol{\theta}} \mathcal{L}}{\max\!\left(1, \frac{\|\nabla_{\boldsymbol{\theta}} \mathcal{L}\|_2}{G_{\max}}\right)}, \quad G_{\max} = 1{,}0
    \label{eq:grad_clip}
\end{equation}

Cette technique est d'autant plus importante dans le contexte multi-agent, où la non-stationnarité induite par l'apprentissage simultané de tous les agents peut engendrer des gradients instables.

% ============================================================
\subsection{Entraînement multi-agent décentralisé}
\label{sec:multi_agent}

\subsubsection{Independent Q-Learning (IQL)}

L'approche adoptée est le \textbf{Independent Q-Learning} (IQL)~\cite{tan1993}, où chaque utilisateur $k$ possède son propre agent DQN qui apprend indépendamment. Les agents interagissent uniquement à travers l'environnement partagé : le SINR de l'utilisateur $k$ dépend des puissances d'émission de tous les utilisateurs via l'interférence :
\begin{equation}
    \gamma_k = \frac{\alpha_k \, p_k}{\sigma^2 + \phi_k \, p_k + \sum_{j \neq k} \omega_{kj} \, p_j}
    \label{eq:sinr_coupling}
\end{equation}

Du point de vue de l'utilisateur $k$, les autres utilisateurs font partie de l'environnement, qui devient \textbf{non-stationnaire} puisque les politiques des autres agents évoluent au cours de l'entraînement. Bien que cela viole formellement les hypothèses de stationnarité du MDP, l'IQL fonctionne remarquablement bien en pratique pour les jeux coopératifs et compétitifs de taille modérée~\cite{lanctot2017}.

\subsubsection{Protocole d'entraînement}

Le protocole d'entraînement est résumé dans l'Algorithme~\ref{alg:training}.

\begin{algorithm}[htbp]
\caption{Entraînement multi-agent DQN décentralisé}
\label{alg:training}
\SetKwInOut{Input}{Entrées}
\SetKwInOut{Output}{Sorties}
\Input{$K = 5$ agents, $N_{\text{ep}} = 10\,000$ épisodes, $T = 200$ pas/épisode}
\Output{Politiques optimales $\{\boldsymbol{\theta}_k^*\}_{k=1}^{K}$}
\BlankLine
Initialiser $K$ réseaux de politique $\{Q(\cdot; \boldsymbol{\theta}_k)\}$ et cibles $\{Q(\cdot; \boldsymbol{\theta}_k^{-})\}$\;
Initialiser $K$ tampons de rejeu $\{\mathcal{D}_k\}$ de capacité $100\,000$\;
$\varepsilon \leftarrow 1{,}0$\;
\BlankLine
\For{épisode $= 1$ \KwTo $N_{\text{ep}}$}{
    Réinitialiser l'environnement (nouvelles positions + canal)\;
    Observer les états initiaux $\{s_k^{(0)}\}_{k=1}^{K}$\;
    \BlankLine
    \For{$t = 0$ \KwTo $T-1$}{
        \For{$k = 1$ \KwTo $K$ \textrm{(en parallèle)}}{
            Sélectionner $a_k^{(t)}$ par politique $\varepsilon$-gloutonne (Éq.~\ref{eq:epsilon_greedy})\;
        }
        Exécuter les actions $\{a_k^{(t)}\}$ dans l'environnement\;
        Observer récompenses $\{r_k^{(t)}\}$ et nouveaux états $\{s_k^{(t+1)}\}$\;
        \BlankLine
        \For{$k = 1$ \KwTo $K$}{
            Stocker $(s_k^{(t)}, a_k^{(t)}, r_k^{(t)}, s_k^{(t+1)}, d^{(t)})$ dans $\mathcal{D}_k$\;
            \If{$|\mathcal{D}_k| \geq 2\,000$ \textbf{et} $t \bmod 16 = 0$}{
                Échantillonner un mini-batch de 256 transitions de $\mathcal{D}_k$\;
                Calculer les cibles Double DQN (Éq.~\ref{eq:double_dqn_target})\;
                Mettre à jour $\boldsymbol{\theta}_k$ par descente de gradient (perte de Huber)\;
                Mise à jour douce : $\boldsymbol{\theta}_k^{-} \leftarrow 0{,}005 \cdot \boldsymbol{\theta}_k + 0{,}995 \cdot \boldsymbol{\theta}_k^{-}$\;
            }
        }
    }
    $\varepsilon \leftarrow \max(0{,}02,\; \varepsilon \times 0{,}9995)$\;
    \BlankLine
    \If{épisode $\bmod$ 500 $= 0$}{
        Sauvegarder les checkpoints $\{\boldsymbol{\theta}_k\}$\;
    }
}
\end{algorithm}

Un détail d'implémentation important est la \textbf{fréquence d'entraînement réduite} : l'optimisation du réseau n'est effectuée que tous les 16 pas de temps ($t \bmod 16 = 0$), ce qui réduit le surcoût de transfert CPU$\leftrightarrow$GPU et évite l'overfitting sur des transitions fortement corrélées au sein d'un même épisode.

\subsubsection{Optimisation}

L'optimiseur utilisé est \textbf{Adam}~\cite{kingma2015} avec un taux d'apprentissage $\eta = 10^{-4}$ :
\begin{equation}
    \boldsymbol{\theta}_{t+1} = \boldsymbol{\theta}_t - \eta \cdot \frac{\hat{m}_t}{\sqrt{\hat{v}_t} + \epsilon_{\text{Adam}}}
    \label{eq:adam}
\end{equation}
où $\hat{m}_t$ et $\hat{v}_t$ sont les estimateurs corrigés du premier et second moment des gradients. Le taux d'apprentissage relativement faible ($10^{-4}$) est choisi pour assurer une convergence stable dans l'environnement multi-agent non-stationnaire.

% ============================================================
\subsection{Sauvegarde et évaluation des modèles}
\label{sec:checkpoints}

\subsubsection{Système de checkpoints}

Des points de contrôle (\textit{checkpoints}) sont sauvegardés tous les 500 épisodes, contenant :
\begin{itemize}
    \item Les paramètres du réseau de politique $\boldsymbol{\theta}_k$ et du réseau cible $\boldsymbol{\theta}_k^{-}$
    \item L'état de l'optimiseur Adam (moments)
    \item La valeur courante de $\varepsilon$
    \item Le compteur d'étapes d'entraînement
    \item L'historique des pertes
\end{itemize}

Cela permet de reprendre l'entraînement en cas d'interruption et d'évaluer les performances intermédiaires de la politique à différents stades de l'apprentissage.

\subsubsection{Protocole d'évaluation}

L'évaluation des modèles entraînés est réalisée en mode \textbf{glouton} ($\varepsilon = 0$) sur $N_{\text{eval}} = 50$ réalisations de canal indépendantes. Pour chaque réalisation, les agents sont exécutés pendant $T = 200$ pas afin de stabiliser les puissances d'émission avant de mesurer l'EE résultante.

\begin{figure*}[htbp]
    \centering
    \includegraphics[width=0.95\textwidth]{figures/training_curves.png}
    \caption{Courbes d'entraînement des agents DQN décentralisés sur $10\,000$ épisodes : (a) efficacité énergétique moyenne, (b) reward cumulatif, (c) perte d'entraînement, (d) SINR moyen, (e) puissance d'émission moyenne, (f) taux d'exploration $\varepsilon$.}
    \label{fig:training_curves}
\end{figure*}

% ============================================================
\subsection{Récapitulatif des hyperparamètres}
\label{sec:hyperparams_summary}

Le Tableau~\ref{tab:hyperparams} regroupe l'ensemble des hyperparamètres utilisés pour l'entraînement.

\begin{table}[H]
    \centering
    \caption{Hyperparamètres de l'entraînement DQN}
    \label{tab:hyperparams}
    \begin{tabular}{@{}llr@{}}
        \toprule
        \textbf{Catégorie} & \textbf{Paramètre} & \textbf{Valeur} \\
        \midrule
        \multirow{4}{*}{Architecture} 
            & Couches cachées & $[128, 256, 128]$ \\
            & Activation & ReLU \\
            & Initialisation & Xavier uniforme \\
            & Params/agent & 70\,944 \\
        \midrule
        \multirow{5}{*}{Optimisation}
            & Optimiseur & Adam \\
            & Taux d'apprentissage $\eta$ & $10^{-4}$ \\
            & Taille du batch $B$ & 256 \\
            & Fonction de perte & Huber ($\delta = 1{,}0$) \\
            & Gradient clipping & $G_{\max} = 1{,}0$ \\
        \midrule
        \multirow{3}{*}{Exploration}
            & $\varepsilon_0$ & 1,0 \\
            & $\varepsilon_{\text{fin}}$ & 0,02 \\
            & Décroissance $\lambda$ & 0,9995/épisode \\
        \midrule
        \multirow{4}{*}{Replay \& cible}
            & Taille du buffer $|\mathcal{D}|$ & 100\,000 \\
            & Buffer minimal $|\mathcal{D}|_{\min}$ & 2\,000 \\
            & Soft update $\tau$ & 0,005 \\
            & Fréquence entraîn. & 1/16 pas \\
        \midrule
        \multirow{3}{*}{Entraînement}
            & Facteur de discount $\gamma$ & 0,95 \\
            & Épisodes $N_{\text{ep}}$ & 10\,000 \\
            & Pas/épisode $T$ & 200 \\
        \midrule
        \multirow{3}{*}{Récompense}
            & $\alpha_{\text{ee}}$ & 1,0 \\
            & $\beta_{\text{qos}}$ & 5,0 \\
            & $\gamma_p$ & 0,02 \\
        \bottomrule
    \end{tabular}
\end{table}


% ============================================================
% ============================================================
%     CHAPITRE : RÉSULTATS ET COMPARAISONS
% ============================================================
% ============================================================

\section{Résultats Expérimentaux et Analyse Comparative}

\label{chap:results}

\section{Introduction}

Ce chapitre présente une analyse comparative détaillée entre l'approche DQN décentralisée développée dans le Chapitre~\ref{chap:dqn} et les méthodes de référence. Nous évaluons les performances sur $N = 50$ réalisations de canal indépendantes, garantissant une comparaison statistiquement significative. Toutes les méthodes sont testées sur les \textbf{mêmes} réalisations de canal (échantillons appariés), assurant l'équité de la comparaison.

\section{Méthodes comparées}
\label{sec:methods}

Cinq méthodes de contrôle de puissance sont évaluées :

\begin{enumerate}
    \item \textbf{Article (BRD + Dinkelbach)} : l'algorithme de référence de Zappone \textit{et al.}~\cite{zappone2015}, basé sur la dynamique de meilleure réponse avec résolution interne par l'algorithme de Dinkelbach (Chapitre~\ref{chap:article_algo}).

    \item \textbf{DQN (entraîné)} : notre approche par apprentissage par renforcement profond, avec $K = 5$ agents DQN indépendants entraînés sur $10\,000$ épisodes (Chapitre~\ref{chap:dqn}).

    \item \textbf{Optimal Uniform} : recherche exhaustive de la puissance uniforme $p^*$ optimale parmi 100 niveaux géométriquement espacés, où tous les utilisateurs émettent à la même puissance. Cette méthode sert de baseline analytique simple.

    \item \textbf{Max Power} : tous les utilisateurs émettent à la puissance maximale $p_{\max} = 23$~dBm ($200$~mW). Ce baseline naïf illustre la sous-optimalité de l'approche brute en termes d'EE.

    \item \textbf{Random Power} : chaque utilisateur choisit une puissance aléatoire uniforme dans $[p_{\min}, p_{\max}]$. Ce baseline représente l'absence totale de stratégie.
\end{enumerate}

\section{Résultats en efficacité énergétique}
\label{sec:ee_results}

\subsection{Performances globales}

Le Tableau~\ref{tab:ee_comparison} présente les statistiques descriptives de l'efficacité énergétique pour chaque méthode sur les 50 essais.

\begin{table*}[htbp]
    \centering
    \caption{Comparaison de l'efficacité énergétique (Mbit/J) — 50 essais}
    \label{tab:ee_comparison}
    \begin{tabular}{@{}lrrrrr@{}}
        \toprule
        \textbf{Méthode} & \textbf{Moyenne} & \textbf{Écart-type} & \textbf{Min} & \textbf{Max} & \textbf{Médiane} \\
        \midrule
        \textbf{DQN (entraîné)}         & \textbf{255,16} & 47,31 & 141,27 & 313,18 & --- \\
        Article (BRD+Dinkelbach)         & 207,01          & 30,67 & 137,99 & 256,16 & --- \\
        Optimal Uniform                  & 213,04          & 52,76 & ---    & ---    & --- \\
        Random Power                     & 23,06           & 15,59 & ---    & ---    & --- \\
        Max Power                        & 10,22           & 2,55  & ---    & ---    & --- \\
        \bottomrule
    \end{tabular}
\end{table*}

L'approche DQN atteint une EE moyenne de $\boldsymbol{255{,}16}$~\textbf{Mbit/J}, surpassant significativement l'algorithme de l'article ($207{,}01$~Mbit/J) avec un \textbf{gain relatif de $+23{,}3$~\%}. Les méthodes naïves (Max Power, Random) sont nettement inférieures, confirmant que le contrôle intelligent de puissance est essentiel pour l'efficacité énergétique.

\begin{figure}[htbp]
    \centering
    \includegraphics[width=0.9\textwidth]{figures/comparison_bar.png}
    \caption{Efficacité énergétique moyenne (Mbit/J) par méthode avec barres d'erreur ($\pm 1$ écart-type). Le DQN surpasse toutes les méthodes de référence.}
    \label{fig:ee_barplot}
\end{figure}

\subsection{Distribution de l'efficacité énergétique}

La Figure~\ref{fig:ee_distribution} présente la distribution de l'EE pour l'approche DQN et l'algorithme de l'article sur les 50 essais.

\begin{figure*}[htbp]
    \centering
    \subfloat[Histogramme comparatif\label{fig:ee_histogram}]{%
        \includegraphics[width=0.48\textwidth]{figures/ee_histogram_dqn_vs_article.png}}
    \hfill
    \subfloat[Fonction de répartition cumulative (CDF)\label{fig:ee_cdf}]{%
        \includegraphics[width=0.48\textwidth]{figures/comparison_cdf.png}}
    \caption{Distribution de l'efficacité énergétique : (a) histogrammes superposés montrant le décalage favorable du DQN ; (b) CDF confirmant la dominance stochastique du DQN sur l'algorithme de l'article.}
    \label{fig:ee_distribution}
\end{figure*}

La CDF (Figure~\ref{fig:ee_cdf}) montre que pour tout seuil d'EE, la probabilité que le DQN dépasse ce seuil est plus élevée que celle de l'algorithme de l'article, ce qui constitue une \textbf{dominance stochastique du premier ordre}. La médiane du DQN se situe aux alentours de 255~Mbit/J contre 218~Mbit/J pour l'article.

\subsection{Comparaison essai par essai}

La Figure~\ref{fig:ee_scatter_trial} présente un graphique de dispersion comparant l'EE du DQN à celle de l'article pour chaque essai individuellement.

\begin{figure}[htbp]
    \centering
    \includegraphics[width=0.65\textwidth]{figures/ee_scatter_dqn_vs_article.png}
    \caption{Comparaison essai par essai de l'EE : DQN (axe Y) vs Article (axe X). Les points au-dessus de la diagonale ($y = x$) indiquent une supériorité du DQN. La grande majorité des essais favorise le DQN.}
    \label{fig:ee_scatter_trial}
\end{figure}

% ============================================================
\section{Analyse de la consommation de puissance}
\label{sec:power_results}

\subsection{Puissance moyenne par méthode}

Le Tableau~\ref{tab:power_comparison} détaille la consommation de puissance moyenne par méthode.

\begin{table}[htbp]
    \centering
    \caption{Consommation de puissance moyenne (mW) par méthode}
    \label{tab:power_comparison}
    \begin{tabular}{@{}lr@{}}
        \toprule
        \textbf{Méthode} & \textbf{Puissance moyenne (mW)} \\
        \midrule
        \textbf{DQN (entraîné)}         & \textbf{15,27} \\
        Article (BRD+Dinkelbach)         & 28,73 \\
        Optimal Uniform                  & 0,02 \\
        Max Power                        & 199,53 \\
        Random Power                     & 100,57 \\
        \bottomrule
    \end{tabular}
\end{table}

L'approche DQN consomme en moyenne $\boldsymbol{15{,}27}$~\textbf{mW}, soit une \textbf{réduction de $46{,}8$~\%} par rapport à l'algorithme de l'article ($28{,}73$~mW). Cette économie de puissance significative est obtenue tout en améliorant l'EE, ce qui démontre que le DQN apprend à opérer dans la zone optimale du compromis débit/consommation.

\begin{figure}[htbp]
    \centering
    \includegraphics[width=0.9\textwidth]{figures/comparison_boxplot.png}
    \caption{Distribution de la puissance d'émission (mW) par méthode sous forme de boîtes à moustaches (échelle logarithmique). Le DQN opère à des puissances significativement inférieures à l'algorithme de l'article.}
    \label{fig:power_boxplot}
\end{figure}

\subsection{Compromis puissance--EE}

La Figure~\ref{fig:power_vs_ee} illustre le compromis fondamental entre la puissance d'émission et l'efficacité énergétique pour toutes les méthodes.

\begin{figure}[htbp]
    \centering
    \includegraphics[width=0.9\textwidth]{figures/scatter_power_vs_ee.png}
    \caption{Diagramme de dispersion puissance (mW, échelle log) vs EE (Mbit/J) pour les 50 essais. Chaque point représente un essai; la couleur indique la méthode. Le DQN (vert) se concentre dans la zone optimale : faible puissance, haute EE.}
    \label{fig:power_vs_ee}
\end{figure}

Ce graphique met en évidence le positionnement optimal du DQN dans le coin supérieur gauche (haute EE, faible puissance), confirmant que l'agent a appris la structure du compromis débit-puissance décrit dans l'article de Zappone \textit{et al.} : l'EE est une fonction quasi-concave de la puissance d'émission, avec un maximum bien défini à des puissances relativement basses.

% ============================================================
\section{Analyse de l'algorithme de l'article}
\label{sec:article_results}

L'algorithme de référence (BRD + Dinkelbach) présente les caractéristiques suivantes sur les 50 essais :

\subsection{Convergence}

\begin{itemize}
    \item \textbf{Fraction convergée} : $100\%$ des essais convergent (tolérance relative $10^{-4}$).
    \item \textbf{Nombre moyen d'itérations} : $5{,}28$ itérations de BRD par essai.
    \item \textbf{Nombre maximum d'itérations} : $12$ (convergence toujours rapide).
\end{itemize}

La convergence rapide et systématique confirme les résultats théoriques de l'article : sous les conditions de faisabilité, la dynamique de meilleure réponse converge vers l'unique GNE.

\subsection{Efficacité énergétique détaillée}

Les statistiques descriptives de l'EE obtenue par l'algorithme de l'article sont présentées dans le Tableau~\ref{tab:article_ee_stats}.

\begin{table}[htbp]
    \centering
    \caption{Statistiques détaillées de l'EE — Algorithme de l'article (50 essais)}
    \label{tab:article_ee_stats}
    \begin{tabular}{@{}lr@{}}
        \toprule
        \textbf{Statistique} & \textbf{Valeur (Mbit/J)} \\
        \midrule
        Moyenne              & 218,07 \\
        Écart-type           & 30,10 \\
        Minimum              & 137,00 \\
        Maximum              & 268,39 \\
        Médiane              & 218,08 \\
        Percentile 5\%       & 179,47 \\
        Percentile 95\%      & 264,72 \\
        \bottomrule
    \end{tabular}
\end{table}

\subsection{Débit spectral et QoS}

\begin{itemize}
    \item \textbf{Débit spectral moyen} : $\bar{R} = 2{,}951$~bits/s/Hz (écart-type~$= 0{,}347$).
    \item \textbf{Débit maximal théorique} : $R_{\max} = \log_2(1 + 125) = 6{,}977$~bits/s/Hz.
    \item \textbf{Contrainte QoS} : $\theta_k = 20\% \times R_{\max} = 1{,}395$~bits/s/Hz.
    \item \textbf{Taux de satisfaction QoS} : $\boldsymbol{74\%}$ des utilisateurs satisfont la contrainte de débit minimal.
\end{itemize}

Le taux de satisfaction QoS de 74\% indique que même l'algorithme optimal ne parvient pas toujours à satisfaire les contraintes de débit dans des conditions de canal défavorables, ce qui souligne la difficulté inhérente du problème de faisabilité dans les systèmes à interférence.

\begin{figure*}[htbp]
    \centering
    \subfloat[Distribution de l'EE\label{fig:article_ee_dist}]{%
        \includegraphics[width=0.48\textwidth]{figures/article_ee_distribution.png}}
    \hfill
    \subfloat[Distribution de la puissance\label{fig:article_power_dist}]{%
        \includegraphics[width=0.48\textwidth]{figures/article_power_distribution.png}}
    \caption{Distributions des métriques de l'algorithme de l'article sur 50 essais : (a) EE avec annotations des percentiles P5 et P95 ; (b) puissance émise montrant une forte variance entre les réalisations de canal.}
    \label{fig:article_distributions}
\end{figure*}

% ============================================================
\section{Analyse de l'entraînement DQN}
\label{sec:training_analysis}

\subsection{Convergence de l'efficacité énergétique}

La Figure~\ref{fig:training_ee} illustre l'évolution de l'EE moyenne au cours de l'entraînement. On distingue trois phases :

\begin{enumerate}
    \item \textbf{Phase d'exploration} (épisodes 1--2\,000) : l'EE est faible et volatile, les agents explorent l'espace d'actions de manière quasi-aléatoire.
    \item \textbf{Phase d'apprentissage rapide} (épisodes 2\,000--6\,000) : l'EE croît rapidement à mesure que les Q-valeurs convergent et que l'exploitation commence à dominer.
    \item \textbf{Phase de stabilisation} (épisodes 6\,000--10\,000) : l'EE se stabilise autour de 250--300~Mbit/J, les agents ayant convergé vers des politiques quasi-optimales.
\end{enumerate}

\begin{figure}[htbp]
    \centering
    \includegraphics[width=0.8\textwidth]{figures/training_ee_convergence.png}
    \caption{Convergence de l'efficacité énergétique moyenne durant l'entraînement. La courbe lissée (fenêtre de 100 épisodes) montre une convergence progressive vers $\sim 250$~Mbit/J. La ligne pointillée rouge indique le niveau atteint par l'algorithme de l'article ($\sim 218$~Mbit/J).}
    \label{fig:training_ee}
\end{figure}

\subsection{Évolution de la perte d'entraînement}

La Figure~\ref{fig:training_loss} montre l'évolution de la perte de Huber au cours de l'entraînement. Après un pic initial correspondant aux premières mises à jour des Q-valeurs (lorsque le réseau cible est encore initialisé aléatoirement), la perte décroît progressivement, indiquant une réduction de l'erreur de prédiction temporelle. La stabilisation de la perte aux alentours de l'épisode 5\,000 confirme que les Q-valeurs convergent.

\begin{figure}[htbp]
    \centering
    \includegraphics[width=0.75\textwidth]{figures/training_loss.png}
    \caption{Perte d'entraînement (Huber loss) lissée sur une fenêtre de 100 épisodes. La décroissance progressive confirme la convergence des Q-valeurs.}
    \label{fig:training_loss}
\end{figure}

\subsection{Évolution de la puissance d'émission}

L'analyse de la puissance d'émission moyenne au cours de l'entraînement (Figure~\ref{fig:training_power}) révèle un comportement caractéristique : les agents apprennent progressivement à réduire leur puissance d'émission, passant d'une puissance élevée (exploration aléatoire) à des puissances optimales basses, cohérentes avec la maximisation de l'EE.

\begin{figure}[htbp]
    \centering
    \includegraphics[width=0.75\textwidth]{figures/training_power_evolution.png}
    \caption{Évolution de la puissance d'émission moyenne (dBm) au cours de l'entraînement. Les agents apprennent à réduire progressivement leur puissance vers le point de fonctionnement optimal.}
    \label{fig:training_power}
\end{figure}

\subsection{Évolution du SINR}

La Figure~\ref{fig:training_sinr} présente l'évolution du SINR moyen au cours de l'entraînement. On observe une stabilisation progressive du SINR à mesure que les agents convergent vers des puissances cohérentes, réduisant les fluctuations d'interférence mutuelle.

\begin{figure}[htbp]
    \centering
    \includegraphics[width=0.75\textwidth]{figures/training_sinr_evolution.png}
    \caption{Évolution du SINR moyen (dB) durant l'entraînement. La stabilisation du SINR reflète l'émergence d'un équilibre entre les agents.}
    \label{fig:training_sinr}
\end{figure}

% ============================================================
\section{Comparaison multi-critères}
\label{sec:multi_criteria}

Le Tableau~\ref{tab:multi_criteria} présente une comparaison synthétique sur l'ensemble des critères de performance.

\begin{table*}[htbp]
    \centering
    \caption{Comparaison multi-critères des méthodes de contrôle de puissance}
    \label{tab:multi_criteria}
    \renewcommand{\arraystretch}{1.2}
    \begin{tabular}{@{}lccccc@{}}
        \toprule
        \textbf{Critère} & \textbf{DQN} & \textbf{Article} & \textbf{Opt. Unif.} & \textbf{Max Power} & \textbf{Random} \\
        \midrule
        EE (Mbit/J) $\uparrow$      & \textbf{255,16} & 207,01 & 213,04 & 10,22 & 23,06 \\
        Puissance (mW) $\downarrow$  & \textbf{15,27}  & 28,73  & 0,02   & 199,53 & 100,57 \\
        Stabilité ($1/\sigma$) $\uparrow$ & 0,021 & \textbf{0,033} & 0,019 & 0,392 & 0,064 \\
        Connaissance modèle          & \ding{55}       & \ding{51} & \ding{55} & \ding{55} & \ding{55} \\
        Décentralisé                  & \ding{51}       & \ding{51} & \ding{55} & N/A   & N/A \\
        Adaptabilité                  & \textbf{Haute}  & Basse    & Nulle & Nulle & Nulle \\
        \bottomrule
    \end{tabular}
\end{table*}

\begin{figure}[htbp]
    \centering
    \includegraphics[width=0.7\textwidth]{figures/radar_chart_comparison.png}
    \caption{Diagramme radar comparatif des méthodes sur quatre critères normalisés : EE, économie de puissance (~$1/P$), stabilité ($1/\sigma$), et débit spectral.}
    \label{fig:radar_comparison}
\end{figure}

% ============================================================
\section{Discussion}
\label{sec:discussion}

\subsection{Avantages du DQN}

\begin{enumerate}
    \item \textbf{Performances supérieures} : le DQN surpasse l'algorithme analytique de $+23{,}3\%$ en EE tout en consommant $46{,}8\%$ moins de puissance. Cet avantage est remarquable car il est obtenu \textit{sans connaissance explicite du modèle de canal}.

    \item \textbf{Absence de modèle} (\textit{model-free}) : contrairement à l'approche BRD + Dinkelbach qui nécessite la connaissance exacte de la structure du SINR et la résolution analytique de l'optimisation fractionnaire, le DQN apprend directement à partir des observations et des récompenses. Cela le rend robuste aux inexactitudes de modélisation et adaptable à des scénarios non couverts par la théorie.

    \item \textbf{Adaptabilité} : les agents DQN peuvent s'adapter à des changements dans l'environnement (nombre d'utilisateurs, topologie, conditions de propagation) par ré-entraînement ou \textit{fine-tuning}, sans nécessiter de re-dérivation des solutions analytiques.
\end{enumerate}

\subsection{Limitations et travaux futurs}

\begin{enumerate}
    \item \textbf{Variance plus élevée} : l'écart-type de l'EE du DQN ($47{,}31$~Mbit/J) est supérieur à celui de l'algorithme de l'article ($30{,}67$~Mbit/J), indiquant une plus grande variabilité des performances selon les réalisations de canal. Cela pourrait être atténué par des techniques de régularisation ou un entraînement plus long.

    \item \textbf{Coût d'entraînement} : l'entraînement de $10\,000$ épisodes requiert un temps significatif, notamment avec $K = 5$ agents indépendants. Cependant, ce coût est assumé hors-ligne (\textit{offline}) et la politique résultante est exécutée en temps réel avec une simple passe avant dans le réseau de neurones.

    \item \textbf{Garanties théoriques} : contrairement à l'algorithme de l'article qui fournit des garanties de convergence vers le GNE, l'approche IQL n'a pas de preuve formelle de convergence dans le cas multi-agent non-stationnaire. Néanmoins, les résultats empiriques démontrent une convergence effective.

    \item \textbf{Passage à l'échelle} : des travaux futurs pourraient explorer des architectures de \textit{parameter sharing} entre agents ou des approches de type \textit{Mean Field} pour étendre le système à un plus grand nombre d'utilisateurs.
\end{enumerate}

% ============================================================
\section{Conclusion du chapitre}

Les résultats expérimentaux démontrent que l'approche par Deep Q-Network constitue une alternative viable et performante à l'algorithme analytique de Zappone \textit{et al.} pour le contrôle de puissance décentralisé dans les réseaux 5G Massive MIMO. Avec un gain de $+23{,}3\%$ en efficacité énergétique et une réduction de $46{,}8\%$ en consommation de puissance, le DQN valide le potentiel de l'apprentissage par renforcement profond pour l'optimisation des ressources radio. L'avantage principal réside dans sa nature \textit{model-free}, qui le rend applicable à des scénarios réalistes où le modèle mathématique exact du canal n'est pas disponible ou est trop complexe pour être exploité analytiquement.

% ============================================================
% BIBLIOGRAPHIE (extraits pertinents)
% ============================================================
% À inclure dans le fichier .bib principal :
%
% @article{zappone2015,
%   author  = {Zappone, A. and Björnson, E. and Sanguinetti, L. and Jorswieck, E.},
%   title   = {Energy-Efficient Power Control: A Look at 5G Wireless Technologies},
%   journal = {IEEE Transactions on Signal Processing},
%   year    = {2015},
%   volume  = {64},
%   number  = {7},
%   pages   = {1668--1683}
% }
%
% @article{mnih2015,
%   author  = {Mnih, V. and Kavukcuoglu, K. and Silver, D. and others},
%   title   = {Human-level control through deep reinforcement learning},
%   journal = {Nature},
%   year    = {2015},
%   volume  = {518},
%   pages   = {529--533}
% }
%
% @inproceedings{vanHasselt2016,
%   author  = {van Hasselt, H. and Guez, A. and Silver, D.},
%   title   = {Deep Reinforcement Learning with Double Q-learning},
%   booktitle = {AAAI},
%   year    = {2016}
% }
%
% @article{kingma2015,
%   author  = {Kingma, D. P. and Ba, J.},
%   title   = {Adam: A Method for Stochastic Optimization},
%   journal = {ICLR},
%   year    = {2015}
% }
%
% @inproceedings{glorot2010,
%   author  = {Glorot, X. and Bengio, Y.},
%   title   = {Understanding the difficulty of training deep feedforward neural networks},
%   booktitle = {AISTATS},
%   year    = {2010}
% }
%
% @inproceedings{tan1993,
%   author  = {Tan, M.},
%   title   = {Multi-Agent Reinforcement Learning: Independent vs. Cooperative Agents},
%   booktitle = {ICML},
%   year    = {1993}
% }
%
% @article{lanctot2017,
%   author  = {Lanctot, M. and Zambaldi, V. and Gruslys, A. and others},
%   title   = {A Unified Game-Theoretic Approach to Multiagent Reinforcement Learning},
%   journal = {NeurIPS},
%   year    = {2017}
% }


%==============================================================================
\section{Conclusion et perspectives}
%==============================================================================

[À développer]

\begin{thebibliography}{99}

\bibitem{3GPP} 3GPP, "Technical Specification Group Radio Access Network; NR; Physical layer; General description (Release 15)," 3GPP TS 38.201, 2018.

\bibitem{Fehske2011} A. Fehske, J. Malmodin, G. Biczok, and G. Fettweis, "The global footprint of mobile communications: The ecological and economic perspective," \textit{IEEE Commun. Mag.}, pp. 55-62, Aug. 2011.

\bibitem{Ericsson2011} Ericsson White Paper, "More than 50 billion connected devices," Ericsson, Tech. Rep. 284 23-3149 Uen, Feb. 2011.

\bibitem{Qualcomm} Qualcomm, "The 1000x data challenge," Tech. Rep.

\bibitem{Zappone2021} A. Zappone, L. Sanguinetti, G. Bacci, E. Jorswieck, and M. Debbah, "Energy-Efficient Power Control: A Look at 5G Wireless Technologies," \textit{IEEE Transactions on Signal Processing}, 2021.

\bibitem{Dinkelbach1967} W. Dinkelbach, "On nonlinear fractional programming," \textit{Management Science}, vol. 13, no. 7, pp. 492-498, Mar. 1967.

\bibitem{Ng2012} D. W. K. Ng, E. S. Lo, and R. Schober, "Energy-efficient resource allocation in multi-cell OFDMA systems with limited backhaul capacity," \textit{IEEE Trans. Wireless Commun.}, vol. 11, no. 10, pp. 3618-3631, Oct. 2012.

\bibitem{Xu2014} Q. Xu, X. Li, H. Ji, and X. Du, "Energy-efficient resource allocation for heterogeneous services in OFDMA downlink networks: Systematic perspective," \textit{IEEE Trans. Veh. Technol.}, vol. 63, no. 5, pp. 2071-2082, June 2014.

\bibitem{Xu2013} J. Xu and L. Qiu, "Energy efficiency optimization for MIMO broadcast channels," \textit{IEEE Trans. Wireless Commun.}, vol. 12, no. 2, pp. 690-701, Feb. 2013.

\bibitem{Du2014} B. Du, C. Pan, W. Zhang, and M. Chen, "Distributed energy-efficient power optimization for CoMP systems with max-min fairness," \textit{IEEE Commun. Lett.}, vol. 18, no. 6, pp. 999-1002, Jun. 2014.

\bibitem{He2013} S. He, Y. Huang, S. Jin, and L. Yang, "Coordinated beamforming for energy efficient transmission in multicell multiuser systems," \textit{IEEE Trans. Commun.}, vol. 61, no. 12, pp. 4961-4971, Dec. 2013.

\bibitem{He2014} S. He, Y. Huang, L. Yang, and B. Ottersten, "Coordinated multicell multiuser precoding for maximizing weighted sum energy efficiency," \textit{IEEE Trans. Signal Process.}, vol. 62, no. 3, pp. 741-751, Feb. 2014.

\bibitem{Venturino2015} L. Venturino, A. Zappone, C. Risi, and S. Buzzi, "Energy-efficient scheduling and power allocation in downlink OFDMA networks with base station coordination," \textit{IEEE Trans. Wireless Commun.}, vol. 14, no. 1, pp. 1-14, 2015.

\bibitem{Miao2011} G. Miao, N. Himayat, G. Li, and S. Talwar, "Distributed interference-aware energy-efficient power optimization," \textit{IEEE Trans. Wireless Commun.}, vol. 10, no. 4, pp. 1323-1333, Apr. 2011.

\bibitem{Scutari2008} G. Scutari, D. Palomar, and S. Barbarossa, "Optimal linear precoding strategies for wideband noncooperative systems based on game theory part I: Nash equilibria," \textit{IEEE Trans. Signal Process.}, vol. 56, no. 3, pp. 1230-1249, March 2008.

\bibitem{Zappone2013} A. Zappone, Z. Chong, E. Jorswieck, and S. Buzzi, "Energy-aware competitive power control in relay-assisted interference wireless networks," \textit{IEEE Trans. Wireless Commun.}, vol. 12, no. 4, pp. 1860-1871, Apr. 2013.

\bibitem{Stupia2015} I. Stupia, L. Sanguinetti, G. Bacci, and L. Vandendorpe, "Power control in networks with heterogeneous users: A quasi-variational inequality approach," \textit{IEEE Trans. Signal Process.}, vol. 63, no. 21, pp. 5691-5705, November 2015.

\bibitem{Bacci2015} G. Bacci, E. Belmega, P. Mertikopoulos, and L. Sanguinetti, "Energy-aware competitive power allocation for heterogeneous networks under QoS constraints," \textit{IEEE Trans. Wireless Commun.}, vol. 14, no. 9, pp. 4728-4742, Sept 2015.

\bibitem{Yates1995} R. D. Yates, "A framework for uplink power control in cellular radio systems," \textit{IEEE J. Sel. Areas Commun.}, vol. 13, no. 7, pp. 1341-1347, September 1995.

\end{thebibliography}

\end{document}
% Prérequis packages : amsmath, amssymb, algorithm2e, graphicx,
%   booktabs, multirow, tabularx, subcaption, hyperref, xcolor

% ============================================================
%            CHAPITRE : APPROCHE PAR DEEP Q-NETWORK
% ============================================================
\section{Contrôle de Puissance par Deep Q-Network}
\label{chap:dqn}

\subsection{Introduction}

Le chapitre précédent a présenté l'approche analytique de Zappone \textit{et al.}~\cite{zappone2015} pour la maximisation décentralisée de l'efficacité énergétique (EE) dans les réseaux 5G Massive MIMO, reposant sur la dynamique de meilleure réponse (BRD) couplée à l'algorithme de Dinkelbach. Bien que cette méthode converge vers l'Équilibre de Nash Généralisé (GNE) sous certaines conditions, elle nécessite une connaissance explicite du modèle de canal et la résolution itérative d'un problème d'optimisation fractionnaire à chaque mise à jour.

Dans ce chapitre, nous proposons une approche alternative basée sur l'apprentissage par renforcement profond (\textit{Deep Reinforcement Learning}), et plus spécifiquement sur le \textbf{Deep Q-Network (DQN)}~\cite{mnih2015}. L'objectif est de permettre à chaque utilisateur $k$ d'apprendre de manière autonome une politique de contrôle de puissance qui maximise son EE individuelle, sans connaissance explicite du modèle mathématique du canal. L'approche est intrinsèquement décentralisée : chaque utilisateur dispose de son propre agent DQN qui prend des décisions indépendantes, tout en interagissant avec les autres agents via l'interférence partagée dans l'environnement sans fil.

% ============================================================
\subsection{Formulation en Processus Décisionnel de Markov}
\label{sec:mdp}

Pour appliquer l'apprentissage par renforcement, nous reformulons le problème de contrôle de puissance comme un \textbf{Processus Décisionnel de Markov} (MDP). Chaque utilisateur $k \in \{1, \dots, K\}$ est modélisé comme un agent autonome interagissant avec un environnement partagé. Le MDP de l'utilisateur $k$ est défini par le tuple $(\mathcal{S}_k, \mathcal{A}_k, \mathcal{P}, \mathcal{R}_k, \gamma)$, dont les composantes sont détaillées ci-après.

\subsubsection{Espace d'états}
\label{sec:state_space}

L'état observé par l'utilisateur $k$ à l'instant $t$ est un vecteur de dimension 6 :
\begin{equation}
    \mathbf{s}_k^{(t)} = \begin{bmatrix}
        s_1^{(t)} \\ s_2^{(t)} \\ s_3^{(t)} \\ s_4^{(t)} \\ s_5^{(t)} \\ s_6^{(t)}
    \end{bmatrix}
    \label{eq:state_vector}
\end{equation}
dont les composantes sont définies comme suit :

\begin{enumerate}
    \item \textbf{SINR normalisé en dB} :
    \begin{equation}
        s_1^{(t)} = \text{clip}\!\left(\frac{\text{SINR}_k^{\text{dB}}}{25}, -1, 1\right), \quad \text{SINR}_k^{\text{dB}} = 10\log_{10}(\gamma_k^{(t)})
        \label{eq:state_sinr}
    \end{equation}
    où la normalisation par 25~dB correspond à la plage pratique du SINR dans le système considéré.

    \item \textbf{Puissance normalisée (échelle logarithmique)} :
    \begin{equation}
        s_2^{(t)} = \frac{\log_{10}(p_k^{(t)}) - \log_{10}(p_{\min})}{\log_{10}(p_{\max}) - \log_{10}(p_{\min})}
        \label{eq:state_power}
    \end{equation}
    Cette représentation en échelle log permet une résolution uniforme sur toute la plage de puissance.

    \item \textbf{Interférence normalisée} :
    \begin{equation}
        s_3^{(t)} = \frac{\log_{10}\left(\sum_{j \neq k} p_j^{(t)}\right) - \log_{10}\left((K\!-\!1)\,p_{\min}\right)}{\log_{10}\left((K\!-\!1)\,p_{\max}\right) - \log_{10}\left((K\!-\!1)\,p_{\min}\right)}
        \label{eq:state_interf}
    \end{equation}

    \item \textbf{Gain de canal équivalent normalisé} :
    \begin{equation}
        s_4^{(t)} = \text{clip}\!\left(\frac{\log_{10}(\mu_k^{(t)})}{10}, 0, 1\right)
        \label{eq:state_mu}
    \end{equation}
    où $\mu_k = \frac{\alpha_k}{\sigma^2 + \sum_{j \neq k} \omega_{kj}\,p_j}$ est le gain de canal équivalent utilisé dans la meilleure réponse de Dinkelbach. Cette composante est cruciale car elle encode l'information nécessaire pour déterminer la puissance optimale.

    \item \textbf{Débit normalisé} :
    \begin{equation}
        s_5^{(t)} = \frac{R_k^{(t)}}{R_{\max}}, \quad R_{\max} = \log_2(1 + \bar{\gamma})
        \label{eq:state_rate}
    \end{equation}
    où $R_{\max}$ est le débit maximal théorique limité par les imperfections matérielles ($\bar{\gamma} = 1/\varepsilon^2$).

    \item \textbf{Marge QoS} :
    \begin{equation}
        s_6^{(t)} = \text{clip}\!\left(\frac{R_k^{(t)} - \theta_k}{\theta_k}, -1, 1\right)
        \label{eq:state_qos}
    \end{equation}
    où $\theta_k$ est le débit minimal requis par la contrainte de qualité de service. Une valeur positive indique que la contrainte est satisfaite.
\end{enumerate}

Toutes les composantes sont normalisées dans l'intervalle $[-1, 1]$ ou $[0, 1]$ afin de faciliter l'apprentissage du réseau de neurones. Le choix d'une représentation en échelle logarithmique pour la puissance et l'interférence reflète la nature multiplicative des signaux radio et améliore la sensibilité de l'agent aux variations dans les régimes de faible puissance, où se situe typiquement le point de fonctionnement optimal en termes d'EE.

\subsubsection{Espace d'actions}
\label{sec:action_space}

L'espace d'actions est \textbf{discret} et composé de $|\mathcal{A}| = 32$ niveaux de puissance répartis de manière logarithmique entre $p_{\min}$ et $p_{\max}$ :
\begin{equation}
    \mathcal{A} = \left\{ p_{\min} \cdot \left(\frac{p_{\max}}{p_{\min}}\right)^{\frac{i}{N_a - 1}} \;\middle|\; i = 0, 1, \dots, N_a - 1 \right\}
    \label{eq:action_space}
\end{equation}
avec $p_{\min} = -20$~dBm ($0{,}01$~mW) et $p_{\max} = 23$~dBm ($200$~mW), soit un rapport de 1:20\,000.

L'espacement logarithmique (géométrique) est motivé par le fait que l'efficacité énergétique optimale se situe typiquement dans la partie inférieure de la plage de puissance. Cette discrétisation offre ainsi une \textbf{résolution fine} dans la zone optimale (entre 0,01 et 5~mW) tout en couvrant l'ensemble de la plage autorisée. L'action $a_k^{(t)} \in \{0, 1, \dots, 31\}$ sélectionne directement le niveau de puissance d'émission $p_k^{(t)} = \mathcal{A}[a_k^{(t)}]$.

\subsubsection{Fonction de récompense}
\label{sec:reward}

La fonction de récompense est conçue pour guider l'agent vers la maximisation de l'efficacité énergétique tout en respectant les contraintes de QoS et en encourageant l'économie de puissance :

\begin{equation}
    r_k^{(t)} = \text{clip}\!\left(\alpha_{\text{ee}} \cdot \frac{\eta_k^{(t)}}{\eta_{\text{ref}}} - \beta_{\text{qos}} \cdot \mathbb{1}_{\text{QoS}} \cdot \left(0{,}3 + 0{,}7 \cdot \min\!\left(\frac{\theta_k - R_k^{(t)}}{\theta_k}, 1\right)\right) - \gamma_p \cdot \frac{p_k^{(t)}}{p_{\max}},\; -3,\; 2\right)
    \label{eq:reward}
\end{equation}

où :
\begin{itemize}
    \item $\eta_k^{(t)} = \frac{B \cdot \log_2(1 + \gamma_k^{(t)})}{p_c + p_k^{(t)}}$ est l'efficacité énergétique en bit/Joule ;
    \item $\eta_{\text{ref}} = 3 \times 10^8$ bit/J (= 300~Mbit/J) est la référence de normalisation, choisie pour que le terme EE soit dans l'intervalle $[0, 1]$ au voisinage de la solution optimale ;
    \item $\alpha_{\text{ee}} = 1{,}0$ est le coefficient de récompense EE ;
    \item $\beta_{\text{qos}} = 5{,}0$ est le coefficient de pénalité QoS, volontairement élevé pour fortement décourager les violations de la contrainte de débit minimal ;
    \item $\gamma_p = 0{,}02$ est le coefficient de pénalité de puissance, encourageant l'économie d'énergie ;
    \item $\mathbb{1}_{\text{QoS}} = \mathbb{1}_{R_k < \theta_k}$ est l'indicatrice de violation QoS.
\end{itemize}

La pénalité QoS est composée d'un terme fixe ($0{,}3$) activé dès qu'il y a violation, et d'un terme proportionnel au déficit relatif $(\theta_k - R_k)/\theta_k$, ce qui permet de différencier les violations légères des violations sévères. Le \textit{soft clipping} dans $[-3, 2]$ assure la stabilité numérique de l'entraînement en évitant les récompenses extrêmes.

\subsubsection{Dynamique de l'environnement}
\label{sec:env_dynamics}

L'environnement simule un réseau 5G Massive MIMO avec les caractéristiques suivantes :

\begin{itemize}
    \item \textbf{Canal dynamique} : à chaque pas de temps, le canal est mis à jour par un évanouissement rapide résiduel :
    \begin{equation}
        d_k^{\text{eff}}(t) = d_k \cdot (1 + \delta_k(t)), \quad \delta_k \sim \mathcal{N}\!\left(0, \frac{1}{\sqrt{M}}\right)
        \label{eq:channel_hardening}
    \end{equation}
    Ce modèle capture l'effet de \textit{channel hardening} typique du Massive MIMO avec $M = 50$ antennes, donnant des variations de $\pm 1$--$2$~dB.

    \item \textbf{Réinitialisation par épisode} : à chaque épisode, les $K = 5$ utilisateurs sont repositionnés aléatoirement dans la cellule (rayon $R_{\text{cell}} = 400$~m, distance minimale $d_{\min} = 35$~m), et un nouveau réalisation de l'affaiblissement grande échelle (perte de trajet + \textit{shadowing}) est tirée :
    \begin{equation}
        d_k = 10^{-(L_k + S_k)/10}, \quad L_k = 10\alpha\log_{10}(r_k), \quad S_k \sim \mathcal{N}(0, \sigma_s^2)
        \label{eq:large_scale_fading}
    \end{equation}
    avec $\alpha = 3{,}76$ et $\sigma_s = 8$~dB.

    \item \textbf{Durée d'épisode} : chaque épisode comprend $T = 200$ pas de temps, durant lesquels les agents interagissent simultanément.
\end{itemize}

% ============================================================
\subsection{Architecture du Deep Q-Network}
\label{sec:dqn_architecture}

\subsubsection{Principe du DQN}

Le Deep Q-Network~\cite{mnih2015} approxime la fonction de valeur action-état $Q^*(s, a)$ par un réseau de neurones profond $Q(s, a; \boldsymbol{\theta})$ paramétré par $\boldsymbol{\theta}$. La politique gloutonne (\textit{greedy}) correspondante est :
\begin{equation}
    \pi^*(s) = \arg\max_{a \in \mathcal{A}} Q(s, a; \boldsymbol{\theta})
    \label{eq:greedy_policy}
\end{equation}

L'objectif est de minimiser l'erreur temporelle différentielle (\textit{TD error}) entre la valeur Q courante et la valeur cible :
\begin{equation}
    \mathcal{L}(\boldsymbol{\theta}) = \mathbb{E}_{(s,a,r,s') \sim \mathcal{D}} \left[ \ell\!\left(r + \gamma \max_{a'} Q(s', a'; \boldsymbol{\theta}^{-}) - Q(s, a; \boldsymbol{\theta})\right) \right]
    \label{eq:dqn_loss}
\end{equation}
où $\mathcal{D}$ est le \textit{replay buffer}, $\boldsymbol{\theta}^{-}$ les paramètres du réseau cible, $\gamma = 0{,}95$ le facteur d'actualisation, et $\ell(\cdot)$ la fonction de perte.

\subsubsection{Architecture du réseau de neurones}

Le réseau de neurones utilisé pour chaque agent DQN est un perceptron multicouche (\textit{Multi-Layer Perceptron}, MLP) composé de trois couches cachées entièrement connectées (\textit{fully connected}). L'architecture est la suivante :

\begin{equation}
    \underbrace{\mathbb{R}^6}_{\text{Entrée (état)}} \xrightarrow{\mathbf{W}_1, \mathbf{b}_1} \underbrace{\mathbb{R}^{128}}_{\text{Couche 1}} \xrightarrow{\text{ReLU}} \underbrace{\mathbb{R}^{256}}_{\text{Couche 2}} \xrightarrow{\text{ReLU}} \underbrace{\mathbb{R}^{128}}_{\text{Couche 3}} \xrightarrow{\text{ReLU}} \underbrace{\mathbb{R}^{32}}_{\text{Sortie (Q-values)}}
    \label{eq:architecture}
\end{equation}

\begin{itemize}
    \item \textbf{Couche d'entrée} : 6 neurones correspondant aux 6 composantes de l'état $\mathbf{s}_k$.
    \item \textbf{Couche cachée 1} : 128 neurones avec activation ReLU.
    \item \textbf{Couche cachée 2} : 256 neurones avec activation ReLU.
    \item \textbf{Couche cachée 3} : 128 neurones avec activation ReLU.
    \item \textbf{Couche de sortie} : 32 neurones (un par action), sans activation — les sorties représentent directement les Q-valeurs $Q(s, a_i; \boldsymbol{\theta})$ pour $i = 0, \dots, 31$.
\end{itemize}

L'architecture en forme de \textit{bottleneck} (128 $\to$ 256 $\to$ 128) permet d'abord une expansion de la représentation dans l'espace latent, puis une compression vers la dimension de sortie, favorisant l'extraction de caractéristiques pertinentes.

Les poids sont initialisés selon la méthode de \textbf{Xavier uniforme}~\cite{glorot2010} :
\begin{equation}
    W_{ij} \sim \mathcal{U}\left(-\sqrt{\frac{6}{n_{\text{in}} + n_{\text{out}}}},\; +\sqrt{\frac{6}{n_{\text{in}} + n_{\text{out}}}}\right)
    \label{eq:xavier_init}
\end{equation}
et les biais sont initialisés à zéro. Le nombre total de paramètres entraînables par agent est de :
\begin{equation}
    N_{\text{params}} = (6 \times 128 + 128) + (128 \times 256 + 256) + (256 \times 128 + 128) + (128 \times 32 + 32) = 70\,944
    \label{eq:n_params}
\end{equation}

\begin{figure*}[htbp]
    \centering
    \includegraphics[width=0.85\textwidth]{figures/dqn_architecture.png}
    \caption{Architecture du réseau DQN pour le contrôle de puissance. L'entrée est le vecteur d'état à 6 dimensions et la sortie fournit les Q-valeurs pour les 32 niveaux de puissance discrets.}
    \label{fig:dqn_architecture}
\end{figure*}

% ============================================================
\subsection{Mécanismes d'apprentissage}
\label{sec:learning_mechanisms}

\subsubsection{Double DQN}
\label{sec:double_dqn}

Afin de réduire la surestimation des Q-valeurs inhérente au DQN standard~\cite{vanHasselt2016}, nous implémentons la variante \textbf{Double DQN}. Le principe consiste à découpler la \textbf{sélection} de l'action (effectuée par le réseau de politique $\boldsymbol{\theta}$) de son \textbf{évaluation} (effectuée par le réseau cible $\boldsymbol{\theta}^{-}$) :

\begin{equation}
    y^{\text{DDQN}} = r + \gamma \cdot Q\!\left(s',\; \underbrace{\arg\max_{a'} Q(s', a'; \boldsymbol{\theta})}_{\text{sélection par réseau politique}};\; \boldsymbol{\theta}^{-}\right)
    \label{eq:double_dqn_target}
\end{equation}

Cette technique réduit significativement le biais d'estimation positif et améliore la stabilité de l'apprentissage, ce qui est particulièrement important dans notre contexte multi-agent où les environnements sont non-stationnaires du point de vue de chaque agent.

\subsubsection{Replay d'expérience}
\label{sec:experience_replay}

Chaque agent maintient un \textbf{tampon de rejeu} (\textit{experience replay buffer}) $\mathcal{D}_k$ de capacité $|\mathcal{D}| = 100\,000$ transitions. Chaque transition stockée est un tuple :
\begin{equation}
    e^{(t)} = \left(s_k^{(t)},\; a_k^{(t)},\; r_k^{(t)},\; s_k^{(t+1)},\; d^{(t)}\right)
    \label{eq:transition}
\end{equation}
où $d^{(t)} \in \{0, 1\}$ est l'indicateur de fin d'épisode. Le tampon fonctionne selon le principe FIFO (\textit{First-In, First-Out}) : lorsque la capacité maximale est atteinte, les transitions les plus anciennes sont remplacées.

À chaque pas d'entraînement, un \textbf{mini-batch} de $B = 256$ transitions est échantillonné uniformément de manière aléatoire :
\begin{equation}
    \{(s_i, a_i, r_i, s'_i, d_i)\}_{i=1}^{B} \sim \text{Uniforme}(\mathcal{D}_k)
    \label{eq:sampling}
\end{equation}

Le replay d'expérience confère deux avantages fondamentaux :
\begin{itemize}
    \item \textbf{Décorrélation temporelle} : les transitions échantillonnées aléatoirement brisent les corrélations séquentielles, satisfaisant l'hypothèse i.i.d. de la descente de gradient stochastique.
    \item \textbf{Efficacité en données} : chaque transition peut être réutilisée dans plusieurs mises à jour, améliorant l'exploitation des données collectées.
\end{itemize}

L'entraînement ne commence qu'une fois le tampon contenant au moins $|\mathcal{D}|_{\min} = 2\,000$ transitions, garantissant une diversité suffisante dans les premiers mini-batchs.

\subsubsection{Réseau cible et mise à jour douce}
\label{sec:target_network}

Pour stabiliser l'apprentissage, nous utilisons un \textbf{réseau cible} (\textit{target network}) $Q(s, a; \boldsymbol{\theta}^{-})$ distinct du réseau de politique $Q(s, a; \boldsymbol{\theta})$. Les paramètres du réseau cible sont mis à jour par \textbf{mise à jour douce} (\textit{soft update} ou \textit{Polyak averaging}) :

\begin{equation}
    \boldsymbol{\theta}^{-} \leftarrow \tau \cdot \boldsymbol{\theta} + (1 - \tau) \cdot \boldsymbol{\theta}^{-}, \quad \tau = 0{,}005
    \label{eq:soft_update}
\end{equation}

Cette mise à jour est effectuée après chaque pas d'optimisation. Le coefficient $\tau = 0{,}005$ assure une transition lente et régulière du réseau cible vers le réseau de politique, évitant les oscillations et divergences liées à une cible en mouvement rapide. Par rapport à la mise à jour dure (\textit{hard update}) périodique, la mise à jour douce offre une stabilité accrue car le réseau cible évolue de manière continue plutôt que par sauts discontinus.

\subsubsection{Fonction de perte de Huber}
\label{sec:huber_loss}

La fonction de perte utilisée pour l'optimisation est la \textbf{perte de Huber} (aussi appelée \textit{Smooth L1 Loss}) :
\begin{equation}
    \ell_{\delta}(x) = \begin{cases}
        \frac{1}{2}x^2 & \text{si } |x| \leq \delta \\
        \delta\left(|x| - \frac{1}{2}\delta\right) & \text{sinon}
    \end{cases}
    \label{eq:huber_loss}
\end{equation}
avec $\delta = 1{,}0$. La perte de Huber combine les avantages de la MSE (sensibilité aux petites erreurs) et de la MAE (robustesse aux valeurs aberrantes), ce qui est crucial dans notre environnement multi-agent où les récompenses peuvent être bruitées par les actions concurrentes des autres agents.

\subsubsection{Politique d'exploration $\varepsilon$-gloutonne}
\label{sec:epsilon_greedy}

L'exploration du space d'actions est réalisée par une politique $\varepsilon$-gloutonne (\textit{$\varepsilon$-greedy}) :
\begin{equation}
    a_k^{(t)} = \begin{cases}
        \text{Uniforme}(\mathcal{A}) & \text{avec probabilité } \varepsilon_t \\
        \arg\max_{a} Q(s_k^{(t)}, a; \boldsymbol{\theta}_k) & \text{avec probabilité } 1 - \varepsilon_t
    \end{cases}
    \label{eq:epsilon_greedy}
\end{equation}

Le taux d'exploration $\varepsilon_t$ décroît de manière exponentielle au fil des épisodes :
\begin{equation}
    \varepsilon_t = \max\!\left(\varepsilon_{\text{fin}},\; \varepsilon_0 \cdot \lambda^t\right), \quad \varepsilon_0 = 1{,}0, \quad \varepsilon_{\text{fin}} = 0{,}02, \quad \lambda = 0{,}9995
    \label{eq:epsilon_decay}
\end{equation}

Ce schéma de décroissance assure :
\begin{itemize}
    \item Une \textbf{phase d'exploration intensive} initiale (épisodes 1--1\,000, $\varepsilon > 0{,}6$) permettant une couverture large de l'espace état-action.
    \item Une \textbf{transition progressive} (épisodes 1\,000--7\,000) où l'agent bascule graduellement vers l'exploitation.
    \item Une \textbf{phase d'exploitation} finale (épisodes 7\,000--10\,000, $\varepsilon \approx 0{,}02$) où l'agent utilise quasi-exclusivement la politique apprise, tout en conservant une exploration résiduelle de 2\% pour maintenir une certaine adaptabilité.
\end{itemize}

L'épisode où $\varepsilon$ passe sous le seuil $0{,}1$ (exploitation dominante) se situe aux alentours de l'épisode $4\,600$, ce qui correspond au régime où les Q-valeurs sont suffisamment stables pour guider les décisions.

\begin{figure}[htbp]
    \centering
    \includegraphics[width=1\textwidth]{figures/epsilon_decay.png}
    \caption{Décroissance exponentielle du taux d'exploration $\varepsilon$ au cours des $10\,000$ épisodes d'entraînement. La ligne pointillée indique le seuil $\varepsilon = 0{,}1$.}
    \label{fig:epsilon_decay}
\end{figure}

\subsubsection{Gradient clipping}
\label{sec:grad_clip}

Pour prévenir l'instabilité numérique causée par des gradients de grande magnitude, une \textbf{troncature du gradient} (\textit{gradient clipping}) est appliquée avant chaque mise à jour des poids :
\begin{equation}
    \nabla_{\boldsymbol{\theta}} \mathcal{L} \leftarrow \frac{\nabla_{\boldsymbol{\theta}} \mathcal{L}}{\max\!\left(1, \frac{\|\nabla_{\boldsymbol{\theta}} \mathcal{L}\|_2}{G_{\max}}\right)}, \quad G_{\max} = 1{,}0
    \label{eq:grad_clip}
\end{equation}

Cette technique est d'autant plus importante dans le contexte multi-agent, où la non-stationnarité induite par l'apprentissage simultané de tous les agents peut engendrer des gradients instables.

% ============================================================
\subsection{Entraînement multi-agent décentralisé}
\label{sec:multi_agent}

\subsubsection{Independent Q-Learning (IQL)}

L'approche adoptée est le \textbf{Independent Q-Learning} (IQL)~\cite{tan1993}, où chaque utilisateur $k$ possède son propre agent DQN qui apprend indépendamment. Les agents interagissent uniquement à travers l'environnement partagé : le SINR de l'utilisateur $k$ dépend des puissances d'émission de tous les utilisateurs via l'interférence :
\begin{equation}
    \gamma_k = \frac{\alpha_k \, p_k}{\sigma^2 + \phi_k \, p_k + \sum_{j \neq k} \omega_{kj} \, p_j}
    \label{eq:sinr_coupling}
\end{equation}

Du point de vue de l'utilisateur $k$, les autres utilisateurs font partie de l'environnement, qui devient \textbf{non-stationnaire} puisque les politiques des autres agents évoluent au cours de l'entraînement. Bien que cela viole formellement les hypothèses de stationnarité du MDP, l'IQL fonctionne remarquablement bien en pratique pour les jeux coopératifs et compétitifs de taille modérée~\cite{lanctot2017}.

\subsubsection{Protocole d'entraînement}

Le protocole d'entraînement est résumé dans l'Algorithme~\ref{alg:training}.

\begin{algorithm}[htbp]
\caption{Entraînement multi-agent DQN décentralisé}
\label{alg:training}
\SetKwInOut{Input}{Entrées}
\SetKwInOut{Output}{Sorties}
\Input{$K = 5$ agents, $N_{\text{ep}} = 10\,000$ épisodes, $T = 200$ pas/épisode}
\Output{Politiques optimales $\{\boldsymbol{\theta}_k^*\}_{k=1}^{K}$}
\BlankLine
Initialiser $K$ réseaux de politique $\{Q(\cdot; \boldsymbol{\theta}_k)\}$ et cibles $\{Q(\cdot; \boldsymbol{\theta}_k^{-})\}$\;
Initialiser $K$ tampons de rejeu $\{\mathcal{D}_k\}$ de capacité $100\,000$\;
$\varepsilon \leftarrow 1{,}0$\;
\BlankLine
\For{épisode $= 1$ \KwTo $N_{\text{ep}}$}{
    Réinitialiser l'environnement (nouvelles positions + canal)\;
    Observer les états initiaux $\{s_k^{(0)}\}_{k=1}^{K}$\;
    \BlankLine
    \For{$t = 0$ \KwTo $T-1$}{
        \For{$k = 1$ \KwTo $K$ \textrm{(en parallèle)}}{
            Sélectionner $a_k^{(t)}$ par politique $\varepsilon$-gloutonne (Éq.~\ref{eq:epsilon_greedy})\;
        }
        Exécuter les actions $\{a_k^{(t)}\}$ dans l'environnement\;
        Observer récompenses $\{r_k^{(t)}\}$ et nouveaux états $\{s_k^{(t+1)}\}$\;
        \BlankLine
        \For{$k = 1$ \KwTo $K$}{
            Stocker $(s_k^{(t)}, a_k^{(t)}, r_k^{(t)}, s_k^{(t+1)}, d^{(t)})$ dans $\mathcal{D}_k$\;
            \If{$|\mathcal{D}_k| \geq 2\,000$ \textbf{et} $t \bmod 16 = 0$}{
                Échantillonner un mini-batch de 256 transitions de $\mathcal{D}_k$\;
                Calculer les cibles Double DQN (Éq.~\ref{eq:double_dqn_target})\;
                Mettre à jour $\boldsymbol{\theta}_k$ par descente de gradient (perte de Huber)\;
                Mise à jour douce : $\boldsymbol{\theta}_k^{-} \leftarrow 0{,}005 \cdot \boldsymbol{\theta}_k + 0{,}995 \cdot \boldsymbol{\theta}_k^{-}$\;
            }
        }
    }
    $\varepsilon \leftarrow \max(0{,}02,\; \varepsilon \times 0{,}9995)$\;
    \BlankLine
    \If{épisode $\bmod$ 500 $= 0$}{
        Sauvegarder les checkpoints $\{\boldsymbol{\theta}_k\}$\;
    }
}
\end{algorithm}

Un détail d'implémentation important est la \textbf{fréquence d'entraînement réduite} : l'optimisation du réseau n'est effectuée que tous les 16 pas de temps ($t \bmod 16 = 0$), ce qui réduit le surcoût de transfert CPU$\leftrightarrow$GPU et évite l'overfitting sur des transitions fortement corrélées au sein d'un même épisode.

\subsubsection{Optimisation}

L'optimiseur utilisé est \textbf{Adam}~\cite{kingma2015} avec un taux d'apprentissage $\eta = 10^{-4}$ :
\begin{equation}
    \boldsymbol{\theta}_{t+1} = \boldsymbol{\theta}_t - \eta \cdot \frac{\hat{m}_t}{\sqrt{\hat{v}_t} + \epsilon_{\text{Adam}}}
    \label{eq:adam}
\end{equation}
où $\hat{m}_t$ et $\hat{v}_t$ sont les estimateurs corrigés du premier et second moment des gradients. Le taux d'apprentissage relativement faible ($10^{-4}$) est choisi pour assurer une convergence stable dans l'environnement multi-agent non-stationnaire.

% ============================================================
\subsection{Sauvegarde et évaluation des modèles}
\label{sec:checkpoints}

\subsubsection{Système de checkpoints}

Des points de contrôle (\textit{checkpoints}) sont sauvegardés tous les 500 épisodes, contenant :
\begin{itemize}
    \item Les paramètres du réseau de politique $\boldsymbol{\theta}_k$ et du réseau cible $\boldsymbol{\theta}_k^{-}$
    \item L'état de l'optimiseur Adam (moments)
    \item La valeur courante de $\varepsilon$
    \item Le compteur d'étapes d'entraînement
    \item L'historique des pertes
\end{itemize}

Cela permet de reprendre l'entraînement en cas d'interruption et d'évaluer les performances intermédiaires de la politique à différents stades de l'apprentissage.

\subsubsection{Protocole d'évaluation}

L'évaluation des modèles entraînés est réalisée en mode \textbf{glouton} ($\varepsilon = 0$) sur $N_{\text{eval}} = 50$ réalisations de canal indépendantes. Pour chaque réalisation, les agents sont exécutés pendant $T = 200$ pas afin de stabiliser les puissances d'émission avant de mesurer l'EE résultante.

\begin{figure*}[htbp]
    \centering
    \includegraphics[width=0.95\textwidth]{figures/training_curves.png}
    \caption{Courbes d'entraînement des agents DQN décentralisés sur $10\,000$ épisodes : (a) efficacité énergétique moyenne, (b) reward cumulatif, (c) perte d'entraînement, (d) SINR moyen, (e) puissance d'émission moyenne, (f) taux d'exploration $\varepsilon$.}
    \label{fig:training_curves}
\end{figure*}

% ============================================================
\subsection{Récapitulatif des hyperparamètres}
\label{sec:hyperparams_summary}

Le Tableau~\ref{tab:hyperparams} regroupe l'ensemble des hyperparamètres utilisés pour l'entraînement.

\begin{table}[H]
    \centering
    \caption{Hyperparamètres de l'entraînement DQN}
    \label{tab:hyperparams}
    \begin{tabular}{@{}llr@{}}
        \toprule
        \textbf{Catégorie} & \textbf{Paramètre} & \textbf{Valeur} \\
        \midrule
        \multirow{4}{*}{Architecture} 
            & Couches cachées & $[128, 256, 128]$ \\
            & Activation & ReLU \\
            & Initialisation & Xavier uniforme \\
            & Params/agent & 70\,944 \\
        \midrule
        \multirow{5}{*}{Optimisation}
            & Optimiseur & Adam \\
            & Taux d'apprentissage $\eta$ & $10^{-4}$ \\
            & Taille du batch $B$ & 256 \\
            & Fonction de perte & Huber ($\delta = 1{,}0$) \\
            & Gradient clipping & $G_{\max} = 1{,}0$ \\
        \midrule
        \multirow{3}{*}{Exploration}
            & $\varepsilon_0$ & 1,0 \\
            & $\varepsilon_{\text{fin}}$ & 0,02 \\
            & Décroissance $\lambda$ & 0,9995/épisode \\
        \midrule
        \multirow{4}{*}{Replay \& cible}
            & Taille du buffer $|\mathcal{D}|$ & 100\,000 \\
            & Buffer minimal $|\mathcal{D}|_{\min}$ & 2\,000 \\
            & Soft update $\tau$ & 0,005 \\
            & Fréquence entraîn. & 1/16 pas \\
        \midrule
        \multirow{3}{*}{Entraînement}
            & Facteur de discount $\gamma$ & 0,95 \\
            & Épisodes $N_{\text{ep}}$ & 10\,000 \\
            & Pas/épisode $T$ & 200 \\
        \midrule
        \multirow{3}{*}{Récompense}
            & $\alpha_{\text{ee}}$ & 1,0 \\
            & $\beta_{\text{qos}}$ & 5,0 \\
            & $\gamma_p$ & 0,02 \\
        \bottomrule
    \end{tabular}
\end{table}


% ============================================================
% ============================================================
%     CHAPITRE : RÉSULTATS ET COMPARAISONS
% ============================================================
% ============================================================

\section{Résultats Expérimentaux et Analyse Comparative}

\label{chap:results}

\section{Introduction}

Ce chapitre présente une analyse comparative détaillée entre l'approche DQN décentralisée développée dans le Chapitre~\ref{chap:dqn} et les méthodes de référence. Nous évaluons les performances sur $N = 50$ réalisations de canal indépendantes, garantissant une comparaison statistiquement significative. Toutes les méthodes sont testées sur les \textbf{mêmes} réalisations de canal (échantillons appariés), assurant l'équité de la comparaison.

\section{Méthodes comparées}
\label{sec:methods}

Cinq méthodes de contrôle de puissance sont évaluées :

\begin{enumerate}
    \item \textbf{Article (BRD + Dinkelbach)} : l'algorithme de référence de Zappone \textit{et al.}~\cite{zappone2015}, basé sur la dynamique de meilleure réponse avec résolution interne par l'algorithme de Dinkelbach (Chapitre~\ref{chap:article_algo}).

    \item \textbf{DQN (entraîné)} : notre approche par apprentissage par renforcement profond, avec $K = 5$ agents DQN indépendants entraînés sur $10\,000$ épisodes (Chapitre~\ref{chap:dqn}).

    \item \textbf{Optimal Uniform} : recherche exhaustive de la puissance uniforme $p^*$ optimale parmi 100 niveaux géométriquement espacés, où tous les utilisateurs émettent à la même puissance. Cette méthode sert de baseline analytique simple.

    \item \textbf{Max Power} : tous les utilisateurs émettent à la puissance maximale $p_{\max} = 23$~dBm ($200$~mW). Ce baseline naïf illustre la sous-optimalité de l'approche brute en termes d'EE.

    \item \textbf{Random Power} : chaque utilisateur choisit une puissance aléatoire uniforme dans $[p_{\min}, p_{\max}]$. Ce baseline représente l'absence totale de stratégie.
\end{enumerate}

\section{Résultats en efficacité énergétique}
\label{sec:ee_results}

\subsection{Performances globales}

Le Tableau~\ref{tab:ee_comparison} présente les statistiques descriptives de l'efficacité énergétique pour chaque méthode sur les 50 essais.

\begin{table*}[htbp]
    \centering
    \caption{Comparaison de l'efficacité énergétique (Mbit/J) — 50 essais}
    \label{tab:ee_comparison}
    \begin{tabular}{@{}lrrrrr@{}}
        \toprule
        \textbf{Méthode} & \textbf{Moyenne} & \textbf{Écart-type} & \textbf{Min} & \textbf{Max} & \textbf{Médiane} \\
        \midrule
        \textbf{DQN (entraîné)}         & \textbf{255,16} & 47,31 & 141,27 & 313,18 & --- \\
        Article (BRD+Dinkelbach)         & 207,01          & 30,67 & 137,99 & 256,16 & --- \\
        Optimal Uniform                  & 213,04          & 52,76 & ---    & ---    & --- \\
        Random Power                     & 23,06           & 15,59 & ---    & ---    & --- \\
        Max Power                        & 10,22           & 2,55  & ---    & ---    & --- \\
        \bottomrule
    \end{tabular}
\end{table*}

L'approche DQN atteint une EE moyenne de $\boldsymbol{255{,}16}$~\textbf{Mbit/J}, surpassant significativement l'algorithme de l'article ($207{,}01$~Mbit/J) avec un \textbf{gain relatif de $+23{,}3$~\%}. Les méthodes naïves (Max Power, Random) sont nettement inférieures, confirmant que le contrôle intelligent de puissance est essentiel pour l'efficacité énergétique.

\begin{figure}[htbp]
    \centering
    \includegraphics[width=0.9\textwidth]{figures/comparison_bar.png}
    \caption{Efficacité énergétique moyenne (Mbit/J) par méthode avec barres d'erreur ($\pm 1$ écart-type). Le DQN surpasse toutes les méthodes de référence.}
    \label{fig:ee_barplot}
\end{figure}

\subsection{Distribution de l'efficacité énergétique}

La Figure~\ref{fig:ee_distribution} présente la distribution de l'EE pour l'approche DQN et l'algorithme de l'article sur les 50 essais.

\begin{figure*}[htbp]
    \centering
    \subfloat[Histogramme comparatif\label{fig:ee_histogram}]{%
        \includegraphics[width=0.48\textwidth]{figures/ee_histogram_dqn_vs_article.png}}
    \hfill
    \subfloat[Fonction de répartition cumulative (CDF)\label{fig:ee_cdf}]{%
        \includegraphics[width=0.48\textwidth]{figures/comparison_cdf.png}}
    \caption{Distribution de l'efficacité énergétique : (a) histogrammes superposés montrant le décalage favorable du DQN ; (b) CDF confirmant la dominance stochastique du DQN sur l'algorithme de l'article.}
    \label{fig:ee_distribution}
\end{figure*}

La CDF (Figure~\ref{fig:ee_cdf}) montre que pour tout seuil d'EE, la probabilité que le DQN dépasse ce seuil est plus élevée que celle de l'algorithme de l'article, ce qui constitue une \textbf{dominance stochastique du premier ordre}. La médiane du DQN se situe aux alentours de 255~Mbit/J contre 218~Mbit/J pour l'article.

\subsection{Comparaison essai par essai}

La Figure~\ref{fig:ee_scatter_trial} présente un graphique de dispersion comparant l'EE du DQN à celle de l'article pour chaque essai individuellement.

\begin{figure}[htbp]
    \centering
    \includegraphics[width=0.65\textwidth]{figures/ee_scatter_dqn_vs_article.png}
    \caption{Comparaison essai par essai de l'EE : DQN (axe Y) vs Article (axe X). Les points au-dessus de la diagonale ($y = x$) indiquent une supériorité du DQN. La grande majorité des essais favorise le DQN.}
    \label{fig:ee_scatter_trial}
\end{figure}

% ============================================================
\section{Analyse de la consommation de puissance}
\label{sec:power_results}

\subsection{Puissance moyenne par méthode}

Le Tableau~\ref{tab:power_comparison} détaille la consommation de puissance moyenne par méthode.

\begin{table}[htbp]
    \centering
    \caption{Consommation de puissance moyenne (mW) par méthode}
    \label{tab:power_comparison}
    \begin{tabular}{@{}lr@{}}
        \toprule
        \textbf{Méthode} & \textbf{Puissance moyenne (mW)} \\
        \midrule
        \textbf{DQN (entraîné)}         & \textbf{15,27} \\
        Article (BRD+Dinkelbach)         & 28,73 \\
        Optimal Uniform                  & 0,02 \\
        Max Power                        & 199,53 \\
        Random Power                     & 100,57 \\
        \bottomrule
    \end{tabular}
\end{table}

L'approche DQN consomme en moyenne $\boldsymbol{15{,}27}$~\textbf{mW}, soit une \textbf{réduction de $46{,}8$~\%} par rapport à l'algorithme de l'article ($28{,}73$~mW). Cette économie de puissance significative est obtenue tout en améliorant l'EE, ce qui démontre que le DQN apprend à opérer dans la zone optimale du compromis débit/consommation.

\begin{figure}[htbp]
    \centering
    \includegraphics[width=0.9\textwidth]{figures/comparison_boxplot.png}
    \caption{Distribution de la puissance d'émission (mW) par méthode sous forme de boîtes à moustaches (échelle logarithmique). Le DQN opère à des puissances significativement inférieures à l'algorithme de l'article.}
    \label{fig:power_boxplot}
\end{figure}

\subsection{Compromis puissance--EE}

La Figure~\ref{fig:power_vs_ee} illustre le compromis fondamental entre la puissance d'émission et l'efficacité énergétique pour toutes les méthodes.

\begin{figure}[htbp]
    \centering
    \includegraphics[width=0.9\textwidth]{figures/scatter_power_vs_ee.png}
    \caption{Diagramme de dispersion puissance (mW, échelle log) vs EE (Mbit/J) pour les 50 essais. Chaque point représente un essai; la couleur indique la méthode. Le DQN (vert) se concentre dans la zone optimale : faible puissance, haute EE.}
    \label{fig:power_vs_ee}
\end{figure}

Ce graphique met en évidence le positionnement optimal du DQN dans le coin supérieur gauche (haute EE, faible puissance), confirmant que l'agent a appris la structure du compromis débit-puissance décrit dans l'article de Zappone \textit{et al.} : l'EE est une fonction quasi-concave de la puissance d'émission, avec un maximum bien défini à des puissances relativement basses.

% ============================================================
\section{Analyse de l'algorithme de l'article}
\label{sec:article_results}

L'algorithme de référence (BRD + Dinkelbach) présente les caractéristiques suivantes sur les 50 essais :

\subsection{Convergence}

\begin{itemize}
    \item \textbf{Fraction convergée} : $100\%$ des essais convergent (tolérance relative $10^{-4}$).
    \item \textbf{Nombre moyen d'itérations} : $5{,}28$ itérations de BRD par essai.
    \item \textbf{Nombre maximum d'itérations} : $12$ (convergence toujours rapide).
\end{itemize}

La convergence rapide et systématique confirme les résultats théoriques de l'article : sous les conditions de faisabilité, la dynamique de meilleure réponse converge vers l'unique GNE.

\subsection{Efficacité énergétique détaillée}

Les statistiques descriptives de l'EE obtenue par l'algorithme de l'article sont présentées dans le Tableau~\ref{tab:article_ee_stats}.

\begin{table}[htbp]
    \centering
    \caption{Statistiques détaillées de l'EE — Algorithme de l'article (50 essais)}
    \label{tab:article_ee_stats}
    \begin{tabular}{@{}lr@{}}
        \toprule
        \textbf{Statistique} & \textbf{Valeur (Mbit/J)} \\
        \midrule
        Moyenne              & 218,07 \\
        Écart-type           & 30,10 \\
        Minimum              & 137,00 \\
        Maximum              & 268,39 \\
        Médiane              & 218,08 \\
        Percentile 5\%       & 179,47 \\
        Percentile 95\%      & 264,72 \\
        \bottomrule
    \end{tabular}
\end{table}

\subsection{Débit spectral et QoS}

\begin{itemize}
    \item \textbf{Débit spectral moyen} : $\bar{R} = 2{,}951$~bits/s/Hz (écart-type~$= 0{,}347$).
    \item \textbf{Débit maximal théorique} : $R_{\max} = \log_2(1 + 125) = 6{,}977$~bits/s/Hz.
    \item \textbf{Contrainte QoS} : $\theta_k = 20\% \times R_{\max} = 1{,}395$~bits/s/Hz.
    \item \textbf{Taux de satisfaction QoS} : $\boldsymbol{74\%}$ des utilisateurs satisfont la contrainte de débit minimal.
\end{itemize}

Le taux de satisfaction QoS de 74\% indique que même l'algorithme optimal ne parvient pas toujours à satisfaire les contraintes de débit dans des conditions de canal défavorables, ce qui souligne la difficulté inhérente du problème de faisabilité dans les systèmes à interférence.

\begin{figure*}[htbp]
    \centering
    \subfloat[Distribution de l'EE\label{fig:article_ee_dist}]{%
        \includegraphics[width=0.48\textwidth]{figures/article_ee_distribution.png}}
    \hfill
    \subfloat[Distribution de la puissance\label{fig:article_power_dist}]{%
        \includegraphics[width=0.48\textwidth]{figures/article_power_distribution.png}}
    \caption{Distributions des métriques de l'algorithme de l'article sur 50 essais : (a) EE avec annotations des percentiles P5 et P95 ; (b) puissance émise montrant une forte variance entre les réalisations de canal.}
    \label{fig:article_distributions}
\end{figure*}

% ============================================================
\section{Analyse de l'entraînement DQN}
\label{sec:training_analysis}

\subsection{Convergence de l'efficacité énergétique}

La Figure~\ref{fig:training_ee} illustre l'évolution de l'EE moyenne au cours de l'entraînement. On distingue trois phases :

\begin{enumerate}
    \item \textbf{Phase d'exploration} (épisodes 1--2\,000) : l'EE est faible et volatile, les agents explorent l'espace d'actions de manière quasi-aléatoire.
    \item \textbf{Phase d'apprentissage rapide} (épisodes 2\,000--6\,000) : l'EE croît rapidement à mesure que les Q-valeurs convergent et que l'exploitation commence à dominer.
    \item \textbf{Phase de stabilisation} (épisodes 6\,000--10\,000) : l'EE se stabilise autour de 250--300~Mbit/J, les agents ayant convergé vers des politiques quasi-optimales.
\end{enumerate}

\begin{figure}[htbp]
    \centering
    \includegraphics[width=0.8\textwidth]{figures/training_ee_convergence.png}
    \caption{Convergence de l'efficacité énergétique moyenne durant l'entraînement. La courbe lissée (fenêtre de 100 épisodes) montre une convergence progressive vers $\sim 250$~Mbit/J. La ligne pointillée rouge indique le niveau atteint par l'algorithme de l'article ($\sim 218$~Mbit/J).}
    \label{fig:training_ee}
\end{figure}

\subsection{Évolution de la perte d'entraînement}

La Figure~\ref{fig:training_loss} montre l'évolution de la perte de Huber au cours de l'entraînement. Après un pic initial correspondant aux premières mises à jour des Q-valeurs (lorsque le réseau cible est encore initialisé aléatoirement), la perte décroît progressivement, indiquant une réduction de l'erreur de prédiction temporelle. La stabilisation de la perte aux alentours de l'épisode 5\,000 confirme que les Q-valeurs convergent.

\begin{figure}[htbp]
    \centering
    \includegraphics[width=0.75\textwidth]{figures/training_loss.png}
    \caption{Perte d'entraînement (Huber loss) lissée sur une fenêtre de 100 épisodes. La décroissance progressive confirme la convergence des Q-valeurs.}
    \label{fig:training_loss}
\end{figure}

\subsection{Évolution de la puissance d'émission}

L'analyse de la puissance d'émission moyenne au cours de l'entraînement (Figure~\ref{fig:training_power}) révèle un comportement caractéristique : les agents apprennent progressivement à réduire leur puissance d'émission, passant d'une puissance élevée (exploration aléatoire) à des puissances optimales basses, cohérentes avec la maximisation de l'EE.

\begin{figure}[htbp]
    \centering
    \includegraphics[width=0.75\textwidth]{figures/training_power_evolution.png}
    \caption{Évolution de la puissance d'émission moyenne (dBm) au cours de l'entraînement. Les agents apprennent à réduire progressivement leur puissance vers le point de fonctionnement optimal.}
    \label{fig:training_power}
\end{figure}

\subsection{Évolution du SINR}

La Figure~\ref{fig:training_sinr} présente l'évolution du SINR moyen au cours de l'entraînement. On observe une stabilisation progressive du SINR à mesure que les agents convergent vers des puissances cohérentes, réduisant les fluctuations d'interférence mutuelle.

\begin{figure}[htbp]
    \centering
    \includegraphics[width=0.75\textwidth]{figures/training_sinr_evolution.png}
    \caption{Évolution du SINR moyen (dB) durant l'entraînement. La stabilisation du SINR reflète l'émergence d'un équilibre entre les agents.}
    \label{fig:training_sinr}
\end{figure}

% ============================================================
\section{Comparaison multi-critères}
\label{sec:multi_criteria}

Le Tableau~\ref{tab:multi_criteria} présente une comparaison synthétique sur l'ensemble des critères de performance.

\begin{table*}[htbp]
    \centering
    \caption{Comparaison multi-critères des méthodes de contrôle de puissance}
    \label{tab:multi_criteria}
    \renewcommand{\arraystretch}{1.2}
    \begin{tabular}{@{}lccccc@{}}
        \toprule
        \textbf{Critère} & \textbf{DQN} & \textbf{Article} & \textbf{Opt. Unif.} & \textbf{Max Power} & \textbf{Random} \\
        \midrule
        EE (Mbit/J) $\uparrow$      & \textbf{255,16} & 207,01 & 213,04 & 10,22 & 23,06 \\
        Puissance (mW) $\downarrow$  & \textbf{15,27}  & 28,73  & 0,02   & 199,53 & 100,57 \\
        Stabilité ($1/\sigma$) $\uparrow$ & 0,021 & \textbf{0,033} & 0,019 & 0,392 & 0,064 \\
        Connaissance modèle          & \ding{55}       & \ding{51} & \ding{55} & \ding{55} & \ding{55} \\
        Décentralisé                  & \ding{51}       & \ding{51} & \ding{55} & N/A   & N/A \\
        Adaptabilité                  & \textbf{Haute}  & Basse    & Nulle & Nulle & Nulle \\
        \bottomrule
    \end{tabular}
\end{table*}

\begin{figure}[htbp]
    \centering
    \includegraphics[width=0.7\textwidth]{figures/radar_chart_comparison.png}
    \caption{Diagramme radar comparatif des méthodes sur quatre critères normalisés : EE, économie de puissance (~$1/P$), stabilité ($1/\sigma$), et débit spectral.}
    \label{fig:radar_comparison}
\end{figure}

% ============================================================
\section{Discussion}
\label{sec:discussion}

\subsection{Avantages du DQN}

\begin{enumerate}
    \item \textbf{Performances supérieures} : le DQN surpasse l'algorithme analytique de $+23{,}3\%$ en EE tout en consommant $46{,}8\%$ moins de puissance. Cet avantage est remarquable car il est obtenu \textit{sans connaissance explicite du modèle de canal}.

    \item \textbf{Absence de modèle} (\textit{model-free}) : contrairement à l'approche BRD + Dinkelbach qui nécessite la connaissance exacte de la structure du SINR et la résolution analytique de l'optimisation fractionnaire, le DQN apprend directement à partir des observations et des récompenses. Cela le rend robuste aux inexactitudes de modélisation et adaptable à des scénarios non couverts par la théorie.

    \item \textbf{Adaptabilité} : les agents DQN peuvent s'adapter à des changements dans l'environnement (nombre d'utilisateurs, topologie, conditions de propagation) par ré-entraînement ou \textit{fine-tuning}, sans nécessiter de re-dérivation des solutions analytiques.
\end{enumerate}

\subsection{Limitations et travaux futurs}

\begin{enumerate}
    \item \textbf{Variance plus élevée} : l'écart-type de l'EE du DQN ($47{,}31$~Mbit/J) est supérieur à celui de l'algorithme de l'article ($30{,}67$~Mbit/J), indiquant une plus grande variabilité des performances selon les réalisations de canal. Cela pourrait être atténué par des techniques de régularisation ou un entraînement plus long.

    \item \textbf{Coût d'entraînement} : l'entraînement de $10\,000$ épisodes requiert un temps significatif, notamment avec $K = 5$ agents indépendants. Cependant, ce coût est assumé hors-ligne (\textit{offline}) et la politique résultante est exécutée en temps réel avec une simple passe avant dans le réseau de neurones.

    \item \textbf{Garanties théoriques} : contrairement à l'algorithme de l'article qui fournit des garanties de convergence vers le GNE, l'approche IQL n'a pas de preuve formelle de convergence dans le cas multi-agent non-stationnaire. Néanmoins, les résultats empiriques démontrent une convergence effective.

    \item \textbf{Passage à l'échelle} : des travaux futurs pourraient explorer des architectures de \textit{parameter sharing} entre agents ou des approches de type \textit{Mean Field} pour étendre le système à un plus grand nombre d'utilisateurs.
\end{enumerate}

% ============================================================
\section{Conclusion du chapitre}

Les résultats expérimentaux démontrent que l'approche par Deep Q-Network constitue une alternative viable et performante à l'algorithme analytique de Zappone \textit{et al.} pour le contrôle de puissance décentralisé dans les réseaux 5G Massive MIMO. Avec un gain de $+23{,}3\%$ en efficacité énergétique et une réduction de $46{,}8\%$ en consommation de puissance, le DQN valide le potentiel de l'apprentissage par renforcement profond pour l'optimisation des ressources radio. L'avantage principal réside dans sa nature \textit{model-free}, qui le rend applicable à des scénarios réalistes où le modèle mathématique exact du canal n'est pas disponible ou est trop complexe pour être exploité analytiquement.

% ============================================================
% BIBLIOGRAPHIE (extraits pertinents)
% ============================================================
% À inclure dans le fichier .bib principal :
%
% @article{zappone2015,
%   author  = {Zappone, A. and Björnson, E. and Sanguinetti, L. and Jorswieck, E.},
%   title   = {Energy-Efficient Power Control: A Look at 5G Wireless Technologies},
%   journal = {IEEE Transactions on Signal Processing},
%   year    = {2015},
%   volume  = {64},
%   number  = {7},
%   pages   = {1668--1683}
% }
%
% @article{mnih2015,
%   author  = {Mnih, V. and Kavukcuoglu, K. and Silver, D. and others},
%   title   = {Human-level control through deep reinforcement learning},
%   journal = {Nature},
%   year    = {2015},
%   volume  = {518},
%   pages   = {529--533}
% }
%
% @inproceedings{vanHasselt2016,
%   author  = {van Hasselt, H. and Guez, A. and Silver, D.},
%   title   = {Deep Reinforcement Learning with Double Q-learning},
%   booktitle = {AAAI},
%   year    = {2016}
% }
%
% @article{kingma2015,
%   author  = {Kingma, D. P. and Ba, J.},
%   title   = {Adam: A Method for Stochastic Optimization},
%   journal = {ICLR},
%   year    = {2015}
% }
%
% @inproceedings{glorot2010,
%   author  = {Glorot, X. and Bengio, Y.},
%   title   = {Understanding the difficulty of training deep feedforward neural networks},
%   booktitle = {AISTATS},
%   year    = {2010}
% }
%
% @inproceedings{tan1993,
%   author  = {Tan, M.},
%   title   = {Multi-Agent Reinforcement Learning: Independent vs. Cooperative Agents},
%   booktitle = {ICML},
%   year    = {1993}
% }
%
% @article{lanctot2017,
%   author  = {Lanctot, M. and Zambaldi, V. and Gruslys, A. and others},
%   title   = {A Unified Game-Theoretic Approach to Multiagent Reinforcement Learning},
%   journal = {NeurIPS},
%   year    = {2017}
% }


%==============================================================================
\section{Conclusion et perspectives}
%==============================================================================

[À développer]

\begin{thebibliography}{99}

\bibitem{3GPP} 3GPP, "Technical Specification Group Radio Access Network; NR; Physical layer; General description (Release 15)," 3GPP TS 38.201, 2018.

\bibitem{Fehske2011} A. Fehske, J. Malmodin, G. Biczok, and G. Fettweis, "The global footprint of mobile communications: The ecological and economic perspective," \textit{IEEE Commun. Mag.}, pp. 55-62, Aug. 2011.

\bibitem{Ericsson2011} Ericsson White Paper, "More than 50 billion connected devices," Ericsson, Tech. Rep. 284 23-3149 Uen, Feb. 2011.

\bibitem{Qualcomm} Qualcomm, "The 1000x data challenge," Tech. Rep.

\bibitem{Zappone2021} A. Zappone, L. Sanguinetti, G. Bacci, E. Jorswieck, and M. Debbah, "Energy-Efficient Power Control: A Look at 5G Wireless Technologies," \textit{IEEE Transactions on Signal Processing}, 2021.

\bibitem{Dinkelbach1967} W. Dinkelbach, "On nonlinear fractional programming," \textit{Management Science}, vol. 13, no. 7, pp. 492-498, Mar. 1967.

\bibitem{Ng2012} D. W. K. Ng, E. S. Lo, and R. Schober, "Energy-efficient resource allocation in multi-cell OFDMA systems with limited backhaul capacity," \textit{IEEE Trans. Wireless Commun.}, vol. 11, no. 10, pp. 3618-3631, Oct. 2012.

\bibitem{Xu2014} Q. Xu, X. Li, H. Ji, and X. Du, "Energy-efficient resource allocation for heterogeneous services in OFDMA downlink networks: Systematic perspective," \textit{IEEE Trans. Veh. Technol.}, vol. 63, no. 5, pp. 2071-2082, June 2014.

\bibitem{Xu2013} J. Xu and L. Qiu, "Energy efficiency optimization for MIMO broadcast channels," \textit{IEEE Trans. Wireless Commun.}, vol. 12, no. 2, pp. 690-701, Feb. 2013.

\bibitem{Du2014} B. Du, C. Pan, W. Zhang, and M. Chen, "Distributed energy-efficient power optimization for CoMP systems with max-min fairness," \textit{IEEE Commun. Lett.}, vol. 18, no. 6, pp. 999-1002, Jun. 2014.

\bibitem{He2013} S. He, Y. Huang, S. Jin, and L. Yang, "Coordinated beamforming for energy efficient transmission in multicell multiuser systems," \textit{IEEE Trans. Commun.}, vol. 61, no. 12, pp. 4961-4971, Dec. 2013.

\bibitem{He2014} S. He, Y. Huang, L. Yang, and B. Ottersten, "Coordinated multicell multiuser precoding for maximizing weighted sum energy efficiency," \textit{IEEE Trans. Signal Process.}, vol. 62, no. 3, pp. 741-751, Feb. 2014.

\bibitem{Venturino2015} L. Venturino, A. Zappone, C. Risi, and S. Buzzi, "Energy-efficient scheduling and power allocation in downlink OFDMA networks with base station coordination," \textit{IEEE Trans. Wireless Commun.}, vol. 14, no. 1, pp. 1-14, 2015.

\bibitem{Miao2011} G. Miao, N. Himayat, G. Li, and S. Talwar, "Distributed interference-aware energy-efficient power optimization," \textit{IEEE Trans. Wireless Commun.}, vol. 10, no. 4, pp. 1323-1333, Apr. 2011.

\bibitem{Scutari2008} G. Scutari, D. Palomar, and S. Barbarossa, "Optimal linear precoding strategies for wideband noncooperative systems based on game theory part I: Nash equilibria," \textit{IEEE Trans. Signal Process.}, vol. 56, no. 3, pp. 1230-1249, March 2008.

\bibitem{Zappone2013} A. Zappone, Z. Chong, E. Jorswieck, and S. Buzzi, "Energy-aware competitive power control in relay-assisted interference wireless networks," \textit{IEEE Trans. Wireless Commun.}, vol. 12, no. 4, pp. 1860-1871, Apr. 2013.

\bibitem{Stupia2015} I. Stupia, L. Sanguinetti, G. Bacci, and L. Vandendorpe, "Power control in networks with heterogeneous users: A quasi-variational inequality approach," \textit{IEEE Trans. Signal Process.}, vol. 63, no. 21, pp. 5691-5705, November 2015.

\bibitem{Bacci2015} G. Bacci, E. Belmega, P. Mertikopoulos, and L. Sanguinetti, "Energy-aware competitive power allocation for heterogeneous networks under QoS constraints," \textit{IEEE Trans. Wireless Commun.}, vol. 14, no. 9, pp. 4728-4742, Sept 2015.

\bibitem{Yates1995} R. D. Yates, "A framework for uplink power control in cellular radio systems," \textit{IEEE J. Sel. Areas Commun.}, vol. 13, no. 7, pp. 1341-1347, September 1995.

\end{thebibliography}

\end{document}
% Prérequis packages : amsmath, amssymb, algorithm2e, graphicx,
%   booktabs, multirow, tabularx, subcaption, hyperref, xcolor

% ============================================================
%            CHAPITRE : APPROCHE PAR DEEP Q-NETWORK
% ============================================================
\section{Contrôle de Puissance par Deep Q-Network}
\label{chap:dqn}

\subsection{Introduction}

Le chapitre précédent a présenté l'approche analytique de Zappone \textit{et al.}~\cite{zappone2015} pour la maximisation décentralisée de l'efficacité énergétique (EE) dans les réseaux 5G Massive MIMO, reposant sur la dynamique de meilleure réponse (BRD) couplée à l'algorithme de Dinkelbach. Bien que cette méthode converge vers l'Équilibre de Nash Généralisé (GNE) sous certaines conditions, elle nécessite une connaissance explicite du modèle de canal et la résolution itérative d'un problème d'optimisation fractionnaire à chaque mise à jour.

Dans ce chapitre, nous proposons une approche alternative basée sur l'apprentissage par renforcement profond (\textit{Deep Reinforcement Learning}), et plus spécifiquement sur le \textbf{Deep Q-Network (DQN)}~\cite{mnih2015}. L'objectif est de permettre à chaque utilisateur $k$ d'apprendre de manière autonome une politique de contrôle de puissance qui maximise son EE individuelle, sans connaissance explicite du modèle mathématique du canal. L'approche est intrinsèquement décentralisée : chaque utilisateur dispose de son propre agent DQN qui prend des décisions indépendantes, tout en interagissant avec les autres agents via l'interférence partagée dans l'environnement sans fil.

% ============================================================
\subsection{Formulation en Processus Décisionnel de Markov}
\label{sec:mdp}

Pour appliquer l'apprentissage par renforcement, nous reformulons le problème de contrôle de puissance comme un \textbf{Processus Décisionnel de Markov} (MDP). Chaque utilisateur $k \in \{1, \dots, K\}$ est modélisé comme un agent autonome interagissant avec un environnement partagé. Le MDP de l'utilisateur $k$ est défini par le tuple $(\mathcal{S}_k, \mathcal{A}_k, \mathcal{P}, \mathcal{R}_k, \gamma)$, dont les composantes sont détaillées ci-après.

\subsubsection{Espace d'états}
\label{sec:state_space}

L'état observé par l'utilisateur $k$ à l'instant $t$ est un vecteur de dimension 6 :
\begin{equation}
    \mathbf{s}_k^{(t)} = \begin{bmatrix}
        s_1^{(t)} \\ s_2^{(t)} \\ s_3^{(t)} \\ s_4^{(t)} \\ s_5^{(t)} \\ s_6^{(t)}
    \end{bmatrix}
    \label{eq:state_vector}
\end{equation}
dont les composantes sont définies comme suit :

\begin{enumerate}
    \item \textbf{SINR normalisé en dB} :
    \begin{equation}
        s_1^{(t)} = \text{clip}\!\left(\frac{\text{SINR}_k^{\text{dB}}}{25}, -1, 1\right), \quad \text{SINR}_k^{\text{dB}} = 10\log_{10}(\gamma_k^{(t)})
        \label{eq:state_sinr}
    \end{equation}
    où la normalisation par 25~dB correspond à la plage pratique du SINR dans le système considéré.

    \item \textbf{Puissance normalisée (échelle logarithmique)} :
    \begin{equation}
        s_2^{(t)} = \frac{\log_{10}(p_k^{(t)}) - \log_{10}(p_{\min})}{\log_{10}(p_{\max}) - \log_{10}(p_{\min})}
        \label{eq:state_power}
    \end{equation}
    Cette représentation en échelle log permet une résolution uniforme sur toute la plage de puissance.

    \item \textbf{Interférence normalisée} :
    \begin{equation}
        s_3^{(t)} = \frac{\log_{10}\left(\sum_{j \neq k} p_j^{(t)}\right) - \log_{10}\left((K\!-\!1)\,p_{\min}\right)}{\log_{10}\left((K\!-\!1)\,p_{\max}\right) - \log_{10}\left((K\!-\!1)\,p_{\min}\right)}
        \label{eq:state_interf}
    \end{equation}

    \item \textbf{Gain de canal équivalent normalisé} :
    \begin{equation}
        s_4^{(t)} = \text{clip}\!\left(\frac{\log_{10}(\mu_k^{(t)})}{10}, 0, 1\right)
        \label{eq:state_mu}
    \end{equation}
    où $\mu_k = \frac{\alpha_k}{\sigma^2 + \sum_{j \neq k} \omega_{kj}\,p_j}$ est le gain de canal équivalent utilisé dans la meilleure réponse de Dinkelbach. Cette composante est cruciale car elle encode l'information nécessaire pour déterminer la puissance optimale.

    \item \textbf{Débit normalisé} :
    \begin{equation}
        s_5^{(t)} = \frac{R_k^{(t)}}{R_{\max}}, \quad R_{\max} = \log_2(1 + \bar{\gamma})
        \label{eq:state_rate}
    \end{equation}
    où $R_{\max}$ est le débit maximal théorique limité par les imperfections matérielles ($\bar{\gamma} = 1/\varepsilon^2$).

    \item \textbf{Marge QoS} :
    \begin{equation}
        s_6^{(t)} = \text{clip}\!\left(\frac{R_k^{(t)} - \theta_k}{\theta_k}, -1, 1\right)
        \label{eq:state_qos}
    \end{equation}
    où $\theta_k$ est le débit minimal requis par la contrainte de qualité de service. Une valeur positive indique que la contrainte est satisfaite.
\end{enumerate}

Toutes les composantes sont normalisées dans l'intervalle $[-1, 1]$ ou $[0, 1]$ afin de faciliter l'apprentissage du réseau de neurones. Le choix d'une représentation en échelle logarithmique pour la puissance et l'interférence reflète la nature multiplicative des signaux radio et améliore la sensibilité de l'agent aux variations dans les régimes de faible puissance, où se situe typiquement le point de fonctionnement optimal en termes d'EE.

\subsubsection{Espace d'actions}
\label{sec:action_space}

L'espace d'actions est \textbf{discret} et composé de $|\mathcal{A}| = 32$ niveaux de puissance répartis de manière logarithmique entre $p_{\min}$ et $p_{\max}$ :
\begin{equation}
    \mathcal{A} = \left\{ p_{\min} \cdot \left(\frac{p_{\max}}{p_{\min}}\right)^{\frac{i}{N_a - 1}} \;\middle|\; i = 0, 1, \dots, N_a - 1 \right\}
    \label{eq:action_space}
\end{equation}
avec $p_{\min} = -20$~dBm ($0{,}01$~mW) et $p_{\max} = 23$~dBm ($200$~mW), soit un rapport de 1:20\,000.

L'espacement logarithmique (géométrique) est motivé par le fait que l'efficacité énergétique optimale se situe typiquement dans la partie inférieure de la plage de puissance. Cette discrétisation offre ainsi une \textbf{résolution fine} dans la zone optimale (entre 0,01 et 5~mW) tout en couvrant l'ensemble de la plage autorisée. L'action $a_k^{(t)} \in \{0, 1, \dots, 31\}$ sélectionne directement le niveau de puissance d'émission $p_k^{(t)} = \mathcal{A}[a_k^{(t)}]$.

\subsubsection{Fonction de récompense}
\label{sec:reward}

La fonction de récompense est conçue pour guider l'agent vers la maximisation de l'efficacité énergétique tout en respectant les contraintes de QoS et en encourageant l'économie de puissance :

\begin{equation}
    r_k^{(t)} = \text{clip}\!\left(\alpha_{\text{ee}} \cdot \frac{\eta_k^{(t)}}{\eta_{\text{ref}}} - \beta_{\text{qos}} \cdot \mathbb{1}_{\text{QoS}} \cdot \left(0{,}3 + 0{,}7 \cdot \min\!\left(\frac{\theta_k - R_k^{(t)}}{\theta_k}, 1\right)\right) - \gamma_p \cdot \frac{p_k^{(t)}}{p_{\max}},\; -3,\; 2\right)
    \label{eq:reward}
\end{equation}

où :
\begin{itemize}
    \item $\eta_k^{(t)} = \frac{B \cdot \log_2(1 + \gamma_k^{(t)})}{p_c + p_k^{(t)}}$ est l'efficacité énergétique en bit/Joule ;
    \item $\eta_{\text{ref}} = 3 \times 10^8$ bit/J (= 300~Mbit/J) est la référence de normalisation, choisie pour que le terme EE soit dans l'intervalle $[0, 1]$ au voisinage de la solution optimale ;
    \item $\alpha_{\text{ee}} = 1{,}0$ est le coefficient de récompense EE ;
    \item $\beta_{\text{qos}} = 5{,}0$ est le coefficient de pénalité QoS, volontairement élevé pour fortement décourager les violations de la contrainte de débit minimal ;
    \item $\gamma_p = 0{,}02$ est le coefficient de pénalité de puissance, encourageant l'économie d'énergie ;
    \item $\mathbb{1}_{\text{QoS}} = \mathbb{1}_{R_k < \theta_k}$ est l'indicatrice de violation QoS.
\end{itemize}

La pénalité QoS est composée d'un terme fixe ($0{,}3$) activé dès qu'il y a violation, et d'un terme proportionnel au déficit relatif $(\theta_k - R_k)/\theta_k$, ce qui permet de différencier les violations légères des violations sévères. Le \textit{soft clipping} dans $[-3, 2]$ assure la stabilité numérique de l'entraînement en évitant les récompenses extrêmes.

\subsubsection{Dynamique de l'environnement}
\label{sec:env_dynamics}

L'environnement simule un réseau 5G Massive MIMO avec les caractéristiques suivantes :

\begin{itemize}
    \item \textbf{Canal dynamique} : à chaque pas de temps, le canal est mis à jour par un évanouissement rapide résiduel :
    \begin{equation}
        d_k^{\text{eff}}(t) = d_k \cdot (1 + \delta_k(t)), \quad \delta_k \sim \mathcal{N}\!\left(0, \frac{1}{\sqrt{M}}\right)
        \label{eq:channel_hardening}
    \end{equation}
    Ce modèle capture l'effet de \textit{channel hardening} typique du Massive MIMO avec $M = 50$ antennes, donnant des variations de $\pm 1$--$2$~dB.

    \item \textbf{Réinitialisation par épisode} : à chaque épisode, les $K = 5$ utilisateurs sont repositionnés aléatoirement dans la cellule (rayon $R_{\text{cell}} = 400$~m, distance minimale $d_{\min} = 35$~m), et un nouveau réalisation de l'affaiblissement grande échelle (perte de trajet + \textit{shadowing}) est tirée :
    \begin{equation}
        d_k = 10^{-(L_k + S_k)/10}, \quad L_k = 10\alpha\log_{10}(r_k), \quad S_k \sim \mathcal{N}(0, \sigma_s^2)
        \label{eq:large_scale_fading}
    \end{equation}
    avec $\alpha = 3{,}76$ et $\sigma_s = 8$~dB.

    \item \textbf{Durée d'épisode} : chaque épisode comprend $T = 200$ pas de temps, durant lesquels les agents interagissent simultanément.
\end{itemize}

% ============================================================
\subsection{Architecture du Deep Q-Network}
\label{sec:dqn_architecture}

\subsubsection{Principe du DQN}

Le Deep Q-Network~\cite{mnih2015} approxime la fonction de valeur action-état $Q^*(s, a)$ par un réseau de neurones profond $Q(s, a; \boldsymbol{\theta})$ paramétré par $\boldsymbol{\theta}$. La politique gloutonne (\textit{greedy}) correspondante est :
\begin{equation}
    \pi^*(s) = \arg\max_{a \in \mathcal{A}} Q(s, a; \boldsymbol{\theta})
    \label{eq:greedy_policy}
\end{equation}

L'objectif est de minimiser l'erreur temporelle différentielle (\textit{TD error}) entre la valeur Q courante et la valeur cible :
\begin{equation}
    \mathcal{L}(\boldsymbol{\theta}) = \mathbb{E}_{(s,a,r,s') \sim \mathcal{D}} \left[ \ell\!\left(r + \gamma \max_{a'} Q(s', a'; \boldsymbol{\theta}^{-}) - Q(s, a; \boldsymbol{\theta})\right) \right]
    \label{eq:dqn_loss}
\end{equation}
où $\mathcal{D}$ est le \textit{replay buffer}, $\boldsymbol{\theta}^{-}$ les paramètres du réseau cible, $\gamma = 0{,}95$ le facteur d'actualisation, et $\ell(\cdot)$ la fonction de perte.

\subsubsection{Architecture du réseau de neurones}

Le réseau de neurones utilisé pour chaque agent DQN est un perceptron multicouche (\textit{Multi-Layer Perceptron}, MLP) composé de trois couches cachées entièrement connectées (\textit{fully connected}). L'architecture est la suivante :

\begin{equation}
    \underbrace{\mathbb{R}^6}_{\text{Entrée (état)}} \xrightarrow{\mathbf{W}_1, \mathbf{b}_1} \underbrace{\mathbb{R}^{128}}_{\text{Couche 1}} \xrightarrow{\text{ReLU}} \underbrace{\mathbb{R}^{256}}_{\text{Couche 2}} \xrightarrow{\text{ReLU}} \underbrace{\mathbb{R}^{128}}_{\text{Couche 3}} \xrightarrow{\text{ReLU}} \underbrace{\mathbb{R}^{32}}_{\text{Sortie (Q-values)}}
    \label{eq:architecture}
\end{equation}

\begin{itemize}
    \item \textbf{Couche d'entrée} : 6 neurones correspondant aux 6 composantes de l'état $\mathbf{s}_k$.
    \item \textbf{Couche cachée 1} : 128 neurones avec activation ReLU.
    \item \textbf{Couche cachée 2} : 256 neurones avec activation ReLU.
    \item \textbf{Couche cachée 3} : 128 neurones avec activation ReLU.
    \item \textbf{Couche de sortie} : 32 neurones (un par action), sans activation — les sorties représentent directement les Q-valeurs $Q(s, a_i; \boldsymbol{\theta})$ pour $i = 0, \dots, 31$.
\end{itemize}

L'architecture en forme de \textit{bottleneck} (128 $\to$ 256 $\to$ 128) permet d'abord une expansion de la représentation dans l'espace latent, puis une compression vers la dimension de sortie, favorisant l'extraction de caractéristiques pertinentes.

Les poids sont initialisés selon la méthode de \textbf{Xavier uniforme}~\cite{glorot2010} :
\begin{equation}
    W_{ij} \sim \mathcal{U}\left(-\sqrt{\frac{6}{n_{\text{in}} + n_{\text{out}}}},\; +\sqrt{\frac{6}{n_{\text{in}} + n_{\text{out}}}}\right)
    \label{eq:xavier_init}
\end{equation}
et les biais sont initialisés à zéro. Le nombre total de paramètres entraînables par agent est de :
\begin{equation}
    N_{\text{params}} = (6 \times 128 + 128) + (128 \times 256 + 256) + (256 \times 128 + 128) + (128 \times 32 + 32) = 70\,944
    \label{eq:n_params}
\end{equation}

\begin{figure*}[htbp]
    \centering
    \includegraphics[width=0.85\textwidth]{figures/dqn_architecture.png}
    \caption{Architecture du réseau DQN pour le contrôle de puissance. L'entrée est le vecteur d'état à 6 dimensions et la sortie fournit les Q-valeurs pour les 32 niveaux de puissance discrets.}
    \label{fig:dqn_architecture}
\end{figure*}

% ============================================================
\subsection{Mécanismes d'apprentissage}
\label{sec:learning_mechanisms}

\subsubsection{Double DQN}
\label{sec:double_dqn}

Afin de réduire la surestimation des Q-valeurs inhérente au DQN standard~\cite{vanHasselt2016}, nous implémentons la variante \textbf{Double DQN}. Le principe consiste à découpler la \textbf{sélection} de l'action (effectuée par le réseau de politique $\boldsymbol{\theta}$) de son \textbf{évaluation} (effectuée par le réseau cible $\boldsymbol{\theta}^{-}$) :

\begin{equation}
    y^{\text{DDQN}} = r + \gamma \cdot Q\!\left(s',\; \underbrace{\arg\max_{a'} Q(s', a'; \boldsymbol{\theta})}_{\text{sélection par réseau politique}};\; \boldsymbol{\theta}^{-}\right)
    \label{eq:double_dqn_target}
\end{equation}

Cette technique réduit significativement le biais d'estimation positif et améliore la stabilité de l'apprentissage, ce qui est particulièrement important dans notre contexte multi-agent où les environnements sont non-stationnaires du point de vue de chaque agent.

\subsubsection{Replay d'expérience}
\label{sec:experience_replay}

Chaque agent maintient un \textbf{tampon de rejeu} (\textit{experience replay buffer}) $\mathcal{D}_k$ de capacité $|\mathcal{D}| = 100\,000$ transitions. Chaque transition stockée est un tuple :
\begin{equation}
    e^{(t)} = \left(s_k^{(t)},\; a_k^{(t)},\; r_k^{(t)},\; s_k^{(t+1)},\; d^{(t)}\right)
    \label{eq:transition}
\end{equation}
où $d^{(t)} \in \{0, 1\}$ est l'indicateur de fin d'épisode. Le tampon fonctionne selon le principe FIFO (\textit{First-In, First-Out}) : lorsque la capacité maximale est atteinte, les transitions les plus anciennes sont remplacées.

À chaque pas d'entraînement, un \textbf{mini-batch} de $B = 256$ transitions est échantillonné uniformément de manière aléatoire :
\begin{equation}
    \{(s_i, a_i, r_i, s'_i, d_i)\}_{i=1}^{B} \sim \text{Uniforme}(\mathcal{D}_k)
    \label{eq:sampling}
\end{equation}

Le replay d'expérience confère deux avantages fondamentaux :
\begin{itemize}
    \item \textbf{Décorrélation temporelle} : les transitions échantillonnées aléatoirement brisent les corrélations séquentielles, satisfaisant l'hypothèse i.i.d. de la descente de gradient stochastique.
    \item \textbf{Efficacité en données} : chaque transition peut être réutilisée dans plusieurs mises à jour, améliorant l'exploitation des données collectées.
\end{itemize}

L'entraînement ne commence qu'une fois le tampon contenant au moins $|\mathcal{D}|_{\min} = 2\,000$ transitions, garantissant une diversité suffisante dans les premiers mini-batchs.

\subsubsection{Réseau cible et mise à jour douce}
\label{sec:target_network}

Pour stabiliser l'apprentissage, nous utilisons un \textbf{réseau cible} (\textit{target network}) $Q(s, a; \boldsymbol{\theta}^{-})$ distinct du réseau de politique $Q(s, a; \boldsymbol{\theta})$. Les paramètres du réseau cible sont mis à jour par \textbf{mise à jour douce} (\textit{soft update} ou \textit{Polyak averaging}) :

\begin{equation}
    \boldsymbol{\theta}^{-} \leftarrow \tau \cdot \boldsymbol{\theta} + (1 - \tau) \cdot \boldsymbol{\theta}^{-}, \quad \tau = 0{,}005
    \label{eq:soft_update}
\end{equation}

Cette mise à jour est effectuée après chaque pas d'optimisation. Le coefficient $\tau = 0{,}005$ assure une transition lente et régulière du réseau cible vers le réseau de politique, évitant les oscillations et divergences liées à une cible en mouvement rapide. Par rapport à la mise à jour dure (\textit{hard update}) périodique, la mise à jour douce offre une stabilité accrue car le réseau cible évolue de manière continue plutôt que par sauts discontinus.

\subsubsection{Fonction de perte de Huber}
\label{sec:huber_loss}

La fonction de perte utilisée pour l'optimisation est la \textbf{perte de Huber} (aussi appelée \textit{Smooth L1 Loss}) :
\begin{equation}
    \ell_{\delta}(x) = \begin{cases}
        \frac{1}{2}x^2 & \text{si } |x| \leq \delta \\
        \delta\left(|x| - \frac{1}{2}\delta\right) & \text{sinon}
    \end{cases}
    \label{eq:huber_loss}
\end{equation}
avec $\delta = 1{,}0$. La perte de Huber combine les avantages de la MSE (sensibilité aux petites erreurs) et de la MAE (robustesse aux valeurs aberrantes), ce qui est crucial dans notre environnement multi-agent où les récompenses peuvent être bruitées par les actions concurrentes des autres agents.

\subsubsection{Politique d'exploration $\varepsilon$-gloutonne}
\label{sec:epsilon_greedy}

L'exploration du space d'actions est réalisée par une politique $\varepsilon$-gloutonne (\textit{$\varepsilon$-greedy}) :
\begin{equation}
    a_k^{(t)} = \begin{cases}
        \text{Uniforme}(\mathcal{A}) & \text{avec probabilité } \varepsilon_t \\
        \arg\max_{a} Q(s_k^{(t)}, a; \boldsymbol{\theta}_k) & \text{avec probabilité } 1 - \varepsilon_t
    \end{cases}
    \label{eq:epsilon_greedy}
\end{equation}

Le taux d'exploration $\varepsilon_t$ décroît de manière exponentielle au fil des épisodes :
\begin{equation}
    \varepsilon_t = \max\!\left(\varepsilon_{\text{fin}},\; \varepsilon_0 \cdot \lambda^t\right), \quad \varepsilon_0 = 1{,}0, \quad \varepsilon_{\text{fin}} = 0{,}02, \quad \lambda = 0{,}9995
    \label{eq:epsilon_decay}
\end{equation}

Ce schéma de décroissance assure :
\begin{itemize}
    \item Une \textbf{phase d'exploration intensive} initiale (épisodes 1--1\,000, $\varepsilon > 0{,}6$) permettant une couverture large de l'espace état-action.
    \item Une \textbf{transition progressive} (épisodes 1\,000--7\,000) où l'agent bascule graduellement vers l'exploitation.
    \item Une \textbf{phase d'exploitation} finale (épisodes 7\,000--10\,000, $\varepsilon \approx 0{,}02$) où l'agent utilise quasi-exclusivement la politique apprise, tout en conservant une exploration résiduelle de 2\% pour maintenir une certaine adaptabilité.
\end{itemize}

L'épisode où $\varepsilon$ passe sous le seuil $0{,}1$ (exploitation dominante) se situe aux alentours de l'épisode $4\,600$, ce qui correspond au régime où les Q-valeurs sont suffisamment stables pour guider les décisions.

\begin{figure}[htbp]
    \centering
    \includegraphics[width=1\textwidth]{figures/epsilon_decay.png}
    \caption{Décroissance exponentielle du taux d'exploration $\varepsilon$ au cours des $10\,000$ épisodes d'entraînement. La ligne pointillée indique le seuil $\varepsilon = 0{,}1$.}
    \label{fig:epsilon_decay}
\end{figure}

\subsubsection{Gradient clipping}
\label{sec:grad_clip}

Pour prévenir l'instabilité numérique causée par des gradients de grande magnitude, une \textbf{troncature du gradient} (\textit{gradient clipping}) est appliquée avant chaque mise à jour des poids :
\begin{equation}
    \nabla_{\boldsymbol{\theta}} \mathcal{L} \leftarrow \frac{\nabla_{\boldsymbol{\theta}} \mathcal{L}}{\max\!\left(1, \frac{\|\nabla_{\boldsymbol{\theta}} \mathcal{L}\|_2}{G_{\max}}\right)}, \quad G_{\max} = 1{,}0
    \label{eq:grad_clip}
\end{equation}

Cette technique est d'autant plus importante dans le contexte multi-agent, où la non-stationnarité induite par l'apprentissage simultané de tous les agents peut engendrer des gradients instables.

% ============================================================
\subsection{Entraînement multi-agent décentralisé}
\label{sec:multi_agent}

\subsubsection{Independent Q-Learning (IQL)}

L'approche adoptée est le \textbf{Independent Q-Learning} (IQL)~\cite{tan1993}, où chaque utilisateur $k$ possède son propre agent DQN qui apprend indépendamment. Les agents interagissent uniquement à travers l'environnement partagé : le SINR de l'utilisateur $k$ dépend des puissances d'émission de tous les utilisateurs via l'interférence :
\begin{equation}
    \gamma_k = \frac{\alpha_k \, p_k}{\sigma^2 + \phi_k \, p_k + \sum_{j \neq k} \omega_{kj} \, p_j}
    \label{eq:sinr_coupling}
\end{equation}

Du point de vue de l'utilisateur $k$, les autres utilisateurs font partie de l'environnement, qui devient \textbf{non-stationnaire} puisque les politiques des autres agents évoluent au cours de l'entraînement. Bien que cela viole formellement les hypothèses de stationnarité du MDP, l'IQL fonctionne remarquablement bien en pratique pour les jeux coopératifs et compétitifs de taille modérée~\cite{lanctot2017}.

\subsubsection{Protocole d'entraînement}

Le protocole d'entraînement est résumé dans l'Algorithme~\ref{alg:training}.

\begin{algorithm}[htbp]
\caption{Entraînement multi-agent DQN décentralisé}
\label{alg:training}
\SetKwInOut{Input}{Entrées}
\SetKwInOut{Output}{Sorties}
\Input{$K = 5$ agents, $N_{\text{ep}} = 10\,000$ épisodes, $T = 200$ pas/épisode}
\Output{Politiques optimales $\{\boldsymbol{\theta}_k^*\}_{k=1}^{K}$}
\BlankLine
Initialiser $K$ réseaux de politique $\{Q(\cdot; \boldsymbol{\theta}_k)\}$ et cibles $\{Q(\cdot; \boldsymbol{\theta}_k^{-})\}$\;
Initialiser $K$ tampons de rejeu $\{\mathcal{D}_k\}$ de capacité $100\,000$\;
$\varepsilon \leftarrow 1{,}0$\;
\BlankLine
\For{épisode $= 1$ \KwTo $N_{\text{ep}}$}{
    Réinitialiser l'environnement (nouvelles positions + canal)\;
    Observer les états initiaux $\{s_k^{(0)}\}_{k=1}^{K}$\;
    \BlankLine
    \For{$t = 0$ \KwTo $T-1$}{
        \For{$k = 1$ \KwTo $K$ \textrm{(en parallèle)}}{
            Sélectionner $a_k^{(t)}$ par politique $\varepsilon$-gloutonne (Éq.~\ref{eq:epsilon_greedy})\;
        }
        Exécuter les actions $\{a_k^{(t)}\}$ dans l'environnement\;
        Observer récompenses $\{r_k^{(t)}\}$ et nouveaux états $\{s_k^{(t+1)}\}$\;
        \BlankLine
        \For{$k = 1$ \KwTo $K$}{
            Stocker $(s_k^{(t)}, a_k^{(t)}, r_k^{(t)}, s_k^{(t+1)}, d^{(t)})$ dans $\mathcal{D}_k$\;
            \If{$|\mathcal{D}_k| \geq 2\,000$ \textbf{et} $t \bmod 16 = 0$}{
                Échantillonner un mini-batch de 256 transitions de $\mathcal{D}_k$\;
                Calculer les cibles Double DQN (Éq.~\ref{eq:double_dqn_target})\;
                Mettre à jour $\boldsymbol{\theta}_k$ par descente de gradient (perte de Huber)\;
                Mise à jour douce : $\boldsymbol{\theta}_k^{-} \leftarrow 0{,}005 \cdot \boldsymbol{\theta}_k + 0{,}995 \cdot \boldsymbol{\theta}_k^{-}$\;
            }
        }
    }
    $\varepsilon \leftarrow \max(0{,}02,\; \varepsilon \times 0{,}9995)$\;
    \BlankLine
    \If{épisode $\bmod$ 500 $= 0$}{
        Sauvegarder les checkpoints $\{\boldsymbol{\theta}_k\}$\;
    }
}
\end{algorithm}

Un détail d'implémentation important est la \textbf{fréquence d'entraînement réduite} : l'optimisation du réseau n'est effectuée que tous les 16 pas de temps ($t \bmod 16 = 0$), ce qui réduit le surcoût de transfert CPU$\leftrightarrow$GPU et évite l'overfitting sur des transitions fortement corrélées au sein d'un même épisode.

\subsubsection{Optimisation}

L'optimiseur utilisé est \textbf{Adam}~\cite{kingma2015} avec un taux d'apprentissage $\eta = 10^{-4}$ :
\begin{equation}
    \boldsymbol{\theta}_{t+1} = \boldsymbol{\theta}_t - \eta \cdot \frac{\hat{m}_t}{\sqrt{\hat{v}_t} + \epsilon_{\text{Adam}}}
    \label{eq:adam}
\end{equation}
où $\hat{m}_t$ et $\hat{v}_t$ sont les estimateurs corrigés du premier et second moment des gradients. Le taux d'apprentissage relativement faible ($10^{-4}$) est choisi pour assurer une convergence stable dans l'environnement multi-agent non-stationnaire.

% ============================================================
\subsection{Sauvegarde et évaluation des modèles}
\label{sec:checkpoints}

\subsubsection{Système de checkpoints}

Des points de contrôle (\textit{checkpoints}) sont sauvegardés tous les 500 épisodes, contenant :
\begin{itemize}
    \item Les paramètres du réseau de politique $\boldsymbol{\theta}_k$ et du réseau cible $\boldsymbol{\theta}_k^{-}$
    \item L'état de l'optimiseur Adam (moments)
    \item La valeur courante de $\varepsilon$
    \item Le compteur d'étapes d'entraînement
    \item L'historique des pertes
\end{itemize}

Cela permet de reprendre l'entraînement en cas d'interruption et d'évaluer les performances intermédiaires de la politique à différents stades de l'apprentissage.

\subsubsection{Protocole d'évaluation}

L'évaluation des modèles entraînés est réalisée en mode \textbf{glouton} ($\varepsilon = 0$) sur $N_{\text{eval}} = 50$ réalisations de canal indépendantes. Pour chaque réalisation, les agents sont exécutés pendant $T = 200$ pas afin de stabiliser les puissances d'émission avant de mesurer l'EE résultante.

\begin{figure*}[htbp]
    \centering
    \includegraphics[width=0.95\textwidth]{figures/training_curves.png}
    \caption{Courbes d'entraînement des agents DQN décentralisés sur $10\,000$ épisodes : (a) efficacité énergétique moyenne, (b) reward cumulatif, (c) perte d'entraînement, (d) SINR moyen, (e) puissance d'émission moyenne, (f) taux d'exploration $\varepsilon$.}
    \label{fig:training_curves}
\end{figure*}

% ============================================================
\subsection{Récapitulatif des hyperparamètres}
\label{sec:hyperparams_summary}

Le Tableau~\ref{tab:hyperparams} regroupe l'ensemble des hyperparamètres utilisés pour l'entraînement.

\begin{table}[H]
    \centering
    \caption{Hyperparamètres de l'entraînement DQN}
    \label{tab:hyperparams}
    \begin{tabular}{@{}llr@{}}
        \toprule
        \textbf{Catégorie} & \textbf{Paramètre} & \textbf{Valeur} \\
        \midrule
        \multirow{4}{*}{Architecture} 
            & Couches cachées & $[128, 256, 128]$ \\
            & Activation & ReLU \\
            & Initialisation & Xavier uniforme \\
            & Params/agent & 70\,944 \\
        \midrule
        \multirow{5}{*}{Optimisation}
            & Optimiseur & Adam \\
            & Taux d'apprentissage $\eta$ & $10^{-4}$ \\
            & Taille du batch $B$ & 256 \\
            & Fonction de perte & Huber ($\delta = 1{,}0$) \\
            & Gradient clipping & $G_{\max} = 1{,}0$ \\
        \midrule
        \multirow{3}{*}{Exploration}
            & $\varepsilon_0$ & 1,0 \\
            & $\varepsilon_{\text{fin}}$ & 0,02 \\
            & Décroissance $\lambda$ & 0,9995/épisode \\
        \midrule
        \multirow{4}{*}{Replay \& cible}
            & Taille du buffer $|\mathcal{D}|$ & 100\,000 \\
            & Buffer minimal $|\mathcal{D}|_{\min}$ & 2\,000 \\
            & Soft update $\tau$ & 0,005 \\
            & Fréquence entraîn. & 1/16 pas \\
        \midrule
        \multirow{3}{*}{Entraînement}
            & Facteur de discount $\gamma$ & 0,95 \\
            & Épisodes $N_{\text{ep}}$ & 10\,000 \\
            & Pas/épisode $T$ & 200 \\
        \midrule
        \multirow{3}{*}{Récompense}
            & $\alpha_{\text{ee}}$ & 1,0 \\
            & $\beta_{\text{qos}}$ & 5,0 \\
            & $\gamma_p$ & 0,02 \\
        \bottomrule
    \end{tabular}
\end{table}


% ============================================================
% ============================================================
%     CHAPITRE : RÉSULTATS ET COMPARAISONS
% ============================================================
% ============================================================

\section{Résultats Expérimentaux et Analyse Comparative}

\label{chap:results}

\section{Introduction}

Ce chapitre présente une analyse comparative détaillée entre l'approche DQN décentralisée développée dans le Chapitre~\ref{chap:dqn} et les méthodes de référence. Nous évaluons les performances sur $N = 50$ réalisations de canal indépendantes, garantissant une comparaison statistiquement significative. Toutes les méthodes sont testées sur les \textbf{mêmes} réalisations de canal (échantillons appariés), assurant l'équité de la comparaison.

\section{Méthodes comparées}
\label{sec:methods}

Cinq méthodes de contrôle de puissance sont évaluées :

\begin{enumerate}
    \item \textbf{Article (BRD + Dinkelbach)} : l'algorithme de référence de Zappone \textit{et al.}~\cite{zappone2015}, basé sur la dynamique de meilleure réponse avec résolution interne par l'algorithme de Dinkelbach (Chapitre~\ref{chap:article_algo}).

    \item \textbf{DQN (entraîné)} : notre approche par apprentissage par renforcement profond, avec $K = 5$ agents DQN indépendants entraînés sur $10\,000$ épisodes (Chapitre~\ref{chap:dqn}).

    \item \textbf{Optimal Uniform} : recherche exhaustive de la puissance uniforme $p^*$ optimale parmi 100 niveaux géométriquement espacés, où tous les utilisateurs émettent à la même puissance. Cette méthode sert de baseline analytique simple.

    \item \textbf{Max Power} : tous les utilisateurs émettent à la puissance maximale $p_{\max} = 23$~dBm ($200$~mW). Ce baseline naïf illustre la sous-optimalité de l'approche brute en termes d'EE.

    \item \textbf{Random Power} : chaque utilisateur choisit une puissance aléatoire uniforme dans $[p_{\min}, p_{\max}]$. Ce baseline représente l'absence totale de stratégie.
\end{enumerate}

\section{Résultats en efficacité énergétique}
\label{sec:ee_results}

\subsection{Performances globales}

Le Tableau~\ref{tab:ee_comparison} présente les statistiques descriptives de l'efficacité énergétique pour chaque méthode sur les 50 essais.

\begin{table*}[htbp]
    \centering
    \caption{Comparaison de l'efficacité énergétique (Mbit/J) — 50 essais}
    \label{tab:ee_comparison}
    \begin{tabular}{@{}lrrrrr@{}}
        \toprule
        \textbf{Méthode} & \textbf{Moyenne} & \textbf{Écart-type} & \textbf{Min} & \textbf{Max} & \textbf{Médiane} \\
        \midrule
        \textbf{DQN (entraîné)}         & \textbf{255,16} & 47,31 & 141,27 & 313,18 & --- \\
        Article (BRD+Dinkelbach)         & 207,01          & 30,67 & 137,99 & 256,16 & --- \\
        Optimal Uniform                  & 213,04          & 52,76 & ---    & ---    & --- \\
        Random Power                     & 23,06           & 15,59 & ---    & ---    & --- \\
        Max Power                        & 10,22           & 2,55  & ---    & ---    & --- \\
        \bottomrule
    \end{tabular}
\end{table*}

L'approche DQN atteint une EE moyenne de $\boldsymbol{255{,}16}$~\textbf{Mbit/J}, surpassant significativement l'algorithme de l'article ($207{,}01$~Mbit/J) avec un \textbf{gain relatif de $+23{,}3$~\%}. Les méthodes naïves (Max Power, Random) sont nettement inférieures, confirmant que le contrôle intelligent de puissance est essentiel pour l'efficacité énergétique.

\begin{figure}[htbp]
    \centering
    \includegraphics[width=0.9\textwidth]{figures/comparison_bar.png}
    \caption{Efficacité énergétique moyenne (Mbit/J) par méthode avec barres d'erreur ($\pm 1$ écart-type). Le DQN surpasse toutes les méthodes de référence.}
    \label{fig:ee_barplot}
\end{figure}

\subsection{Distribution de l'efficacité énergétique}

La Figure~\ref{fig:ee_distribution} présente la distribution de l'EE pour l'approche DQN et l'algorithme de l'article sur les 50 essais.

\begin{figure*}[htbp]
    \centering
    \subfloat[Histogramme comparatif\label{fig:ee_histogram}]{%
        \includegraphics[width=0.48\textwidth]{figures/ee_histogram_dqn_vs_article.png}}
    \hfill
    \subfloat[Fonction de répartition cumulative (CDF)\label{fig:ee_cdf}]{%
        \includegraphics[width=0.48\textwidth]{figures/comparison_cdf.png}}
    \caption{Distribution de l'efficacité énergétique : (a) histogrammes superposés montrant le décalage favorable du DQN ; (b) CDF confirmant la dominance stochastique du DQN sur l'algorithme de l'article.}
    \label{fig:ee_distribution}
\end{figure*}

La CDF (Figure~\ref{fig:ee_cdf}) montre que pour tout seuil d'EE, la probabilité que le DQN dépasse ce seuil est plus élevée que celle de l'algorithme de l'article, ce qui constitue une \textbf{dominance stochastique du premier ordre}. La médiane du DQN se situe aux alentours de 255~Mbit/J contre 218~Mbit/J pour l'article.

\subsection{Comparaison essai par essai}

La Figure~\ref{fig:ee_scatter_trial} présente un graphique de dispersion comparant l'EE du DQN à celle de l'article pour chaque essai individuellement.

\begin{figure}[htbp]
    \centering
    \includegraphics[width=0.65\textwidth]{figures/ee_scatter_dqn_vs_article.png}
    \caption{Comparaison essai par essai de l'EE : DQN (axe Y) vs Article (axe X). Les points au-dessus de la diagonale ($y = x$) indiquent une supériorité du DQN. La grande majorité des essais favorise le DQN.}
    \label{fig:ee_scatter_trial}
\end{figure}

% ============================================================
\section{Analyse de la consommation de puissance}
\label{sec:power_results}

\subsection{Puissance moyenne par méthode}

Le Tableau~\ref{tab:power_comparison} détaille la consommation de puissance moyenne par méthode.

\begin{table}[htbp]
    \centering
    \caption{Consommation de puissance moyenne (mW) par méthode}
    \label{tab:power_comparison}
    \begin{tabular}{@{}lr@{}}
        \toprule
        \textbf{Méthode} & \textbf{Puissance moyenne (mW)} \\
        \midrule
        \textbf{DQN (entraîné)}         & \textbf{15,27} \\
        Article (BRD+Dinkelbach)         & 28,73 \\
        Optimal Uniform                  & 0,02 \\
        Max Power                        & 199,53 \\
        Random Power                     & 100,57 \\
        \bottomrule
    \end{tabular}
\end{table}

L'approche DQN consomme en moyenne $\boldsymbol{15{,}27}$~\textbf{mW}, soit une \textbf{réduction de $46{,}8$~\%} par rapport à l'algorithme de l'article ($28{,}73$~mW). Cette économie de puissance significative est obtenue tout en améliorant l'EE, ce qui démontre que le DQN apprend à opérer dans la zone optimale du compromis débit/consommation.

\begin{figure}[htbp]
    \centering
    \includegraphics[width=0.9\textwidth]{figures/comparison_boxplot.png}
    \caption{Distribution de la puissance d'émission (mW) par méthode sous forme de boîtes à moustaches (échelle logarithmique). Le DQN opère à des puissances significativement inférieures à l'algorithme de l'article.}
    \label{fig:power_boxplot}
\end{figure}

\subsection{Compromis puissance--EE}

La Figure~\ref{fig:power_vs_ee} illustre le compromis fondamental entre la puissance d'émission et l'efficacité énergétique pour toutes les méthodes.

\begin{figure}[htbp]
    \centering
    \includegraphics[width=0.9\textwidth]{figures/scatter_power_vs_ee.png}
    \caption{Diagramme de dispersion puissance (mW, échelle log) vs EE (Mbit/J) pour les 50 essais. Chaque point représente un essai; la couleur indique la méthode. Le DQN (vert) se concentre dans la zone optimale : faible puissance, haute EE.}
    \label{fig:power_vs_ee}
\end{figure}

Ce graphique met en évidence le positionnement optimal du DQN dans le coin supérieur gauche (haute EE, faible puissance), confirmant que l'agent a appris la structure du compromis débit-puissance décrit dans l'article de Zappone \textit{et al.} : l'EE est une fonction quasi-concave de la puissance d'émission, avec un maximum bien défini à des puissances relativement basses.

% ============================================================
\section{Analyse de l'algorithme de l'article}
\label{sec:article_results}

L'algorithme de référence (BRD + Dinkelbach) présente les caractéristiques suivantes sur les 50 essais :

\subsection{Convergence}

\begin{itemize}
    \item \textbf{Fraction convergée} : $100\%$ des essais convergent (tolérance relative $10^{-4}$).
    \item \textbf{Nombre moyen d'itérations} : $5{,}28$ itérations de BRD par essai.
    \item \textbf{Nombre maximum d'itérations} : $12$ (convergence toujours rapide).
\end{itemize}

La convergence rapide et systématique confirme les résultats théoriques de l'article : sous les conditions de faisabilité, la dynamique de meilleure réponse converge vers l'unique GNE.

\subsection{Efficacité énergétique détaillée}

Les statistiques descriptives de l'EE obtenue par l'algorithme de l'article sont présentées dans le Tableau~\ref{tab:article_ee_stats}.

\begin{table}[htbp]
    \centering
    \caption{Statistiques détaillées de l'EE — Algorithme de l'article (50 essais)}
    \label{tab:article_ee_stats}
    \begin{tabular}{@{}lr@{}}
        \toprule
        \textbf{Statistique} & \textbf{Valeur (Mbit/J)} \\
        \midrule
        Moyenne              & 218,07 \\
        Écart-type           & 30,10 \\
        Minimum              & 137,00 \\
        Maximum              & 268,39 \\
        Médiane              & 218,08 \\
        Percentile 5\%       & 179,47 \\
        Percentile 95\%      & 264,72 \\
        \bottomrule
    \end{tabular}
\end{table}

\subsection{Débit spectral et QoS}

\begin{itemize}
    \item \textbf{Débit spectral moyen} : $\bar{R} = 2{,}951$~bits/s/Hz (écart-type~$= 0{,}347$).
    \item \textbf{Débit maximal théorique} : $R_{\max} = \log_2(1 + 125) = 6{,}977$~bits/s/Hz.
    \item \textbf{Contrainte QoS} : $\theta_k = 20\% \times R_{\max} = 1{,}395$~bits/s/Hz.
    \item \textbf{Taux de satisfaction QoS} : $\boldsymbol{74\%}$ des utilisateurs satisfont la contrainte de débit minimal.
\end{itemize}

Le taux de satisfaction QoS de 74\% indique que même l'algorithme optimal ne parvient pas toujours à satisfaire les contraintes de débit dans des conditions de canal défavorables, ce qui souligne la difficulté inhérente du problème de faisabilité dans les systèmes à interférence.

\begin{figure*}[htbp]
    \centering
    \subfloat[Distribution de l'EE\label{fig:article_ee_dist}]{%
        \includegraphics[width=0.48\textwidth]{figures/article_ee_distribution.png}}
    \hfill
    \subfloat[Distribution de la puissance\label{fig:article_power_dist}]{%
        \includegraphics[width=0.48\textwidth]{figures/article_power_distribution.png}}
    \caption{Distributions des métriques de l'algorithme de l'article sur 50 essais : (a) EE avec annotations des percentiles P5 et P95 ; (b) puissance émise montrant une forte variance entre les réalisations de canal.}
    \label{fig:article_distributions}
\end{figure*}

% ============================================================
\section{Analyse de l'entraînement DQN}
\label{sec:training_analysis}

\subsection{Convergence de l'efficacité énergétique}

La Figure~\ref{fig:training_ee} illustre l'évolution de l'EE moyenne au cours de l'entraînement. On distingue trois phases :

\begin{enumerate}
    \item \textbf{Phase d'exploration} (épisodes 1--2\,000) : l'EE est faible et volatile, les agents explorent l'espace d'actions de manière quasi-aléatoire.
    \item \textbf{Phase d'apprentissage rapide} (épisodes 2\,000--6\,000) : l'EE croît rapidement à mesure que les Q-valeurs convergent et que l'exploitation commence à dominer.
    \item \textbf{Phase de stabilisation} (épisodes 6\,000--10\,000) : l'EE se stabilise autour de 250--300~Mbit/J, les agents ayant convergé vers des politiques quasi-optimales.
\end{enumerate}

\begin{figure}[htbp]
    \centering
    \includegraphics[width=0.8\textwidth]{figures/training_ee_convergence.png}
    \caption{Convergence de l'efficacité énergétique moyenne durant l'entraînement. La courbe lissée (fenêtre de 100 épisodes) montre une convergence progressive vers $\sim 250$~Mbit/J. La ligne pointillée rouge indique le niveau atteint par l'algorithme de l'article ($\sim 218$~Mbit/J).}
    \label{fig:training_ee}
\end{figure}

\subsection{Évolution de la perte d'entraînement}

La Figure~\ref{fig:training_loss} montre l'évolution de la perte de Huber au cours de l'entraînement. Après un pic initial correspondant aux premières mises à jour des Q-valeurs (lorsque le réseau cible est encore initialisé aléatoirement), la perte décroît progressivement, indiquant une réduction de l'erreur de prédiction temporelle. La stabilisation de la perte aux alentours de l'épisode 5\,000 confirme que les Q-valeurs convergent.

\begin{figure}[htbp]
    \centering
    \includegraphics[width=0.75\textwidth]{figures/training_loss.png}
    \caption{Perte d'entraînement (Huber loss) lissée sur une fenêtre de 100 épisodes. La décroissance progressive confirme la convergence des Q-valeurs.}
    \label{fig:training_loss}
\end{figure}

\subsection{Évolution de la puissance d'émission}

L'analyse de la puissance d'émission moyenne au cours de l'entraînement (Figure~\ref{fig:training_power}) révèle un comportement caractéristique : les agents apprennent progressivement à réduire leur puissance d'émission, passant d'une puissance élevée (exploration aléatoire) à des puissances optimales basses, cohérentes avec la maximisation de l'EE.

\begin{figure}[htbp]
    \centering
    \includegraphics[width=0.75\textwidth]{figures/training_power_evolution.png}
    \caption{Évolution de la puissance d'émission moyenne (dBm) au cours de l'entraînement. Les agents apprennent à réduire progressivement leur puissance vers le point de fonctionnement optimal.}
    \label{fig:training_power}
\end{figure}

\subsection{Évolution du SINR}

La Figure~\ref{fig:training_sinr} présente l'évolution du SINR moyen au cours de l'entraînement. On observe une stabilisation progressive du SINR à mesure que les agents convergent vers des puissances cohérentes, réduisant les fluctuations d'interférence mutuelle.

\begin{figure}[htbp]
    \centering
    \includegraphics[width=0.75\textwidth]{figures/training_sinr_evolution.png}
    \caption{Évolution du SINR moyen (dB) durant l'entraînement. La stabilisation du SINR reflète l'émergence d'un équilibre entre les agents.}
    \label{fig:training_sinr}
\end{figure}

% ============================================================
\section{Comparaison multi-critères}
\label{sec:multi_criteria}

Le Tableau~\ref{tab:multi_criteria} présente une comparaison synthétique sur l'ensemble des critères de performance.

\begin{table*}[htbp]
    \centering
    \caption{Comparaison multi-critères des méthodes de contrôle de puissance}
    \label{tab:multi_criteria}
    \renewcommand{\arraystretch}{1.2}
    \begin{tabular}{@{}lccccc@{}}
        \toprule
        \textbf{Critère} & \textbf{DQN} & \textbf{Article} & \textbf{Opt. Unif.} & \textbf{Max Power} & \textbf{Random} \\
        \midrule
        EE (Mbit/J) $\uparrow$      & \textbf{255,16} & 207,01 & 213,04 & 10,22 & 23,06 \\
        Puissance (mW) $\downarrow$  & \textbf{15,27}  & 28,73  & 0,02   & 199,53 & 100,57 \\
        Stabilité ($1/\sigma$) $\uparrow$ & 0,021 & \textbf{0,033} & 0,019 & 0,392 & 0,064 \\
        Connaissance modèle          & \ding{55}       & \ding{51} & \ding{55} & \ding{55} & \ding{55} \\
        Décentralisé                  & \ding{51}       & \ding{51} & \ding{55} & N/A   & N/A \\
        Adaptabilité                  & \textbf{Haute}  & Basse    & Nulle & Nulle & Nulle \\
        \bottomrule
    \end{tabular}
\end{table*}

\begin{figure}[htbp]
    \centering
    \includegraphics[width=0.7\textwidth]{figures/radar_chart_comparison.png}
    \caption{Diagramme radar comparatif des méthodes sur quatre critères normalisés : EE, économie de puissance (~$1/P$), stabilité ($1/\sigma$), et débit spectral.}
    \label{fig:radar_comparison}
\end{figure}

% ============================================================
\section{Discussion}
\label{sec:discussion}

\subsection{Avantages du DQN}

\begin{enumerate}
    \item \textbf{Performances supérieures} : le DQN surpasse l'algorithme analytique de $+23{,}3\%$ en EE tout en consommant $46{,}8\%$ moins de puissance. Cet avantage est remarquable car il est obtenu \textit{sans connaissance explicite du modèle de canal}.

    \item \textbf{Absence de modèle} (\textit{model-free}) : contrairement à l'approche BRD + Dinkelbach qui nécessite la connaissance exacte de la structure du SINR et la résolution analytique de l'optimisation fractionnaire, le DQN apprend directement à partir des observations et des récompenses. Cela le rend robuste aux inexactitudes de modélisation et adaptable à des scénarios non couverts par la théorie.

    \item \textbf{Adaptabilité} : les agents DQN peuvent s'adapter à des changements dans l'environnement (nombre d'utilisateurs, topologie, conditions de propagation) par ré-entraînement ou \textit{fine-tuning}, sans nécessiter de re-dérivation des solutions analytiques.
\end{enumerate}

\subsection{Limitations et travaux futurs}

\begin{enumerate}
    \item \textbf{Variance plus élevée} : l'écart-type de l'EE du DQN ($47{,}31$~Mbit/J) est supérieur à celui de l'algorithme de l'article ($30{,}67$~Mbit/J), indiquant une plus grande variabilité des performances selon les réalisations de canal. Cela pourrait être atténué par des techniques de régularisation ou un entraînement plus long.

    \item \textbf{Coût d'entraînement} : l'entraînement de $10\,000$ épisodes requiert un temps significatif, notamment avec $K = 5$ agents indépendants. Cependant, ce coût est assumé hors-ligne (\textit{offline}) et la politique résultante est exécutée en temps réel avec une simple passe avant dans le réseau de neurones.

    \item \textbf{Garanties théoriques} : contrairement à l'algorithme de l'article qui fournit des garanties de convergence vers le GNE, l'approche IQL n'a pas de preuve formelle de convergence dans le cas multi-agent non-stationnaire. Néanmoins, les résultats empiriques démontrent une convergence effective.

    \item \textbf{Passage à l'échelle} : des travaux futurs pourraient explorer des architectures de \textit{parameter sharing} entre agents ou des approches de type \textit{Mean Field} pour étendre le système à un plus grand nombre d'utilisateurs.
\end{enumerate}

% ============================================================
\section{Conclusion du chapitre}

Les résultats expérimentaux démontrent que l'approche par Deep Q-Network constitue une alternative viable et performante à l'algorithme analytique de Zappone \textit{et al.} pour le contrôle de puissance décentralisé dans les réseaux 5G Massive MIMO. Avec un gain de $+23{,}3\%$ en efficacité énergétique et une réduction de $46{,}8\%$ en consommation de puissance, le DQN valide le potentiel de l'apprentissage par renforcement profond pour l'optimisation des ressources radio. L'avantage principal réside dans sa nature \textit{model-free}, qui le rend applicable à des scénarios réalistes où le modèle mathématique exact du canal n'est pas disponible ou est trop complexe pour être exploité analytiquement.

% ============================================================
% BIBLIOGRAPHIE (extraits pertinents)
% ============================================================
% À inclure dans le fichier .bib principal :
%
% @article{zappone2015,
%   author  = {Zappone, A. and Björnson, E. and Sanguinetti, L. and Jorswieck, E.},
%   title   = {Energy-Efficient Power Control: A Look at 5G Wireless Technologies},
%   journal = {IEEE Transactions on Signal Processing},
%   year    = {2015},
%   volume  = {64},
%   number  = {7},
%   pages   = {1668--1683}
% }
%
% @article{mnih2015,
%   author  = {Mnih, V. and Kavukcuoglu, K. and Silver, D. and others},
%   title   = {Human-level control through deep reinforcement learning},
%   journal = {Nature},
%   year    = {2015},
%   volume  = {518},
%   pages   = {529--533}
% }
%
% @inproceedings{vanHasselt2016,
%   author  = {van Hasselt, H. and Guez, A. and Silver, D.},
%   title   = {Deep Reinforcement Learning with Double Q-learning},
%   booktitle = {AAAI},
%   year    = {2016}
% }
%
% @article{kingma2015,
%   author  = {Kingma, D. P. and Ba, J.},
%   title   = {Adam: A Method for Stochastic Optimization},
%   journal = {ICLR},
%   year    = {2015}
% }
%
% @inproceedings{glorot2010,
%   author  = {Glorot, X. and Bengio, Y.},
%   title   = {Understanding the difficulty of training deep feedforward neural networks},
%   booktitle = {AISTATS},
%   year    = {2010}
% }
%
% @inproceedings{tan1993,
%   author  = {Tan, M.},
%   title   = {Multi-Agent Reinforcement Learning: Independent vs. Cooperative Agents},
%   booktitle = {ICML},
%   year    = {1993}
% }
%
% @article{lanctot2017,
%   author  = {Lanctot, M. and Zambaldi, V. and Gruslys, A. and others},
%   title   = {A Unified Game-Theoretic Approach to Multiagent Reinforcement Learning},
%   journal = {NeurIPS},
%   year    = {2017}
% }


%==============================================================================
\section{Conclusion et perspectives}
%==============================================================================

[À développer]

\begin{thebibliography}{99}

\bibitem{3GPP} 3GPP, "Technical Specification Group Radio Access Network; NR; Physical layer; General description (Release 15)," 3GPP TS 38.201, 2018.

\bibitem{Fehske2011} A. Fehske, J. Malmodin, G. Biczok, and G. Fettweis, "The global footprint of mobile communications: The ecological and economic perspective," \textit{IEEE Commun. Mag.}, pp. 55-62, Aug. 2011.

\bibitem{Ericsson2011} Ericsson White Paper, "More than 50 billion connected devices," Ericsson, Tech. Rep. 284 23-3149 Uen, Feb. 2011.

\bibitem{Qualcomm} Qualcomm, "The 1000x data challenge," Tech. Rep.

\bibitem{Zappone2021} A. Zappone, L. Sanguinetti, G. Bacci, E. Jorswieck, and M. Debbah, "Energy-Efficient Power Control: A Look at 5G Wireless Technologies," \textit{IEEE Transactions on Signal Processing}, 2021.

\bibitem{Dinkelbach1967} W. Dinkelbach, "On nonlinear fractional programming," \textit{Management Science}, vol. 13, no. 7, pp. 492-498, Mar. 1967.

\bibitem{Ng2012} D. W. K. Ng, E. S. Lo, and R. Schober, "Energy-efficient resource allocation in multi-cell OFDMA systems with limited backhaul capacity," \textit{IEEE Trans. Wireless Commun.}, vol. 11, no. 10, pp. 3618-3631, Oct. 2012.

\bibitem{Xu2014} Q. Xu, X. Li, H. Ji, and X. Du, "Energy-efficient resource allocation for heterogeneous services in OFDMA downlink networks: Systematic perspective," \textit{IEEE Trans. Veh. Technol.}, vol. 63, no. 5, pp. 2071-2082, June 2014.

\bibitem{Xu2013} J. Xu and L. Qiu, "Energy efficiency optimization for MIMO broadcast channels," \textit{IEEE Trans. Wireless Commun.}, vol. 12, no. 2, pp. 690-701, Feb. 2013.

\bibitem{Du2014} B. Du, C. Pan, W. Zhang, and M. Chen, "Distributed energy-efficient power optimization for CoMP systems with max-min fairness," \textit{IEEE Commun. Lett.}, vol. 18, no. 6, pp. 999-1002, Jun. 2014.

\bibitem{He2013} S. He, Y. Huang, S. Jin, and L. Yang, "Coordinated beamforming for energy efficient transmission in multicell multiuser systems," \textit{IEEE Trans. Commun.}, vol. 61, no. 12, pp. 4961-4971, Dec. 2013.

\bibitem{He2014} S. He, Y. Huang, L. Yang, and B. Ottersten, "Coordinated multicell multiuser precoding for maximizing weighted sum energy efficiency," \textit{IEEE Trans. Signal Process.}, vol. 62, no. 3, pp. 741-751, Feb. 2014.

\bibitem{Venturino2015} L. Venturino, A. Zappone, C. Risi, and S. Buzzi, "Energy-efficient scheduling and power allocation in downlink OFDMA networks with base station coordination," \textit{IEEE Trans. Wireless Commun.}, vol. 14, no. 1, pp. 1-14, 2015.

\bibitem{Miao2011} G. Miao, N. Himayat, G. Li, and S. Talwar, "Distributed interference-aware energy-efficient power optimization," \textit{IEEE Trans. Wireless Commun.}, vol. 10, no. 4, pp. 1323-1333, Apr. 2011.

\bibitem{Scutari2008} G. Scutari, D. Palomar, and S. Barbarossa, "Optimal linear precoding strategies for wideband noncooperative systems based on game theory part I: Nash equilibria," \textit{IEEE Trans. Signal Process.}, vol. 56, no. 3, pp. 1230-1249, March 2008.

\bibitem{Zappone2013} A. Zappone, Z. Chong, E. Jorswieck, and S. Buzzi, "Energy-aware competitive power control in relay-assisted interference wireless networks," \textit{IEEE Trans. Wireless Commun.}, vol. 12, no. 4, pp. 1860-1871, Apr. 2013.

\bibitem{Stupia2015} I. Stupia, L. Sanguinetti, G. Bacci, and L. Vandendorpe, "Power control in networks with heterogeneous users: A quasi-variational inequality approach," \textit{IEEE Trans. Signal Process.}, vol. 63, no. 21, pp. 5691-5705, November 2015.

\bibitem{Bacci2015} G. Bacci, E. Belmega, P. Mertikopoulos, and L. Sanguinetti, "Energy-aware competitive power allocation for heterogeneous networks under QoS constraints," \textit{IEEE Trans. Wireless Commun.}, vol. 14, no. 9, pp. 4728-4742, Sept 2015.

\bibitem{Yates1995} R. D. Yates, "A framework for uplink power control in cellular radio systems," \textit{IEEE J. Sel. Areas Commun.}, vol. 13, no. 7, pp. 1341-1347, September 1995.

\end{thebibliography}

\end{document}